% !TEX program = lualatex
% ---------------------------------------------------------------------------------------------
%  Readback note 29: The tower method at a place: floor measures, gauges and the bounds on each floor
%
%  Build:  python3 tools/rb.py note <slug>
%
%  One unit of a Lean artifact, rewritten as ordinary mathematics. The prose is hand-written. Lean names,
%  line ranges and counts come from data.tex; quoted Lean statements are copied from the artifact by the build.
%
%  Rules of the text: standard terminology; every symbol is explained where it first appears and listed again
%  in the notation at the end; no letter has two meanings; a proof is complete or says what it takes for granted;
%  Lean names appear only in the LEAN lines and in the concordance.
%
%  The signature rule (the build enforces it): a function, relation or predicate that the artifact defines is
%  written only through its notation macro, which the build makes from the spec (`notation`) with one slot for
%  each explicit argument of the Lean definition. No argument is fixed and dropped, also not "for brevity" and
%  not for one section. Every argument and every hypothesis of a stated declaration stands inside its
%  statement, or in a Setting that the statement cites by number; only implicit and instance arguments (the
%  ambient objects) may stand in the conventions.
% ---------------------------------------------------------------------------------------------
\documentclass[11pt,a4paper]{article}
\ifdefined\pdfvariable\pdfvariable suppressoptionalinfo \numexpr 1+2+4+8\relax\fi
\usepackage[a4paper,hmargin=35mm,top=22mm,bottom=27mm,footskip=12mm,nohead]{geometry}

\usepackage{fontspec}
\directlua{
  luaotfload.add_fallback("gfallback", {"STIXTwoMath-Regular.otf:mode=node;", "DejaVuSans.ttf:mode=node;"})
  luaotfload.add_fallback("gmonofallback", {"DejaVuSansMono.ttf:mode=node;", "STIXTwoMath-Regular.otf:mode=node;"})
}
\setmainfont{NewCM10}[Extension=.otf,UprightFont=*-Book,ItalicFont=*-BookItalic,
  BoldFont=*-Bold,BoldItalicFont=*-BoldItalic,SmallCapsFont=*-Book,SmallCapsFeatures={Letters=SmallCaps},
  RawFeature={fallback=gfallback}]
\setsansfont{NewCMSans10}[Extension=.otf,UprightFont=*-Book,ItalicFont=*-BookOblique,
  BoldFont=*-Bold,BoldItalicFont=*-BoldOblique,RawFeature={fallback=gfallback}]
\setmonofont{JetBrainsMono}[Extension=.otf,UprightFont=*-Regular,ItalicFont=*-Italic,
  BoldFont=*-SemiBold,BoldItalicFont=*-SemiBoldItalic,Scale=MatchLowercase,
  RawFeature={-calt;-liga;fallback=gmonofallback}]
\usepackage{amsmath}
\usepackage{amsthm}
\usepackage[math-style=TeX]{unicode-math}
\setmathfont{NewCMMath-Book.otf}
\emergencystretch=2.5em\relax

\usepackage{xcolor}
\usepackage{tikz}
\usetikzlibrary{calc,arrows.meta,positioning,backgrounds}
\usepackage[protrusion=true,expansion=true]{microtype}
\usepackage{enumitem}
\usepackage{needspace}
\usepackage{booktabs}
\usepackage{array}
\usepackage{longtable}
\usepackage{fancyhdr}
\usepackage{lastpage}
\usepackage{multicol}
\usepackage[hidelinks,pdfusetitle]{hyperref}

\definecolor{ink}{HTML}{0F172A}
\definecolor{inksoft}{HTML}{475569}
\definecolor{inkfaint}{HTML}{94A3B8}
\definecolor{rulegrey}{HTML}{CBD5E1}
\definecolor{arith}{HTML}{C2410C}
\definecolor{est}{HTML}{7C3AED}
\definecolor{ana}{HTML}{0891B2}
\definecolor{ok}{HTML}{047857}
\definecolor{leancol}{HTML}{3730A3}

\title{The tower method at a place: floor measures, gauges and the bounds on each floor}
\author{grh-compact}
\date{}

% ------------------------------------------------------------------------------------- macros
\newcommand{\leanfont}{\ttfamily\fontsize{8.6}{10.5}\selectfont\color{leancol}}
\newcommand{\lean}[1]{\csname LN@#1\endcsname}
\newcommand{\fillin}[1]{{\color{arith}\itshape[#1]}}  % a place of the template that is not written yet; the build reports every one that is left
\newcommand{\leanquote}[1]{\csname LQ@#1\endcsname}
% generated by rb.py note; do not edit
\newcommand{\ConcDecls}{42}
\newcommand{\Pin}{72fead10}
\newcommand{\PinText}{that of commit {\leanfont 72fead10}}
\newcommand{\ArtifactName}{grh-compact}
\expandafter\newcommand\csname LN@sieveVector\endcsname{\href{https://github.com/lsernau/grh-lean/blob/72fead102f663d5c7a407964a13dc0cf9d977207/grh-compact/src/New/RealPlace/BandLimits.lean\#L17-L17}{{\leanfont sieveVector}}}
\expandafter\newcommand\csname LN@continuoussieveVector\endcsname{\href{https://github.com/lsernau/grh-lean/blob/72fead102f663d5c7a407964a13dc0cf9d977207/grh-compact/src/New/RealPlace/BandLimits.lean\#L26-L35}{{\leanfont IsCutBessel.\allowbreak{}continuous\_sieveVector}}}
\expandafter\newcommand\csname LN@hasBandLimitsofisCutBessel\endcsname{\href{https://github.com/lsernau/grh-lean/blob/72fead102f663d5c7a407964a13dc0cf9d977207/grh-compact/src/New/RealPlace/BandLimits.lean\#L43-L125}{{\leanfont hasBandLimits\_of\_isCutBessel}}}
\expandafter\newcommand\csname LQ@hasBandLimitsofisCutBessel\endcsname{\leanquoteline{0}{lemma hasBandLimits\_of\_isCutBessel [Nonempty I]}\leanquoteline{4}{\{deg\} (hd : 0 < deg) (hS : SchurLe W (.ofReal deg⁻¹ * card I))}\leanquoteline{2}{(hD : deg * 𝔰.bandConst ≤ b \textasciicircum{} (4 * 𝔰.r \textasciicircum{} 2)) : ∃ J, HasBandLimits b F deg J}}
\expandafter\newcommand\csname LN@IsCutBessel\endcsname{\href{https://github.com/lsernau/grh-lean/blob/72fead102f663d5c7a407964a13dc0cf9d977207/grh-compact/src/New/RealPlace/ComplexShells.lean\#L18-L18}{{\leanfont IsCutBessel}}}
\expandafter\newcommand\csname LN@formConst\endcsname{\href{https://github.com/lsernau/grh-lean/blob/72fead102f663d5c7a407964a13dc0cf9d977207/grh-compact/src/New/RealPlace/ComplexShells.lean\#L26-L26}{{\leanfont IsCutBessel.\allowbreak{}formConst}}}
\expandafter\newcommand\csname LN@isCutcomplexLevel\endcsname{\href{https://github.com/lsernau/grh-lean/blob/72fead102f663d5c7a407964a13dc0cf9d977207/grh-compact/src/New/RealPlace/ComplexShells.lean\#L28-L30}{{\leanfont IsCutBessel.\allowbreak{}isCut\_complexLevel}}}
\expandafter\newcommand\csname LN@volumeExponenteq\endcsname{\href{https://github.com/lsernau/grh-lean/blob/72fead102f663d5c7a407964a13dc0cf9d977207/grh-compact/src/New/RealPlace/ComplexShells.lean\#L32-L33}{{\leanfont IsCutBessel.\allowbreak{}volumeExponent\_eq}}}
\expandafter\newcommand\csname LN@existsgaugeisHead\endcsname{\href{https://github.com/lsernau/grh-lean/blob/72fead102f663d5c7a407964a13dc0cf9d977207/grh-compact/src/New/RealPlace/ComplexShells.lean\#L38-L48}{{\leanfont IsCutBessel.\allowbreak{}exists\_gauge\_isHead}}}
\expandafter\newcommand\csname LN@formConstmullintegralle\endcsname{\href{https://github.com/lsernau/grh-lean/blob/72fead102f663d5c7a407964a13dc0cf9d977207/grh-compact/src/New/RealPlace/ComplexShells.lean\#L56-L101}{{\leanfont IsCutBessel.\allowbreak{}formConst\_mul\_lintegral\_le}}}
\expandafter\newcommand\csname LN@lintegralformMajorantnetop\endcsname{\href{https://github.com/lsernau/grh-lean/blob/72fead102f663d5c7a407964a13dc0cf9d977207/grh-compact/src/New/RealPlace/ComplexShells.lean\#L106-L110}{{\leanfont IsCutBessel.\allowbreak{}lintegral\_formMajorant\_ne\_top}}}
\expandafter\newcommand\csname LN@normfloorFormle\endcsname{\href{https://github.com/lsernau/grh-lean/blob/72fead102f663d5c7a407964a13dc0cf9d977207/grh-compact/src/New/RealPlace/ComplexShells.lean\#L115-L144}{{\leanfont IsCutBessel.\allowbreak{}norm\_floorForm\_le}}}
\expandafter\newcommand\csname LN@headDensity\endcsname{\href{https://github.com/lsernau/grh-lean/blob/72fead102f663d5c7a407964a13dc0cf9d977207/grh-compact/src/New/RealPlace/LocalTheorem.lean\#L32-L33}{{\leanfont Level.\allowbreak{}headDensity}}}
\expandafter\newcommand\csname LN@translatePairing\endcsname{\href{https://github.com/lsernau/grh-lean/blob/72fead102f663d5c7a407964a13dc0cf9d977207/grh-compact/src/New/RealPlace/LocalTheorem.lean\#L36-L36}{{\leanfont translatePairing}}}
\expandafter\newcommand\csname LN@IsBessel\endcsname{\href{https://github.com/lsernau/grh-lean/blob/72fead102f663d5c7a407964a13dc0cf9d977207/grh-compact/src/New/RealPlace/LocalTheorem.lean\#L40-L41}{{\leanfont Place.\allowbreak{}IsBessel}}}
\expandafter\newcommand\csname LN@IsCut\endcsname{\href{https://github.com/lsernau/grh-lean/blob/72fead102f663d5c7a407964a13dc0cf9d977207/grh-compact/src/New/RealPlace/LocalTheorem.lean\#L45-L46}{{\leanfont Level.\allowbreak{}IsCut}}}
\expandafter\newcommand\csname LN@pairingNorm\endcsname{\href{https://github.com/lsernau/grh-lean/blob/72fead102f663d5c7a407964a13dc0cf9d977207/grh-compact/src/New/RealPlace/LocalTheorem.lean\#L50-L50}{{\leanfont Level.\allowbreak{}pairingNorm}}}
\expandafter\newcommand\csname LN@formIntegrandOne\endcsname{\href{https://github.com/lsernau/grh-lean/blob/72fead102f663d5c7a407964a13dc0cf9d977207/grh-compact/src/New/RealPlace/LocalTheorem.lean\#L53-L54}{{\leanfont Level.\allowbreak{}formIntegrandOne}}}
\expandafter\newcommand\csname LN@formIntegrand\endcsname{\href{https://github.com/lsernau/grh-lean/blob/72fead102f663d5c7a407964a13dc0cf9d977207/grh-compact/src/New/RealPlace/LocalTheorem.lean\#L59-L60}{{\leanfont Level.\allowbreak{}formIntegrand}}}
\expandafter\newcommand\csname LN@floorForm\endcsname{\href{https://github.com/lsernau/grh-lean/blob/72fead102f663d5c7a407964a13dc0cf9d977207/grh-compact/src/New/RealPlace/LocalTheorem.lean\#L64-L65}{{\leanfont Level.\allowbreak{}floorForm}}}
\expandafter\newcommand\csname LN@formModulus\endcsname{\href{https://github.com/lsernau/grh-lean/blob/72fead102f663d5c7a407964a13dc0cf9d977207/grh-compact/src/New/RealPlace/LocalTheorem.lean\#L69-L69}{{\leanfont Level.\allowbreak{}formModulus}}}
\expandafter\newcommand\csname LN@formMajorant\endcsname{\href{https://github.com/lsernau/grh-lean/blob/72fead102f663d5c7a407964a13dc0cf9d977207/grh-compact/src/New/RealPlace/LocalTheorem.lean\#L74-L76}{{\leanfont Level.\allowbreak{}formMajorant}}}
\expandafter\newcommand\csname LN@selfNormProduct\endcsname{\href{https://github.com/lsernau/grh-lean/blob/72fead102f663d5c7a407964a13dc0cf9d977207/grh-compact/src/New/RealPlace/LocalTheorem.lean\#L80-L81}{{\leanfont Gauge.\allowbreak{}selfNormProduct}}}
\expandafter\newcommand\csname LN@frameGramdiagmul\endcsname{\href{https://github.com/lsernau/grh-lean/blob/72fead102f663d5c7a407964a13dc0cf9d977207/grh-compact/src/New/RealPlace/LocalTheorem.lean\#L92-L96}{{\leanfont frameGram\_diag\_mul}}}
\expandafter\newcommand\csname LN@pivotModulussqleofcapped\endcsname{\href{https://github.com/lsernau/grh-lean/blob/72fead102f663d5c7a407964a13dc0cf9d977207/grh-compact/src/New/RealPlace/LocalTheorem.lean\#L101-L115}{{\leanfont Gauge.\allowbreak{}pivotModulus\_sq\_le\_of\_capped}}}
\expandafter\newcommand\csname LN@lintegralselfNormle\endcsname{\href{https://github.com/lsernau/grh-lean/blob/72fead102f663d5c7a407964a13dc0cf9d977207/grh-compact/src/New/RealPlace/LocalTheorem.lean\#L120-L175}{{\leanfont Gauge.\allowbreak{}lintegral\_selfNorm\_le}}}
\expandafter\newcommand\csname LN@enormfloorFormle\endcsname{\href{https://github.com/lsernau/grh-lean/blob/72fead102f663d5c7a407964a13dc0cf9d977207/grh-compact/src/New/RealPlace/LocalTheorem.lean\#L178-L181}{{\leanfont Level.\allowbreak{}enorm\_floorForm\_le}}}
\expandafter\newcommand\csname LN@stackMeasuretop\endcsname{\href{https://github.com/lsernau/grh-lean/blob/72fead102f663d5c7a407964a13dc0cf9d977207/grh-compact/src/New/RealPlace/LocalTheorem.lean\#L183-L184}{{\leanfont Level.\allowbreak{}stackMeasure\_top}}}
\expandafter\newcommand\csname LN@isCompacthasMatSize\endcsname{\href{https://github.com/lsernau/grh-lean/blob/72fead102f663d5c7a407964a13dc0cf9d977207/grh-compact/src/New/RealPlace/LocalTheorem.lean\#L186-L187}{{\leanfont Gauge.\allowbreak{}isCompact\_hasMatSize}}}
\expandafter\newcommand\csname LN@isNearofnezero\endcsname{\href{https://github.com/lsernau/grh-lean/blob/72fead102f663d5c7a407964a13dc0cf9d977207/grh-compact/src/New/RealPlace/LocalTheorem.lean\#L191-L192}{{\leanfont Level.\allowbreak{}IsCut.\allowbreak{}isNear\_of\_ne\_zero}}}
\expandafter\newcommand\csname LN@formModulusleformMajorant\endcsname{\href{https://github.com/lsernau/grh-lean/blob/72fead102f663d5c7a407964a13dc0cf9d977207/grh-compact/src/New/RealPlace/LocalTheorem.lean\#L197-L203}{{\leanfont Level.\allowbreak{}IsCut.\allowbreak{}formModulus\_le\_formMajorant}}}
\expandafter\newcommand\csname LN@hasMatSizeofnezero\endcsname{\href{https://github.com/lsernau/grh-lean/blob/72fead102f663d5c7a407964a13dc0cf9d977207/grh-compact/src/New/RealPlace/LocalTheorem.lean\#L209-L213}{{\leanfont Gauge.\allowbreak{}hasMatSize\_of\_ne\_zero}}}
\expandafter\newcommand\csname LN@innerbesselCorrelation\endcsname{\href{https://github.com/lsernau/grh-lean/blob/72fead102f663d5c7a407964a13dc0cf9d977207/grh-compact/src/New/RealPlace/LocalTheorem.lean\#L218-L227}{{\leanfont Gauge.\allowbreak{}inner\_besselCorrelation}}}
\expandafter\newcommand\csname LN@isPairingKerneltranslatePairing\endcsname{\href{https://github.com/lsernau/grh-lean/blob/72fead102f663d5c7a407964a13dc0cf9d977207/grh-compact/src/New/RealPlace/LocalTheorem.lean\#L231-L251}{{\leanfont Gauge.\allowbreak{}isPairingKernel\_translatePairing}}}
\expandafter\newcommand\csname LN@isPairingKernelself\endcsname{\href{https://github.com/lsernau/grh-lean/blob/72fead102f663d5c7a407964a13dc0cf9d977207/grh-compact/src/New/RealPlace/LocalTheorem.lean\#L255-L256}{{\leanfont Gauge.\allowbreak{}isPairingKernel\_self}}}
\expandafter\newcommand\csname LN@formIntegrandOneeq\endcsname{\href{https://github.com/lsernau/grh-lean/blob/72fead102f663d5c7a407964a13dc0cf9d977207/grh-compact/src/New/RealPlace/LocalTheorem.lean\#L261-L264}{{\leanfont Level.\allowbreak{}formIntegrandOne\_eq}}}
\expandafter\newcommand\csname LN@integralsuccflooreq\endcsname{\href{https://github.com/lsernau/grh-lean/blob/72fead102f663d5c7a407964a13dc0cf9d977207/grh-compact/src/New/RealPlace/LocalTheorem.lean\#L270-L303}{{\leanfont Level.\allowbreak{}integral\_succ\_floor\_eq}}}
\expandafter\newcommand\csname LN@integralfloorFormsucc\endcsname{\href{https://github.com/lsernau/grh-lean/blob/72fead102f663d5c7a407964a13dc0cf9d977207/grh-compact/src/New/RealPlace/LocalTheorem.lean\#L310-L318}{{\leanfont Gauge.\allowbreak{}integral\_floorForm\_succ}}}
\expandafter\newcommand\csname LN@lintegralpositionsle\endcsname{\href{https://github.com/lsernau/grh-lean/blob/72fead102f663d5c7a407964a13dc0cf9d977207/grh-compact/src/New/RealPlace/LocalTheorem.lean\#L326-L343}{{\leanfont Gauge.\allowbreak{}lintegral\_positions\_le}}}
\expandafter\newcommand\csname LN@sqlintegralleselfNormProduct\endcsname{\href{https://github.com/lsernau/grh-lean/blob/72fead102f663d5c7a407964a13dc0cf9d977207/grh-compact/src/New/RealPlace/LocalTheorem.lean\#L349-L363}{{\leanfont Gauge.\allowbreak{}sq\_lintegral\_le\_selfNormProduct}}}
\expandafter\newcommand\csname LN@formMajorantsqle\endcsname{\href{https://github.com/lsernau/grh-lean/blob/72fead102f663d5c7a407964a13dc0cf9d977207/grh-compact/src/New/RealPlace/LocalTheorem.lean\#L368-L380}{{\leanfont Gauge.\allowbreak{}formMajorant\_sq\_le}}}
\expandafter\newcommand\csname LN@formModulussqleofchar\endcsname{\href{https://github.com/lsernau/grh-lean/blob/72fead102f663d5c7a407964a13dc0cf9d977207/grh-compact/src/New/RealPlace/LocalTheorem.lean\#L389-L416}{{\leanfont Gauge.\allowbreak{}formModulus\_sq\_le\_of\_char}}}
\expandafter\newcommand\csname LN@formModulussqleofheads\endcsname{\href{https://github.com/lsernau/grh-lean/blob/72fead102f663d5c7a407964a13dc0cf9d977207/grh-compact/src/New/RealPlace/LocalTheorem.lean\#L421-L458}{{\leanfont Gauge.\allowbreak{}formModulus\_sq\_le\_of\_heads}}}
\expandafter\newcommand\csname LQ@formModulussqleofheads\endcsname{\leanquoteline{0}{lemma Gauge.formModulus\_sq\_le\_of\_heads \{i\} (hi : j = i + 1) (h : ξ =}\leanquoteline{2}{F.g⁻¹ * ((p i)⁻¹.1 - (q i)⁻¹.1)) :}\leanquoteline{2}{F.formModulus j p q W x 1 \textasciicircum{} 2 ≤ (P.c / P.md ξ) \textasciicircum{} card ι * Γ.selfNormProduct j p q W W x}}
\newcommand{\ConcordanceRows}{\crow{sieveVector}{\href{https://github.com/lsernau/grh-lean/blob/72fead102f663d5c7a407964a13dc0cf9d977207/grh-compact/src/New/RealPlace/BandLimits.lean\#L17-L17}{{\leanfont sieveVector}}}{BandLimits.lean}{17--17}
\crow{continuoussieveVector}{\href{https://github.com/lsernau/grh-lean/blob/72fead102f663d5c7a407964a13dc0cf9d977207/grh-compact/src/New/RealPlace/BandLimits.lean\#L26-L35}{{\leanfont IsCutBessel.\allowbreak{}continuous\_sieveVector}}}{BandLimits.lean}{26--35}
\crow{hasBandLimitsofisCutBessel}{\href{https://github.com/lsernau/grh-lean/blob/72fead102f663d5c7a407964a13dc0cf9d977207/grh-compact/src/New/RealPlace/BandLimits.lean\#L43-L125}{{\leanfont hasBandLimits\_of\_isCutBessel}}}{BandLimits.lean}{43--125}
\crow{IsCutBessel}{\href{https://github.com/lsernau/grh-lean/blob/72fead102f663d5c7a407964a13dc0cf9d977207/grh-compact/src/New/RealPlace/ComplexShells.lean\#L18-L18}{{\leanfont IsCutBessel}}}{ComplexShells.lean}{18--18}
\crow{formConst}{\href{https://github.com/lsernau/grh-lean/blob/72fead102f663d5c7a407964a13dc0cf9d977207/grh-compact/src/New/RealPlace/ComplexShells.lean\#L26-L26}{{\leanfont IsCutBessel.\allowbreak{}formConst}}}{ComplexShells.lean}{26--26}
\crow{isCutcomplexLevel}{\href{https://github.com/lsernau/grh-lean/blob/72fead102f663d5c7a407964a13dc0cf9d977207/grh-compact/src/New/RealPlace/ComplexShells.lean\#L28-L30}{{\leanfont IsCutBessel.\allowbreak{}isCut\_complexLevel}}}{ComplexShells.lean}{28--30}
\crow{volumeExponenteq}{\href{https://github.com/lsernau/grh-lean/blob/72fead102f663d5c7a407964a13dc0cf9d977207/grh-compact/src/New/RealPlace/ComplexShells.lean\#L32-L33}{{\leanfont IsCutBessel.\allowbreak{}volumeExponent\_eq}}}{ComplexShells.lean}{32--33}
\crow{existsgaugeisHead}{\href{https://github.com/lsernau/grh-lean/blob/72fead102f663d5c7a407964a13dc0cf9d977207/grh-compact/src/New/RealPlace/ComplexShells.lean\#L38-L48}{{\leanfont IsCutBessel.\allowbreak{}exists\_gauge\_isHead}}}{ComplexShells.lean}{38--48}
\crow{formConstmullintegralle}{\href{https://github.com/lsernau/grh-lean/blob/72fead102f663d5c7a407964a13dc0cf9d977207/grh-compact/src/New/RealPlace/ComplexShells.lean\#L56-L101}{{\leanfont IsCutBessel.\allowbreak{}formConst\_mul\_lintegral\_le}}}{ComplexShells.lean}{56--101}
\crow{lintegralformMajorantnetop}{\href{https://github.com/lsernau/grh-lean/blob/72fead102f663d5c7a407964a13dc0cf9d977207/grh-compact/src/New/RealPlace/ComplexShells.lean\#L106-L110}{{\leanfont IsCutBessel.\allowbreak{}lintegral\_formMajorant\_ne\_top}}}{ComplexShells.lean}{106--110}
\crow{normfloorFormle}{\href{https://github.com/lsernau/grh-lean/blob/72fead102f663d5c7a407964a13dc0cf9d977207/grh-compact/src/New/RealPlace/ComplexShells.lean\#L115-L144}{{\leanfont IsCutBessel.\allowbreak{}norm\_floorForm\_le}}}{ComplexShells.lean}{115--144}
\crow{headDensity}{\href{https://github.com/lsernau/grh-lean/blob/72fead102f663d5c7a407964a13dc0cf9d977207/grh-compact/src/New/RealPlace/LocalTheorem.lean\#L32-L33}{{\leanfont Level.\allowbreak{}headDensity}}}{LocalTheorem.lean}{32--33}
\crow{translatePairing}{\href{https://github.com/lsernau/grh-lean/blob/72fead102f663d5c7a407964a13dc0cf9d977207/grh-compact/src/New/RealPlace/LocalTheorem.lean\#L36-L36}{{\leanfont translatePairing}}}{LocalTheorem.lean}{36--36}
\crow{IsBessel}{\href{https://github.com/lsernau/grh-lean/blob/72fead102f663d5c7a407964a13dc0cf9d977207/grh-compact/src/New/RealPlace/LocalTheorem.lean\#L40-L41}{{\leanfont Place.\allowbreak{}IsBessel}}}{LocalTheorem.lean}{40--41}
\crow{IsCut}{\href{https://github.com/lsernau/grh-lean/blob/72fead102f663d5c7a407964a13dc0cf9d977207/grh-compact/src/New/RealPlace/LocalTheorem.lean\#L45-L46}{{\leanfont Level.\allowbreak{}IsCut}}}{LocalTheorem.lean}{45--46}
\crow{pairingNorm}{\href{https://github.com/lsernau/grh-lean/blob/72fead102f663d5c7a407964a13dc0cf9d977207/grh-compact/src/New/RealPlace/LocalTheorem.lean\#L50-L50}{{\leanfont Level.\allowbreak{}pairingNorm}}}{LocalTheorem.lean}{50--50}
\crow{formIntegrandOne}{\href{https://github.com/lsernau/grh-lean/blob/72fead102f663d5c7a407964a13dc0cf9d977207/grh-compact/src/New/RealPlace/LocalTheorem.lean\#L53-L54}{{\leanfont Level.\allowbreak{}formIntegrandOne}}}{LocalTheorem.lean}{53--54}
\crow{formIntegrand}{\href{https://github.com/lsernau/grh-lean/blob/72fead102f663d5c7a407964a13dc0cf9d977207/grh-compact/src/New/RealPlace/LocalTheorem.lean\#L59-L60}{{\leanfont Level.\allowbreak{}formIntegrand}}}{LocalTheorem.lean}{59--60}
\crow{floorForm}{\href{https://github.com/lsernau/grh-lean/blob/72fead102f663d5c7a407964a13dc0cf9d977207/grh-compact/src/New/RealPlace/LocalTheorem.lean\#L64-L65}{{\leanfont Level.\allowbreak{}floorForm}}}{LocalTheorem.lean}{64--65}
\crow{formModulus}{\href{https://github.com/lsernau/grh-lean/blob/72fead102f663d5c7a407964a13dc0cf9d977207/grh-compact/src/New/RealPlace/LocalTheorem.lean\#L69-L69}{{\leanfont Level.\allowbreak{}formModulus}}}{LocalTheorem.lean}{69--69}
\crow{formMajorant}{\href{https://github.com/lsernau/grh-lean/blob/72fead102f663d5c7a407964a13dc0cf9d977207/grh-compact/src/New/RealPlace/LocalTheorem.lean\#L74-L76}{{\leanfont Level.\allowbreak{}formMajorant}}}{LocalTheorem.lean}{74--76}
\crow{selfNormProduct}{\href{https://github.com/lsernau/grh-lean/blob/72fead102f663d5c7a407964a13dc0cf9d977207/grh-compact/src/New/RealPlace/LocalTheorem.lean\#L80-L81}{{\leanfont Gauge.\allowbreak{}selfNormProduct}}}{LocalTheorem.lean}{80--81}
\crow{frameGramdiagmul}{\href{https://github.com/lsernau/grh-lean/blob/72fead102f663d5c7a407964a13dc0cf9d977207/grh-compact/src/New/RealPlace/LocalTheorem.lean\#L92-L96}{{\leanfont frameGram\_diag\_mul}}}{LocalTheorem.lean}{92--96}
\crow{pivotModulussqleofcapped}{\href{https://github.com/lsernau/grh-lean/blob/72fead102f663d5c7a407964a13dc0cf9d977207/grh-compact/src/New/RealPlace/LocalTheorem.lean\#L101-L115}{{\leanfont Gauge.\allowbreak{}pivotModulus\_sq\_le\_of\_capped}}}{LocalTheorem.lean}{101--115}
\crow{lintegralselfNormle}{\href{https://github.com/lsernau/grh-lean/blob/72fead102f663d5c7a407964a13dc0cf9d977207/grh-compact/src/New/RealPlace/LocalTheorem.lean\#L120-L175}{{\leanfont Gauge.\allowbreak{}lintegral\_selfNorm\_le}}}{LocalTheorem.lean}{120--175}
\crow{enormfloorFormle}{\href{https://github.com/lsernau/grh-lean/blob/72fead102f663d5c7a407964a13dc0cf9d977207/grh-compact/src/New/RealPlace/LocalTheorem.lean\#L178-L181}{{\leanfont Level.\allowbreak{}enorm\_floorForm\_le}}}{LocalTheorem.lean}{178--181}
\crow{stackMeasuretop}{\href{https://github.com/lsernau/grh-lean/blob/72fead102f663d5c7a407964a13dc0cf9d977207/grh-compact/src/New/RealPlace/LocalTheorem.lean\#L183-L184}{{\leanfont Level.\allowbreak{}stackMeasure\_top}}}{LocalTheorem.lean}{183--184}
\crow{isCompacthasMatSize}{\href{https://github.com/lsernau/grh-lean/blob/72fead102f663d5c7a407964a13dc0cf9d977207/grh-compact/src/New/RealPlace/LocalTheorem.lean\#L186-L187}{{\leanfont Gauge.\allowbreak{}isCompact\_hasMatSize}}}{LocalTheorem.lean}{186--187}
\crow{isNearofnezero}{\href{https://github.com/lsernau/grh-lean/blob/72fead102f663d5c7a407964a13dc0cf9d977207/grh-compact/src/New/RealPlace/LocalTheorem.lean\#L191-L192}{{\leanfont Level.\allowbreak{}IsCut.\allowbreak{}isNear\_of\_ne\_zero}}}{LocalTheorem.lean}{191--192}
\crow{formModulusleformMajorant}{\href{https://github.com/lsernau/grh-lean/blob/72fead102f663d5c7a407964a13dc0cf9d977207/grh-compact/src/New/RealPlace/LocalTheorem.lean\#L197-L203}{{\leanfont Level.\allowbreak{}IsCut.\allowbreak{}formModulus\_le\_formMajorant}}}{LocalTheorem.lean}{197--203}
\crow{hasMatSizeofnezero}{\href{https://github.com/lsernau/grh-lean/blob/72fead102f663d5c7a407964a13dc0cf9d977207/grh-compact/src/New/RealPlace/LocalTheorem.lean\#L209-L213}{{\leanfont Gauge.\allowbreak{}hasMatSize\_of\_ne\_zero}}}{LocalTheorem.lean}{209--213}
\crow{innerbesselCorrelation}{\href{https://github.com/lsernau/grh-lean/blob/72fead102f663d5c7a407964a13dc0cf9d977207/grh-compact/src/New/RealPlace/LocalTheorem.lean\#L218-L227}{{\leanfont Gauge.\allowbreak{}inner\_besselCorrelation}}}{LocalTheorem.lean}{218--227}
\crow{isPairingKerneltranslatePairing}{\href{https://github.com/lsernau/grh-lean/blob/72fead102f663d5c7a407964a13dc0cf9d977207/grh-compact/src/New/RealPlace/LocalTheorem.lean\#L231-L251}{{\leanfont Gauge.\allowbreak{}isPairingKernel\_translatePairing}}}{LocalTheorem.lean}{231--251}
\crow{isPairingKernelself}{\href{https://github.com/lsernau/grh-lean/blob/72fead102f663d5c7a407964a13dc0cf9d977207/grh-compact/src/New/RealPlace/LocalTheorem.lean\#L255-L256}{{\leanfont Gauge.\allowbreak{}isPairingKernel\_self}}}{LocalTheorem.lean}{255--256}
\crow{formIntegrandOneeq}{\href{https://github.com/lsernau/grh-lean/blob/72fead102f663d5c7a407964a13dc0cf9d977207/grh-compact/src/New/RealPlace/LocalTheorem.lean\#L261-L264}{{\leanfont Level.\allowbreak{}formIntegrandOne\_eq}}}{LocalTheorem.lean}{261--264}
\crow{integralsuccflooreq}{\href{https://github.com/lsernau/grh-lean/blob/72fead102f663d5c7a407964a13dc0cf9d977207/grh-compact/src/New/RealPlace/LocalTheorem.lean\#L270-L303}{{\leanfont Level.\allowbreak{}integral\_succ\_floor\_eq}}}{LocalTheorem.lean}{270--303}
\crow{integralfloorFormsucc}{\href{https://github.com/lsernau/grh-lean/blob/72fead102f663d5c7a407964a13dc0cf9d977207/grh-compact/src/New/RealPlace/LocalTheorem.lean\#L310-L318}{{\leanfont Gauge.\allowbreak{}integral\_floorForm\_succ}}}{LocalTheorem.lean}{310--318}
\crow{lintegralpositionsle}{\href{https://github.com/lsernau/grh-lean/blob/72fead102f663d5c7a407964a13dc0cf9d977207/grh-compact/src/New/RealPlace/LocalTheorem.lean\#L326-L343}{{\leanfont Gauge.\allowbreak{}lintegral\_positions\_le}}}{LocalTheorem.lean}{326--343}
\crow{sqlintegralleselfNormProduct}{\href{https://github.com/lsernau/grh-lean/blob/72fead102f663d5c7a407964a13dc0cf9d977207/grh-compact/src/New/RealPlace/LocalTheorem.lean\#L349-L363}{{\leanfont Gauge.\allowbreak{}sq\_lintegral\_le\_selfNormProduct}}}{LocalTheorem.lean}{349--363}
\crow{formMajorantsqle}{\href{https://github.com/lsernau/grh-lean/blob/72fead102f663d5c7a407964a13dc0cf9d977207/grh-compact/src/New/RealPlace/LocalTheorem.lean\#L368-L380}{{\leanfont Gauge.\allowbreak{}formMajorant\_sq\_le}}}{LocalTheorem.lean}{368--380}
\crow{formModulussqleofchar}{\href{https://github.com/lsernau/grh-lean/blob/72fead102f663d5c7a407964a13dc0cf9d977207/grh-compact/src/New/RealPlace/LocalTheorem.lean\#L389-L416}{{\leanfont Gauge.\allowbreak{}formModulus\_sq\_le\_of\_char}}}{LocalTheorem.lean}{389--416}
\crow{formModulussqleofheads}{\href{https://github.com/lsernau/grh-lean/blob/72fead102f663d5c7a407964a13dc0cf9d977207/grh-compact/src/New/RealPlace/LocalTheorem.lean\#L421-L458}{{\leanfont Gauge.\allowbreak{}formModulus\_sq\_le\_of\_heads}}}{LocalTheorem.lean}{421--458}}
% the notations of the defined functions: one slot for each explicit argument of the Lean declaration
\newcommand{\sieveVector}[2]{\Xi(#1;#2)}
\newcommand{\IsCutBessel}[7]{\operatorname{CB}(#1,#2,#3,#4,#5,#6,#7)}
\newcommand{\formConst}[1]{c_{\mathrm{f}}(#1)}
\newcommand{\headDensity}[4]{\mathsf{h}_{#1}(#4;#2,#3)}
\newcommand{\translatePairing}[5]{\langle #1(#4),#2(#3#5)\rangle}
\newcommand{\IsBessel}[2]{#1\in\operatorname{Bes}(#2)}
\newcommand{\IsCut}[1]{#1\in\operatorname{Cut}}
\newcommand{\pairingNorm}[2]{\operatorname{PN}(#1,#2)}
\newcommand{\formIntegrandOne}[6]{\mathsf{I}^{1}_{#1}(#2;#3,#4,#5)(#6)}
\newcommand{\formIntegrand}[7]{\mathsf{I}_{#1}(#2,#3;#4,#5,#6)(#7)}
\newcommand{\floorForm}[5]{\mathcal{I}_{#1}(#2,#3;#4,#5)}
\newcommand{\formModulus}[6]{\mathcal{I}^{\mathrm{mod}}_{#1}(#2,#3;#4,#5,#6)}
\newcommand{\formMajorant}[6]{\mathcal{I}^{\mathrm{maj}}_{#1}(#2,#3;#4,#5,#6)}
\newcommand{\selfNormProduct}[7]{\operatorname{SNP}_{#1}(#7;#2,#3;#4,#5,#6)}
\newcommand{\UnitMembers}{239}
\newcommand{\UnitInNote}{42}
\newcommand{\UnitProofs}{108}
\newcommand{\UnitSteps}{639}
\newcommand{\UnitAllSteps}{665}


\linespread{1.08}                       % 11 pt on 14.7 pt
\setlength{\parindent}{0pt}
\setlength{\parskip}{5.5pt plus 1.5pt minus 0.5pt}
\setlength{\abovedisplayskip}{7.5pt plus 2pt minus 3pt}
\setlength{\belowdisplayskip}{7.5pt plus 2pt minus 3pt}
\setlength{\abovedisplayshortskip}{3pt plus 2pt}
\setlength{\belowdisplayshortskip}{6pt plus 2pt minus 2pt}
\emergencystretch=1.5em
\clubpenalty=10000 \widowpenalty=10000 \displaywidowpenalty=10000
\relpenalty=7000 \binoppenalty=9000
\raggedbottom
\pagestyle{fancy}\fancyhf{}\renewcommand{\headrulewidth}{0pt}
\fancyfoot[C]{\sffamily\fontsize{7.6}{9}\selectfont\color{inkfaint}\thepage\,/\,\pageref*{LastPage}}
\setlist{topsep=2pt,partopsep=0pt,parsep=0pt,itemsep=4pt,leftmargin=2.1em,labelsep=0.55em}

\newcommand{\Q}{\mathbb{Q}}
\newcommand{\Z}{\mathbb{Z}}
\newcommand{\F}{\mathbb{F}}
\newcommand{\len}{\operatorname{len}}
\newcommand{\cok}{\operatorname{cok}}
\newcommand{\End}{\operatorname{End}}
\newcommand{\Hom}{\operatorname{Hom}}
\newcommand{\Gal}{\operatorname{Gal}}
\newcommand{\diag}{\operatorname{diag}}
\newcommand{\kw}[1]{\emph{#1}}
\newcommand{\m}[1]{\mbox{$#1$}}        % a short formula that must not be broken across lines
\newcommand{\head}[1]{\par\addvspace{17pt plus 4pt minus 2pt}\Needspace{6\baselineskip}%
  \noindent{\fontsize{13.2}{16}\selectfont\bfseries\color{ink}#1}\par\nopagebreak\vspace{-1pt}}
\newcommand{\sect}[2]{\refstepcounter{section}\label{#2}\par\addvspace{19pt plus 4pt minus 2pt}\Needspace{7\baselineskip}%
  \noindent\llap{{\fontsize{13.2}{16}\selectfont\bfseries\color{inkfaint}\thesection}\hspace{9pt}}%
  {\fontsize{13.2}{16}\selectfont\bfseries\color{ink}#1}\par\nopagebreak\vspace{-1pt}}
\newcommand{\subhead}[1]{\par\addvspace{9pt plus 2pt}\Needspace{4\baselineskip}\noindent\textbf{#1}\enspace\ignorespaces}
% where a statement lives in the Lean archive
\newcommand{\leantag}{{\sffamily\bfseries\fontsize{6.3}{8}\selectfont\color{leancol}\addfontfeatures{LetterSpace=7}LEAN}}
\newcommand{\sep}{\unskip\nobreak\hspace{5pt}{\color{inkfaint}\textperiodcentered}\hspace{5pt plus 2pt}\ignorespaces}
\newcommand{\leanline}[1]{\par\nopagebreak\vspace{-2.5pt}{\fontsize{8.9}{12.4}\selectfont\color{inksoft}\raggedright
  \noindent\llap{\leantag\hspace{9pt}}#1\par}\vspace{1pt}}
% a displayed statement in a tinted box: #1 colour, #2 text
\newcommand{\statement}[2]{\par\addvspace{7pt}\begingroup\setlength{\fboxsep}{8.5pt}\setlength{\fboxrule}{0.6pt}%
  \noindent\fcolorbox{#1!30!white}{#1!3!white}{\parbox{\dimexpr\linewidth-2\fboxsep-2\fboxrule\relax}{#2}}%
  \endgroup\par\addvspace{7pt}}
\newcommand{\nt}[2]{\par\noindent\hangindent=1.1em\hangafter=1 #1\enspace{\color{inksoft}#2}\par\vspace{1pt}}
\newcommand{\legdot}[1]{\tikz[baseline=-0.62ex]\fill[#1] (0,0) circle (2.5pt);\hspace{3pt}}
\newcommand{\crow}[4]{\csname CS@#1\endcsname & #2 & \textcolor{inksoft}{#3} & \textcolor{inksoft}{#4}\\}
\newcommand{\leanquoteline}[2]{\par\noindent\hangindent=\dimexpr#1\fontdimen2\font+2.2em\relax\hangafter=1
  \hspace*{#1\fontdimen2\font}#2}
\newcommand{\quoted}[1]{\par\vspace{3pt}{\ttfamily\fontsize{7.5}{10.4}\selectfont\color{leancol}\raggedright
  \setlength{\parskip}{0pt}\leanquote{#1}\par}\vspace{4pt}}
\newcommand{\bx}[3]{{\sffamily\bfseries\fontsize{7.3}{9.2}\selectfont\color{#1!82!black}#2}\\[1.5pt]#3}
\newcommand{\cs}[2]{\expandafter\newcommand\csname CS@#1\endcsname{#2}}

% theorems
\newtheoremstyle{plainx}{7pt plus 2pt}{7pt plus 2pt}{\itshape}{}{\bfseries}{.}{0.5em}{}
\newtheoremstyle{uprightx}{7pt plus 2pt}{7pt plus 2pt}{}{}{\bfseries}{.}{0.5em}{}
\theoremstyle{plainx}
\newtheorem{theorem}{Theorem}[section]
\newtheorem{lemma}[theorem]{Lemma}
\newtheorem{proposition}[theorem]{Proposition}
\newtheorem{corollary}[theorem]{Corollary}
\theoremstyle{uprightx}
\newtheorem{definition}[theorem]{Definition}
\newtheorem{setting}[theorem]{Setting}          % hypotheses shared by several statements; each of them cites it by number
\newtheorem{remark}[theorem]{Remark}
\newtheorem{example}[theorem]{Example}
\renewcommand{\qedsymbol}{{\setlength{\fboxsep}{0pt}\setlength{\fboxrule}{0.5pt}\fbox{\rule{0pt}{0.46em}\hspace{0.46em}}}}
% what holds throughout the note: the ambient objects only (implicit and instance arguments), never a hypothesis
\newenvironment{conventions}{\par\addvspace{5pt}\noindent\textbf{Conventions.}\enspace\ignorespaces}{\par\addvspace{3pt}}
\makeatletter
\renewenvironment{proof}[1][Proof]{\par\pushQED{\qed}\normalfont\topsep0pt\partopsep0pt
  \trivlist\item[\hskip\labelsep\itshape #1.]\ignorespaces}{\popQED\endtrivlist\@endpefalse}
\makeatother

% notations of definitions of unit 132 that this note uses and another note of the unit states
% (made by shared/mkspec.py from shared/notation-132.json; the same TeX as in the note that states them)
\newcommand{\RealPlaceSizes}[4]{\operatorname{RPS}(#1,#2,#3,#4)}
\newcommand{\haarRatio}[1]{\lambda(#1)}
\newcommand{\gauge}[2]{\Gamma(#1;#2)}
\newcommand{\HasFibreBound}[5]{\operatorname{Fib}_{#2,#4}(#1;#3)\le #5}
\newcommand{\gramSection}[1]{\gamma(#1)}
\newcommand{\invCorner}[1]{\operatorname{ic}(#1)}
\newcommand{\cornerDefect}[1]{\operatorname{cd}(#1)}
\newcommand{\IsBoundedInCell}[2]{#1\in\operatorname{Cell}_{\le #2}}
\newcommand{\diagLast}[1]{\operatorname{d}_{\varrho}(#1)}
\newcommand{\sectionCoeff}[2]{\mathsf{c}(#1,#2)}
\newcommand{\sectionFrame}[3]{\mathsf{f}(#1,#2,#3)}
\newcommand{\sectionMatrix}[3]{\mathsf{m}(#1,#2,#3)}
\newcommand{\Levels}[3]{\operatorname{Lev}(#1,#2,#3)}
\newcommand{\floorMeasure}[2]{\theta_{#2}(#1)}
\newcommand{\stackMeasure}[1]{\mu^{#1}}
\newcommand{\groupMeasure}[1]{\mathsf{M}(#1)}
\newcommand{\upperMeasure}[2]{\mathsf{M}_{\ge #2}(#1)}
\newcommand{\IsNear}[3]{#3\in\operatorname{Near}_{#2}(#1)}
\newcommand{\orthonormalPart}[1]{\mathsf{O}(#1)}
\newcommand{\closePairs}[1]{\operatorname{Cl}(#1)}
\newcommand{\floorValues}[4]{\operatorname{Val}_{#2}(#1;#3,#4)}
\newcommand{\volumeExponent}[1]{\mathsf{s}(#1)}
\newcommand{\floorWeight}[4]{\mathsf{w}_{#2}(#1;#3,#4)}
\newcommand{\floorModulus}[3]{\mathsf{d}_{#2}(#1;#3)}
\newcommand{\FloorMassLe}[4]{\operatorname{Mass}_{#2}(#1;#4)\le #3}
\newcommand{\upperPairing}[5]{\mathsf{Av}_{#2}(#1;#3)(#4,#5)}
\newcommand{\IsPairingKernel}[4]{(#3,#4)\in\operatorname{PK}_{#2}(#1)}
\newcommand{\invWeight}[3]{\mathsf{w}^{-1}_{#2}(#1;#3)}
\newcommand{\pivotModulus}[5]{\mathsf{Pv}_{#2}(#1;#3)(#4,#5)}
\newcommand{\weightedSquare}[5]{\mathsf{Sq}_{#2}(#1;#3,#4)(#5)}
\newcommand{\weightedSquareSum}[5]{\mathsf{Sm}_{#4}(#1;#2,#3)(#5)}
\newcommand{\selfNorm}[4]{\operatorname{SN}_{#2}(#1;#3,#4)}
\newcommand{\Gauges}[1]{\operatorname{Gau}(#1)}
\newcommand{\positionsOfSize}[2]{\operatorname{Pos}(#1;#2)}
\newcommand{\Sizes}[1]{\operatorname{Siz}(#1)}
\newcommand{\weylRatio}[3]{w_{#1}(#3)^{-1}w_{#1}(#2)}
\newcommand{\vecsOfSize}[1]{\mathrm{Vec}_{#1}}
\newcommand{\HasMatSize}[2]{\operatorname{sz}(#2)\le #1}
\newcommand{\IsHead}[2]{#2\in\operatorname{Hd}_{#1}}
\newcommand{\frameVolume}[1]{\Delta(#1)}
\newcommand{\frameXPart}[1]{\mathsf{x}_0(#1)}
\newcommand{\floorConst}[2]{C_{#2}(#1)}
\newcommand{\globalConst}[1]{C_{\mathrm{gl}}(#1)}

% hand macros of the three notes of unit 132 (objects of the background library; letters of shared/PLAN-132.md)
\newcommand{\cL}{\mathcal{L}}
\newcommand{\CC}{\mathbb{C}}
\newcommand{\RR}{\mathbb{R}}
\newcommand{\NN}{\mathbb{N}}
\newcommand{\dd}{\mathrm{d}}
\newcommand{\GT}[1]{G(#1)}
\newcommand{\JT}[1]{J(#1)}
\newcommand{\Hl}[2]{\mathcal{H}_{#1}(#2)}
\newcommand{\Ql}[2]{\mathcal{Q}_{#1}(#2)}
\newcommand{\Nb}[2]{\mathcal{N}_{<#1}(#2)}
\newcommand{\Nf}[2]{\mathcal{N}_{#1}(#2)}
\newcommand{\UT}[1]{\mathcal{U}(#1)}
\newcommand{\modA}{\operatorname{mod}_A}
\newcommand{\ZT}[2]{Z_{#1}(#2)}
\newcommand{\sd}{\operatorname{sd}}
\newcommand{\wl}[2]{w_{#1}(#2)}
\newcommand{\wo}[2]{w^{\circ}_{#1}(#2)}
\newcommand{\ip}[2]{\langle #1,#2\rangle}
\newcommand{\gL}[1]{g_{#1}}
\newcommand{\BoxL}[1]{\operatorname{Box}(#1)}
\newcommand{\Posn}[1]{\operatorname{Posn}(#1)}
\newcommand{\Ball}[1]{\mathfrak{B}(#1)}
\newcommand{\etab}[1]{\eta_{<#1}}
\newcommand{\vol}{\operatorname{vol}}
\newcommand{\cLC}[1]{\mathcal{L}_{\CC}(#1)}
\newcommand{\Gone}[1]{G_1(#1)}
\newcommand{\Jone}[1]{J_1(#1)}
\newcommand{\volU}[1]{\vol\mathrm{U}(#1)}
\newcommand{\hw}[3]{\operatorname{hw}_{#1,#2}(#3)}
\newcommand{\damp}[2]{\ell_{#1}(#2)}
\newcommand{\bandc}[1]{\Theta(#1)}
\newcommand{\cellc}[1]{c_G(#1)}
\newcommand{\corr}[3]{J^{#1}_{#2}(#3)}
\newcommand{\boost}[2]{\chi_{#1}(#2)}
\newcommand{\quadric}[2]{\delta^{(#2)}(#1)}
\newcommand{\Inp}[1]{\textup{(#1)}}
% where each declaration of the concordance is treated
\cs{IsBessel}{Definition~\ref{d:bessel}}
\cs{IsCut}{Definition~\ref{d:cut}}
\cs{isNearofnezero}{Definition~\ref{d:cut}}
\cs{translatePairing}{Definition~\ref{d:kernel}}
\cs{pairingNorm}{Definition~\ref{d:pn}}
\cs{headDensity}{Definition~\ref{d:hd}}
\cs{formIntegrandOne}{Definition~\ref{d:integrand}}
\cs{formIntegrand}{Definition~\ref{d:integrand}}
\cs{floorForm}{Definition~\ref{d:form}}
\cs{formModulus}{Definition~\ref{d:form}}
\cs{formMajorant}{Definition~\ref{d:form}}
\cs{selfNormProduct}{Definition~\ref{d:snp}}
\cs{IsCutBessel}{Definition~\ref{d:cutbessel}}
\cs{volumeExponenteq}{Definition~\ref{d:cutbessel}}
\cs{formConst}{Definition~\ref{d:formconst}}
\cs{sieveVector}{Definition~\ref{d:sieve}}
\cs{formMajorantsqle}{Theorem A (a)}
\cs{formModulussqleofchar}{Theorem A (b)}
\cs{formModulussqleofheads}{Theorem A (c)}
\cs{integralfloorFormsucc}{Theorem B}
\cs{lintegralselfNormle}{Theorem C}
\cs{lintegralpositionsle}{Theorem D}
\cs{hasBandLimitsofisCutBessel}{Theorem E}
\cs{stackMeasuretop}{Theorem E}
\cs{isPairingKerneltranslatePairing}{Lemma~\ref{l:pk}}
\cs{isPairingKernelself}{Lemma~\ref{l:pk}}
\cs{formIntegrandOneeq}{Lemma~\ref{l:one}}
\cs{enormfloorFormle}{Lemma~\ref{l:fmm}}
\cs{formModulusleformMajorant}{Lemma~\ref{l:fmm}}
\cs{pivotModulussqleofcapped}{Lemma~\ref{l:cap}}
\cs{sqlintegralleselfNormProduct}{Lemma~\ref{l:snp}}
\cs{frameGramdiagmul}{Lemma~\ref{l:diag}}
\cs{integralsuccflooreq}{Lemma~\ref{l:succ}}
\cs{isCutcomplexLevel}{Lemma~\ref{l:iscut}}
\cs{existsgaugeisHead}{Lemma~\ref{l:heads}}
\cs{formConstmullintegralle}{Lemma~\ref{l:shells}}
\cs{lintegralformMajorantnetop}{Lemma~\ref{l:fin}}
\cs{normfloorFormle}{Lemma~\ref{l:norm}}
\cs{hasMatSizeofnezero}{Lemma~\ref{l:supp}}
\cs{isCompacthasMatSize}{Lemma~\ref{l:supp}}
\cs{innerbesselCorrelation}{Lemma~\ref{l:corr}}
\cs{continuoussieveVector}{Lemma~\ref{l:cont}}

\begin{document}
\color{ink}

{\raggedright
{\sffamily\bfseries\fontsize{7.2}{9}\selectfont\color{leancol}\addfontfeatures{LetterSpace=8}READBACK NOTE 29\par}
\vspace{5pt}
{\fontsize{21.5}{25.5}\selectfont\bfseries The tower method at a place: floor measures, gauges and the bounds on each floor\par}
\vspace{5pt}
{\fontsize{12.4}{16}\selectfont\itshape\color{inksoft} One unit of the Lean artifact \ArtifactName, read back into
mathematics\par}}
\vspace{8pt}
{\color{rulegrey}\hrule height 0.6pt}
\vspace{7pt}
\thispagestyle{empty}

\textbf{Summary.} A unitary group $\GT{T}$ over a local ring $A$ is cut into \emph{floors}: a tower of subgroups, each with
a measure, whose convolution is a Haar measure. A \emph{cut Bessel function} $W$ on $\GT{T}$ is paired with its
translates by Weyl elements, and the pairing is averaged floor by floor; this gives \emph{floor forms}
$\floorForm{j}{p}{q}{W}{\omega}$, one family for each floor~$j$. The \emph{local theorem} bounds them by one quantity,
the self-norm product, in three situations: by a volume, by Plancherel's inequality against a character of the floor,
and by Plancherel's inequality on a Heisenberg group when the two Weyl elements differ. A descent identity writes
each form as an integral of the forms of the next floor. At the complex place the bounds are summed over shells of
gauges and give the \emph{band limits} that the second moment of Poincar\'e series (Theorem~B of the write-up) needs.

\statement{leancol}{\fontsize{9.6}{13.2}\selectfont\textbf{What this note is.} The artifact \ArtifactName{} is a
body of Lean code whose final theorems are the Riemann hypothesis and the generalized Riemann hypothesis for Dirichlet
$L$-functions. Its authors report that it compiles and that Lean's kernel accepts these theorems from the three
standard axioms; nothing was built for this note. That is a statement about a formal text: no mathematician has read
the argument through, and this note does not say that the Riemann hypothesis is proved. The note treats one part of
one unit of the artifact (the folder of the analysis at the real places): \UnitInNote{} of the unit's \UnitMembers{}
declarations, those of three files; two companion notes treat the others. Proved here in full, from stated inputs:
Theorems A to~D. Theorem~E is proved up to two comparisons of explicit exponents, which are named where they occur.
Taken for granted, with precise statements: the results of the two companion notes (Inputs \Inp{T1}--\Inp{T14} and
\Inp{C1}--\Inp{C4}), definitions of the unit on the second moment, and classical results of the background library
(Section~\ref{s:inputs}). Every statement of a result of this note was compared, argument by argument, with its Lean
text. The Lean text read is \PinText; Appendix~A is the concordance of \ConcDecls{} declarations.}

\sect{The statements}{s:intro}
The objects are defined in Section~\ref{s:setting}; here is what the letters stand for. $A$ is a \emph{place}: a
locally compact ring with involution, with a module $\modA$, an additive character $\psi_A$ and a Plancherel constant
$c_A$; the model cases are $\CC$ and products of non-archimedean local fields. $V$ is a hermitian space over $A$ with
$\varrho$ hyperbolic planes and an anisotropic part, $\GT{T}$ its unitary group, and $\Hl{0}{T}=\GT{T}\supset
\Hl{1}{T}\supset\dots$ the tower of the stabilisers of the first hyperbolic planes. A \emph{level} $\cL$ gives a measure
$\floorMeasure{\cL}{l}$ on each floor $\Hl{l}{T}$ and a measure $\groupMeasure{\cL}$ on $\GT{T}$, their convolution;
a \emph{gauge} $\Gamma$ on $\cL$ measures the size of matrix entries by a filtration $\Gamma_0\subseteq\Gamma_1\subseteq
\dots$ of~$A$ and has a constant $\Lambda(\Gamma)\ge2$. A \emph{head} is a sequence $p\colon\NN\to A^{\times}$ whose
partial products have inverses of size~$1$; $\wl{j}{p}$ is the Weyl element that exchanges the first $j$ hyperbolic
pairs with the entries of~$p$.

\begin{setting}[the local theorem]\label{set:local}
Let $A$, $\mathfrak{s}$, $T$, $\Omega$, $\mathfrak{m}$, $I$ be as in the conventions (Section~\ref{s:setting}), let
$\cL$ be a level in $\Levels{\Omega}{T}{\mathfrak{m}}$ whose measure $\groupMeasure{\cL}$ is a Haar measure of
$\GT{T}$, and let $\Gamma$ be a gauge in $\Gauges{\cL}$. Let $W\colon\GT{T}\to\CC^I$ be a cut Bessel function for
$\cL$: \m{\IsCut{W}}. Let $j$ be a natural number with $j\le\varrho$, let $p,q\colon\NN\to A^{\times}$ be
sequences of units such that $p$ is a head of length $j$ for $\Gamma$, that is \m{\IsHead{j}{p}} for $\Gamma$, and let
$x\in\Posn{\cL}$ be a position.
\end{setting}

\statement{arith}{\phantomsection\label{t:A}\textbf{Theorem A (the local theorem).}\enspace In
Setting~\ref{set:local}, the three quantities on the left are bounded by the self-norm product
\[\begin{aligned}
\selfNormProduct{j}{p}{q}{W}{W'}{x}{\Gamma}={}&\Bigl(\prod_{a<j}\floorModulus{\cL}{a}{p_a^{-1}}\Bigr)^{2}\,
\globalConst{\Gamma}^{3}\;\selfNorm{\cL}{0}{\translatePairing{W}{W}{1}{\cdot}{\cdot}}{x}\\
&\cdot\selfNorm{\cL}{j}{\translatePairing{W}{W'}{d}{\cdot}{\cdot}}{x},\qquad d=\weylRatio{j}{p}{q}.
\end{aligned}\]
\textbf{(a) Volume.} Let $W'\colon\GT{T}\to\CC^I$ satisfy \m{\IsBessel{W'}{\gL{\cL}^{-1}}} and vanish at every $y$
that is not near at~$j$ (if $W'(y)\ne0$ then \m{\IsNear{\cL}{j}{y}}). Let $V\in[0,\infty]$ with $V\ge1$ and
\m{\FloorMassLe{\cL}{j}{V}{1}}, and let $\sigma\in[0,\infty]$. Suppose that, if $\sigma<1$, then for every
$y\in\GT{T}$ with $\ip{W(y)}{W(yx)}\ne0$, \m{\IsNear{\cL}{j}{\wl{j}{p}^{-1}y}} and
\m{\IsNear{\cL}{j}{\wl{j}{q}^{-1}yx}} there is $a$ with $j\le a<\varrho$ and
$\vol\bigl(\floorValues{\cL}{a}{y}{x}\bigr)\le\sigma$. Then
\[\formMajorant{j}{p}{q}{W}{W'}{x}^{2}\ \le\ \sigma^{\volumeExponent{\cL}}\;V^{2}\;\selfNormProduct{j}{p}{q}{W}{W'}{x}{\Gamma}.\]
\textbf{(b) A character on the floor.} Let $\xi\in A^{\times}$ and $\eta\in A$ be such that
$\psi_A(\eta z)=\psi_A(\xi z)$ for all $z\in\Gamma_{200n^2}$, and let $\omega_\eta(y)=\psi_A\bigl(\eta\,\ZT{T}{y}_{jj}\bigr)$
for $y\in\GT{T}$. Then
\[\formModulus{j}{p}{q}{W}{x}{\omega_\eta}^{2}\ \le\ \Bigl(\frac{c_A}{\modA(\xi)}\Bigr)^{|\Omega|}\;\selfNormProduct{j}{p}{q}{W}{W}{x}{\Gamma}.\]
\textbf{(c) Unequal heads.} Let $i$ be a natural number with $j=i+1$, and let $\xi\in A^{\times}$ with
$\xi=\gL{\cL}^{-1}\bigl(p_i^{-1}-q_i^{-1}\bigr)$. Then
\[\formModulus{j}{p}{q}{W}{x}{1}^{2}\ \le\ \Bigl(\frac{c_A}{\modA(\xi)}\Bigr)^{|\Omega|}\;\selfNormProduct{j}{p}{q}{W}{W}{x}{\Gamma}.\]}
\leanline{\lean{formMajorantsqle}\sep\lean{formModulussqleofchar}\sep\lean{formModulussqleofheads}}

So a form is small when the pair of positions is spread out (a), when it is tested against a character of large
frequency~$\xi$ in the diagonal entry of floor~$j$ (b), and when the last entries of the two heads differ by a large
amount (c); in (b) and (c) the gain is the factor $\modA(\xi)^{-|\Omega|}$, where $|\Omega|$ is the number of
anisotropic directions that the floor leaves free.

\begin{setting}[a cut function and a gauge]\label{set:cut}
Let $A$, $\mathfrak{s}$, $T$, $\Omega$, $\mathfrak{m}$, $I$ be as in the conventions, let $\cL$ be a level in
$\Levels{\Omega}{T}{\mathfrak{m}}$, let $\Gamma$ be a gauge in $\Gauges{\cL}$, and let $W\colon\GT{T}\to\CC^I$ be a cut
Bessel function for $\cL$: \m{\IsCut{W}}.
\end{setting}

\statement{arith}{\phantomsection\label{t:B}\textbf{Theorem B (descent).}\enspace In Setting~\ref{set:cut}, let
$p,q\colon\NN\to A^{\times}$, let $a$ be a natural number with $a<\varrho$, let $t\in\CC$ and $\eta\in A$, and suppose
that $\int\formMajorant{a}{p}{q}{W}{W}{x}\,\dd\groupMeasure{\cL}(x)<\infty$. For $c\in A$ let $u(c)=c$ if $c$ is a
unit and $u(c)=1$ otherwise, and let $p[a\mapsto u(c)]$ be the sequence $p$ with its entry $a$ replaced by $u(c)$.
Then the function
\[f(c)=\psi_A(\eta c)\cdot\Bigl(t\cdot\floorForm{a+1}{p[a\mapsto u(c)]}{q[a\mapsto u(c)]}{W}{1}\Bigr)\]
is integrable on $A$, and
\[\int_A f(c)\,\dd c\;=\;t\cdot\floorForm{a}{p}{q}{W}{\omega_\eta},\qquad \omega_\eta(y)=\psi_A\bigl(\eta\,\ZT{T}{y}_{aa}\bigr).\]}
\leanline{\lean{integralfloorFormsucc}}

For $\eta=0$ this says that the form of floor $a$ is the integral of the forms of floor $a+1$ over a common entry at
the place~$a$; for $\eta\ne0$ it identifies the Fourier transform in that entry, which Theorem~A(b) bounds.

\statement{arith}{\phantomsection\label{t:C}\textbf{Theorem C (the self-norm, integrated over positions).}\enspace
Let $A$, $\mathfrak{s}$, $T$, $\Omega$, $\mathfrak{m}$ be as in the conventions, let $\cL$ be a level in
$\Levels{\Omega}{T}{\mathfrak{m}}$ whose measure $\groupMeasure{\cL}$ is a Haar measure, and $\Gamma$ a gauge in
$\Gauges{\cL}$. Let $j$ be a natural number with $j\le\varrho$, let $k\colon\GT{T}\times\GT{T}\to\CC$ and
$d\in\GT{T}$ form a pairing kernel, \m{\IsPairingKernel{\cL}{j}{k}{d}}, and let
\[c=\Bigl(\prod_{l<\varrho}\floorConst{\Gamma}{l}\Bigr)\cdot\iint|k(g,y)|^{2}\,\dd\upperMeasure{\cL}{0}(y)\,\dd\upperMeasure{\cL}{0}(g).\]
Then
\[\int_{\positionsOfSize{\cL}{\Gamma}}\selfNorm{\cL}{j}{k}{x}\;\dd\groupMeasure{\cL}(x)\ \le\ n\cdot c.\]}
\leanline{\lean{lintegralselfNormle}}

\statement{arith}{\phantomsection\label{t:D}\textbf{Theorem D (a bound on a set of positions).}\enspace In
Setting~\ref{set:cut}, suppose that $\groupMeasure{\cL}$ is a Haar measure. Let $\Gamma'$ be a second gauge in
$\Gauges{\cL}$, let $W'\colon\GT{T}\to\CC^I$ satisfy \m{\IsBessel{W'}{\gL{\cL}^{-1}}}, let $j\le\varrho$ and let
$p,q\colon\NN\to A^{\times}$ with $q$ a head of length $j$ for $\Gamma$: \m{\IsHead{j}{q}} for $\Gamma$. Let
$\Theta,e\in[0,\infty]$, let $S\subseteq\GT{T}$ and $f\colon\GT{T}\to[0,\infty]$, and suppose that
\[\begin{gathered}\pairingNorm{W}{W}\cdot\pairingNorm{W}{W'}\le\Theta^2,\qquad S\subseteq\positionsOfSize{\cL}{\Gamma'},\\
f(x)^{2}\le e^{2}\cdot\selfNormProduct{j}{p}{q}{W}{W'}{x}{\Gamma}\ \text{ for all $x\in S$}.\end{gathered}\]
Then
\[\begin{aligned}\headDensity{j}{p}{q}{\cL}\int_{S}f\,\dd\groupMeasure{\cL}\ \le\ {}&\Bigl(e\prod_{l<\varrho}\floorModulus{\cL}{l}{\gL{\cL}}\Bigr)\,
c_A^{\,|\Omega|\varrho}\\ &\cdot\Lambda(\Gamma)^{1500000\,n^{3}}\;\Lambda(\Gamma')^{50000\,n^{2}}\;\Theta\;\Lambda(\Gamma).\end{aligned}\]}
\leanline{\lean{lintegralpositionsle}}

The last theorem is at the complex place. There $\psi_{\CC}(c)=\exp(c-\bar c)$, the level $\cLC{\mathfrak{h}}$ and its
gauges are those of the companion note on the complex place (Inputs \Inp{C1}--\Inp{C3}), and a pair of a Bessel function
and its cut is given by the unit on the second moment (Definition~\ref{d:cutbessel}). The notion of band limits is
recalled in Section~\ref{s:inputs}: it is a system of forms $J_0,\dots,J_{\varrho+1}$ on sequences of non-zero complex
numbers, each the integral of the next over a common entry, with decay in the frequency and off the diagonal.

\statement{arith}{\phantomsection\label{t:E}\textbf{Theorem E (band limits at the complex place).}\enspace Let
$\mathfrak{s}$ be a shape, $I$ a non-empty finite set, $\Phi,\Phi_0\colon\Gone{\mathfrak{s}}\to\CC^I$ continuous
functions, $\mathrm{rep}$ a unitary representation of $\mathrm{U}(\kappa)$ on $\CC^I$, $\rho$ and $b$ real numbers,
$s_0$ and $\mathsf{n}$ natural numbers, and let $\mathfrak{c}$ be cut Bessel data:
\m{\IsCutBessel{\Phi}{\mathrm{rep}}{\rho}{s_0}{b}{\Phi_0}{\mathsf{n}}} (Definition~\ref{d:cutbessel}; $\Phi$ is the
cut, $\Phi_0$ the full function). Let $\deg$ be a real number with $\deg>0$ such that $\Phi_0$ has Schur's bound with
constant $\deg^{-1}|I|$ (Section~\ref{s:inputs}), and suppose that $\deg\cdot\bandc{\mathfrak{s}}\le b^{4r^2}$. Then
there is a family $J=(J_i)_{i\in\NN}$ of functions of two sequences $p,q\colon\NN\to\CC^{\times}$ such that $\Phi$
has the band limits with respect to $b$, $\deg$ and $J$. One such family is
\[J_i(p,q)=\formConst{\mathfrak{c}}\cdot\floorForm{i}{p}{q}{\Phi}{1}\ \ (i\le\varrho),\qquad
J_i(p,q)=\ip{\sieveVector{\Phi}{p}}{\sieveVector{\Phi}{q}}\ \ (i>\varrho),\]
with the floor forms of the level $\cLC{\mathfrak{h}}$ of the sizes $\mathfrak{h}$ that underlie $\mathfrak{c}$.}
\leanline{\lean{hasBandLimitsofisCutBessel}}

\sect{The setting}{s:setting}
\begin{conventions}\phantomsection\label{c:conventions}
$A$ is a commutative ring with an involution $x\mapsto\bar x$, a topology and a Haar measure $\dd x$, and it is a
\emph{place} in the sense of Section~\ref{s:space}. $\mathfrak{s}=(\varrho,\kappa,k_0)$ is a shape, with
$n=\varrho+1$, $k=|\kappa|$ and $r=2\varrho+k+1$. $T$ is an $\mathfrak{s}$-form over~$A$ whose split Gram matrix
$\JT{T}$ is a hermitian unit. $\Omega$ is a finite set~$\Omega$ and $\mathfrak{m}$ a natural number; $\cL$ denotes a level in
$\Levels{\Omega}{T}{\mathfrak{m}}$. $I$ is a finite set~$I$, and $\CC^I$ has the inner product
$\ip{v}{w}=\sum_{i\in I}\overline{v_i}\,w_i$ and the norm $|v|$. In statements at the complex place, $A=\CC$ with complex
conjugation, $T$ is the form of $\Gone{\mathfrak{s}}$, and $\Phi$ denotes a continuous function
$\Gone{\mathfrak{s}}\to\CC^I$. Integrals of functions with values in $[0,\infty]$ are Lebesgue integrals and need no
integrability; a product $0\cdot\infty$ is $0$. Subtraction of natural numbers in exponents is truncated at~$0$.
\end{conventions}

\subhead{The space, its groups, and a place.}\phantomsection\label{s:space} These objects belong to the background
library (units 112, 113 and 119 of the artifact); the companion Note~T (\emph{The tower method at a place: floor
measures, gauges and the bounds on each floor}) defines them in full. A \emph{shape} is a triple
$\mathfrak{s}=(\varrho,\kappa,k_0)$ of a natural number, a finite set and an element $k_0\in\kappa$. The free module
$V=A^{r}$ has the basis $e_0,\dots,e_{\varrho-1}$, $x_i$ ($i\in\kappa$), $z$, $e^*_{\varrho-1},\dots,e^*_0$. An
\emph{$\mathfrak{s}$-form} $T$ is a square matrix on the indices of the $x_i$ and of~$z$; its \emph{split Gram matrix}
$\JT{T}$ pairs $e_a$ with $e^*_a$ (entries~$1$), restricts to $T$ on $V_\varrho=\langle x_i,z\rangle$, and is zero
elsewhere. It is a \emph{hermitian unit} if it is hermitian with invertible determinant. We write
$h_T(v,w)=v^{*}\JT{T}w$, with $v^{*}$ the conjugate transpose, and $\GT{T}$ for the group of the invertible matrices
$g$ with $g^{*}\JT{T}g=\JT{T}$. The \emph{floor} of $e_a$ and of $e^*_a$ is $a$, that of $z$ is $\varrho$, that of
the $x_i$ is $n$; $V_l$ is the span of the basis vectors of floor at least~$l$. The \emph{frame} is
$f_0,\dots,f_\varrho$ with $f_a=e_a$ for $a<\varrho$ and $f_\varrho=z$, and
$\ZT{T}{g}=\bigl(h_T(f_a,gf_b)\bigr)_{0\le a,b\le\varrho}$ is the frame Gram matrix of $g\in\GT{T}$. The subgroups are:
\begin{itemize}
\item $\Hl{l}{T}$, the elements that fix every basis vector of floor less than $l$; so $\Hl{0}{T}=\GT{T}$,
  $\Hl{\varrho}{T}$ is the unitary group of $V_\varrho$, $\Hl{n}{T}$ that of $\langle x_i\rangle$, and
  $\Hl{l}{T}=1$ for $l\ge n+1$;
\item $\Ql{l}{T}$, the elements $g$ with $gf_a-f_a\in\langle e_0,\dots,e_{a-1}\rangle$ for every $a<l$, $a\le\varrho$
  (and, for $l\ge n+1$, also $gx_i-x_i\in\langle e_0,\dots,e_{\varrho-1}\rangle$);
\item $\UT{T}=\Ql{n+1}{T}$, the unipotent radical of the stabiliser of the flag
  $\langle e_0\rangle\subset\langle e_0,e_1\rangle\subset\dots$, and $\Nb{l}{T}$, the unipotent radical of the
  stabiliser of the first $l$ steps of this flag, so that $\Nb{l}{T}=\UT{T}$ for $l\ge\varrho$;
\item $\Nf{l}{T}=\Hl{l}{T}\cap\Nb{l+1}{T}$, a Heisenberg group.
\end{itemize}
They satisfy $\Hl{l}{T}\subseteq\Ql{l}{T}$, $\Nb{l}{T}\subseteq\UT{T}\subseteq\Ql{l'}{T}$ for all $l$, $l'$, the
families $\Hl{l}{T}$ and $\Ql{l}{T}$ decrease with~$l$, $\Hl{l}{T}$ normalises $\Nb{l}{T}$, every element of
$\Ql{l}{T}$ is a product $\gamma\nu$ with $\gamma\in\Hl{l}{T}$ and $\nu\in\Nb{l}{T}$, and for $b\le\varrho$ the
group $\Ql{b+1}{T}\cap\Hl{b}{T}$ is the stabiliser of $f_b$ in $\Hl{b}{T}$. For $u\in\GT{T}$, $\sd(u)$ is the sum of
the entries of $u$ at the places $(f_a,f_{a+1})$, $a<\varrho$. For $p\colon\NN\to A^{\times}$ the \emph{Weyl
element} $\wl{l}{p}\in\GT{T}$ maps $e_a\mapsto p_ae^*_a$ and $e^*_a\mapsto\bar p_a^{-1}e_a$ for $a<\min(l,\varrho)$
and fixes the other basis vectors; it commutes with $\Hl{l}{T}$, and $\wl{l}{\bar p}=\wl{l}{p}^{-1}$. The \emph{floor
Weyl element} $\wo{a}{c}$, for $a<\varrho$ and $c\in A^{\times}$, does the same with the one pair $(e_a,e^*_a)$,
so that $\wl{a+1}{p}=\wo{a}{p_a}\wl{a}{p}$ and $\ZT{T}{\wo{a}{c}}_{aa}=c$.

A \emph{place} is a commutative ring $A$ with involution that is a locally compact Polish topological ring with a Haar
measure, together with: an element $\tau_A$ with $\tau_A+\bar\tau_A=1$; a unit $\delta_A$ with
$\bar\delta_A=-\delta_A$; a Haar measure on the skew elements; a measurable multiplicative \emph{module}
$\modA\colon A\to[0,\infty]$, invariant under $x\mapsto\bar x$ and $x\mapsto-x$, such that multiplication by a matrix
$M$ with invertible determinant multiplies Haar measure on $A^{N}$ by $\modA(\det M)^{-1}$; a continuous additive
character $\psi_A\colon A\to\CC$ of absolute value~$1$; a constant $c_A\in[0,\infty]$; and further axioms (almost
every element is a unit; a compatibility of the measure on skew elements; a local bound for $\modA^{-1}$; a
Plancherel kernel) that enter this note only through the inputs. For $A=\CC$: $\modA(x)=|x|^2$,
$\psi_{\CC}(c)=\exp(c-\bar c)$, $c_{\CC}=\pi^{2}$.

\subhead{The tower (recalled from Note~T).}\phantomsection\label{s:tower} A \emph{level}
$\cL\in\Levels{\Omega}{T}{\mathfrak{m}}$ consists of: a \emph{hyperbolic basis} of $V_\varrho$, that is a basis matrix
$E(\cL)$ made of $\mathfrak{m}+1$ hyperbolic pairs and of orthonormal vectors indexed by $\Omega$ (so
$k+1=|\Omega|+2\mathfrak{m}+2$); two $s$-finite measures on $\GT{T}$, written $\floorMeasure{\cL}{\varrho}$ and
$\floorMeasure{\cL}{n}$, the second non-zero; a unit $\gL{\cL}=\overline{\gL{\cL}}$ of $A$; an open set
$\Ball{\cL}\subseteq V$ that contains $0$ and has compact closure; a compact set $\BoxL{\cL}\subseteq\GT{T}$, the
\emph{box}; and a measurable set $\Posn{\cL}\subseteq\GT{T}$ of \emph{positions}. It is required that
$\floorMeasure{\cL}{\varrho}$ is carried by $\Hl{\varrho}{T}$ and that its image under $\sigma\mapsto\sigma^{-1}z$ is
the quadric measure $\quadric{\cL}{\varrho}$ of Note~T on $\{v\in V_\varrho:h_T(v,v)=h_T(z,z)\}$; and that
$\floorMeasure{\cL}{n}$ is carried by $\Hl{n}{T}$, is invariant under right translations by $\Hl{n}{T}$, and gives
mass at most $1$ to $\{h:u\,h\,g\in\BoxL{\cL}\text{ for some }u\in\UT{T}\}$, for every $g$. For $l<\varrho$ the
\emph{floor measure} is $\floorMeasure{\cL}{l}=\varpi_l*\eta_l$, where $*$ is convolution of measures on $\GT{T}$,
$\varpi_l$ is the image of the measure $\modA(c)^{k+1+2(\varrho-l-1)}\dd c$ on $A$ under $c\mapsto\wo{l}{c}^{-1}$,
and $\eta_l$ is the \emph{Heisenberg measure}, a Haar measure of $\Nf{l}{T}$ in the coordinates of Note~T; and
$\floorMeasure{\cL}{l}=\delta_1$, the unit mass at~$1$, for $l>n$. We put $\eta_l=\delta_1$ for $l\ge\varrho$ and
$\etab{j}=\eta_0*\dots*\eta_{j-1}$. Then
\[\begin{gathered}
\stackMeasure{\cL}_{a,i}=(\eta_{a+i-1}*\floorMeasure{\cL}{a+i-1})*\dots*(\eta_a*\floorMeasure{\cL}{a}),\\
\groupMeasure{\cL}=\stackMeasure{\cL}_{0,\varrho+2},\qquad
\upperMeasure{\cL}{l}=\floorMeasure{\cL}{n}*\dots*\floorMeasure{\cL}{l},
\end{gathered}\]
with $\upperMeasure{\cL}{l}=\delta_1$ for $l>n$. The measure $\floorMeasure{\cL}{l}$ is carried by $\Hl{l}{T}$, and
$\groupMeasure{\cL}$ is right invariant.

A point $g\in\GT{T}$ is \emph{near at $j$}, \m{\IsNear{\cL}{j}{g}}, if $\xi g\in\BoxL{\cL}$ for some
$\xi\in\Ql{j}{T}$. For $a<\varrho$ and $x,y\in\GT{T}$, the set $\floorValues{\cL}{a}{x}{y}\subseteq A$ consists of the
coordinates at $e^*_a$ of the vectors $w\in V$ whose coordinates at $e^*_b$, $b<a$, vanish and for which
$u_1=x^{-1}w$ and $u_2=(xy)^{-1}w$ satisfy $u_1\in\Ball{\cL}$ and $\gL{\cL}^{-1}(u_1-u_2)\in\Ball{\cL}$. With the
\emph{volume exponent} $\volumeExponent{\cL}=|\Omega|-(3n+1+\mathfrak{m})$, the \emph{weight} and the \emph{floor
module} are
\[\begin{gathered}
\floorWeight{\cL}{j}{x}{y}=\prod_{a<\min(j,\varrho)}\vol\bigl(\floorValues{\cL}{a}{x}{y}\bigr)^{\volumeExponent{\cL}},\\
\floorModulus{\cL}{l}{g}=\modA(g)^{k+1+2(\varrho-(l+1))}\,\modA(g)^{1/2}\quad(g\in A^{\times}),
\end{gathered}\]
and $\invWeight{\cL}{j}{x}$ is the function $z\mapsto\floorWeight{\cL}{j}{z}{x}^{-1}$ on the points near at $j$,
extended by~$0$. For a function $c\colon\GT{T}\to[0,\infty]$ and $K\in[0,\infty]$, \m{\FloorMassLe{\cL}{j}{K}{c}}
means that $\int c(\sigma g)\,\dd\floorMeasure{\cL}{j}(\sigma)\le K$ for every $g$, the integral taken over the
$\sigma$ with $\sigma g$ near at $j+1$.

\subhead{Kernels averaged over floors (recalled from Note~T).} For a measurable
$k\colon\GT{T}\times\GT{T}\to\CC$ the \emph{averaged kernel}, its \emph{pivot}, the \emph{weighted square} and the
\emph{weighted square sum} at a point $x\in\GT{T}$ are
\begin{gather*}
\upperPairing{\cL}{j}{k}{g}{y}=\int k(\sigma g,\sigma y)\,\dd\upperMeasure{\cL}{j}(\sigma),\\
\pivotModulus{\cL}{j}{k}{g}{y}=\int\bigl|\upperPairing{\cL}{j+1}{k}{\sigma g}{\sigma y}\bigr|\,\dd\floorMeasure{\cL}{j}(\sigma),\\
\weightedSquare{\cL}{j}{k}{x}{g}=\floorWeight{\cL}{j}{g}{x}\cdot\bigl|\upperPairing{\cL}{j+1}{k}{g}{gx}\bigr|^{2},\\
\weightedSquareSum{\cL}{k}{x}{j}{g}=\int\Bigl(\weightedSquare{\cL}{j}{k}{x}{\sigma g}+\weightedSquareSum{\cL}{k}{x}{j+1}{\sigma g}\Bigr)\dd\floorMeasure{\cL}{j}(\sigma)
\end{gather*}
for $j\le\varrho$, with $\weightedSquareSum{\cL}{k}{x}{j}{g}=0$ for $j>\varrho$,
and the \emph{self-norm} is
$\selfNorm{\cL}{j}{k}{x}=\int\weightedSquareSum{\cL}{k}{x}{j}{g}\,\dd(\floorMeasure{\cL}{j-1}*\dots*\floorMeasure{\cL}{0})(g)$.
A pair of such a $k$ and an element $w\in\GT{T}$ is a \emph{pairing kernel at $j$}, \m{\IsPairingKernel{\cL}{j}{k}{w}},
if: $k$ is measurable; for all $l$, all $\nu,\nu'\in\Nb{l}{T}$ and all $g,y$,
\[\upperPairing{\cL}{l+1}{k}{\nu g}{\nu'y}=\psi_A\Bigl(\gL{\cL}^{-1}\bigl(\sd(w\nu'w^{-1})-\sd(\nu)\bigr)\Bigr)\,\upperPairing{\cL}{l+1}{k}{g}{y};\]
$w$ commutes with $\Hl{j}{T}$; for every $l$ the integrals $\int|k(\sigma g,\sigma y)|\,\dd\upperMeasure{\cL}{l}(\sigma)$
are bounded by a finite constant independent of $g$, $y$; and $\upperPairing{\cL}{l}{k}{g}{y}\ne0$ implies
\m{\IsNear{\cL}{l}{g}}.

\subhead{Sizes and gauges (recalled from Note~T).} A \emph{gauge} $\Gamma\in\Gauges{\cL}$ consists of sets
$\Gamma_0,\Gamma_1,\dots\subseteq A$, compact, containing $0$ and $1$, stable under $x\mapsto-x$ and $x\mapsto\bar x$,
with $\Gamma_i\Gamma_{i'}\subseteq\Gamma_{i+i'}$ and with sums of at most $r+1$ elements of $\Gamma_i$ in
$\Gamma_{i+1}$, and of a constant $\Lambda(\Gamma)\in[2,\infty)$, subject to the axioms listed in Note~T. Those used
here by name: the entries of $\JT{T}$ and $\JT{T}^{-1}$ lie in $\Gamma_1$; $\mathfrak{m}\le9n$ and $c_A\ge1$;
$\gL{\cL}\in\Gamma_1$; the entries of the elements of $\BoxL{\cL}$ lie in $\Gamma_1$; the coordinates of the vectors
of $\Ball{\cL}$ lie in $\Gamma_1$; $\modA(u)\le\Lambda(\Gamma)^{i}$ for units $u\in\Gamma_i$; and
$\vol(\Gamma_i)\le\Lambda(\Gamma)^{i+1}$. A matrix $X$ over $A$ \emph{has size at most $i$ for $\Gamma$},
\m{\HasMatSize{i}{X}}, if all its entries lie in $\Gamma_i$; for $g\in\GT{T}$ this refers to the matrix of~$g$;
$\vecsOfSize{i}$ is the set of the vectors with all coordinates in $\Gamma_i$. A sequence $p\colon\NN\to A^{\times}$
is a \emph{head of length $j$ for $\Gamma$}, \m{\IsHead{j}{p}}, if $(p_0\cdots p_{l-1})^{-1}\in\Gamma_1$ for all
$l\le j$; then $\bar p$ is one too. The predicates $\operatorname{sz}$, $\operatorname{Hd}$ and the sets
$\mathrm{Vec}$ refer to a gauge that the Lean text takes as an implicit argument; every statement below names it.
Further,
\[\begin{gathered}
\positionsOfSize{\cL}{\Gamma}=\bigl\{x\in\Posn{\cL}:\HasMatSize{1}{\gL{\cL}^{-1}(x-1)}\ \text{for }\Gamma\bigr\},\\
\floorConst{\Gamma}{l}=\floorModulus{\cL}{l}{\gL{\cL}}\,c_A^{|\Omega|}\,\Lambda(\Gamma)^{4600n},\qquad
\globalConst{\Gamma}=\Lambda(\Gamma)^{10^{6}n^{3}}.
\end{gathered}\]

\subhead{Bessel functions and the kernels of two functions.}
\begin{definition}[Bessel functions]\label{d:bessel}
Let $W\colon\GT{T}\to\CC^I$ and $\eta\in A^{\times}$. Then $W$ is a \emph{Bessel function of frequency $\eta$},
\m{\IsBessel{W}{\eta}}, if $W$ is measurable and bounded and
\[W(ux)=\psi_A\bigl(\eta\,\sd(u)\bigr)\,W(x)\qquad\text{for all }u\in\UT{T},\ x\in\GT{T}.\]
\end{definition}
\leanline{\lean{IsBessel}}
\begin{definition}[cut functions]\label{d:cut}
Let $\cL$ be a level in $\Levels{\Omega}{T}{\mathfrak{m}}$ and $W\colon\GT{T}\to\CC^I$. Then $W$ is \emph{cut for
$\cL$}, \m{\IsCut{W}}, if \m{\IsBessel{W}{\gL{\cL}^{-1}}} and, for every $x$ with $W(x)\ne0$, there is
$u\in\UT{T}$ with $ux\in\BoxL{\cL}$. The level is an implicit argument of this predicate in the Lean text.
\end{definition}
\leanline{\lean{IsCut}\sep\lean{isNearofnezero}}
Since $\UT{T}\subseteq\Ql{j}{T}$ for every $j$, a cut function vanishes at every point that is not near at~$j$: if
$W(x)\ne0$ then \m{\IsNear{\cL}{j}{x}}, for every~$j$.
\begin{definition}[the kernel of two functions]\label{d:kernel}
For $W,W'\colon\GT{T}\to\CC^I$ and $d,x,y\in\GT{T}$, $\translatePairing{W}{W'}{d}{x}{y}$ is the inner product of
$W(x)$ and $W'(dy)$ in $\CC^I$. As a function of $(x,y)$ it is written $\translatePairing{W}{W'}{d}{\cdot}{\cdot}$.
\end{definition}
\leanline{\lean{translatePairing}}
\begin{definition}[pairing norm]\label{d:pn}
For a level $\cL$ in $\Levels{\Omega}{T}{\mathfrak{m}}$ (an implicit argument) and $W,W'\colon\GT{T}\to\CC^I$,
\[\pairingNorm{W}{W'}=\iint\bigl|\ip{W(x)}{W'(y)}\bigr|^{2}\,\dd\upperMeasure{\cL}{0}(y)\,\dd\upperMeasure{\cL}{0}(x)\ \in[0,\infty].\]
\end{definition}
\leanline{\lean{pairingNorm}}

\subhead{The floor forms.}
\begin{definition}[head density]\label{d:hd}
For a natural number $j$, sequences $p,q\colon\NN\to A^{\times}$ and a level $\cL$ in
$\Levels{\Omega}{T}{\mathfrak{m}}$,
\[\headDensity{j}{p}{q}{\cL}=\prod_{a<j}\Bigl(\modA(p_aq_a)^{1/2}\Bigr)^{k+1+2(\varrho-(a+1))}\ \in[0,\infty].\]
It depends on $\cL$ only through $A$ and $\mathfrak{s}$.
\end{definition}
\leanline{\lean{headDensity}}
\begin{definition}[the integrands]\label{d:integrand}
Let $\cL$ be a level in $\Levels{\Omega}{T}{\mathfrak{m}}$ (an implicit argument of the three notions), $j$ a natural
number, $W\colon\GT{T}\to\CC^I$ and $x,d,w,g\in\GT{T}$. Put
\[
\formIntegrandOne{j}{W}{x}{d}{w}{g}=\upperPairing{\cL}{j}{\translatePairing{W}{W}{1}{\cdot}{\cdot}}{g}{gx}\cdot
\overline{\upperPairing{\cL}{j}{\translatePairing{W}{W}{d}{\cdot}{\cdot}}{w^{-1}g}{w^{-1}gx}}.
\]
For $p,q\colon\NN\to A^{\times}$ and $\omega\colon\GT{T}\to\CC$ put, with $d=\weylRatio{j}{p}{q}$,
\[
\formIntegrand{j}{p}{q}{W}{x}{\omega}{g}=\iint\omega(\sigma\tau^{-1})\;
\formIntegrandOne{j+1}{W}{x}{d}{\sigma\wl{j}{p}\tau^{-1}}{\sigma g}\;\dd\floorMeasure{\cL}{j}(\tau)\,\dd\floorMeasure{\cL}{j}(\sigma).
\]
(Integrals of complex functions are Bochner integrals, equal to $0$ when the function is not integrable.)
\end{definition}
\leanline{\lean{formIntegrandOne}\sep\lean{formIntegrand}}
\begin{definition}[floor form, modulus, majorant]\label{d:form}
In the situation of Definition~\ref{d:integrand}, the \emph{floor form} of floor $j$ is the complex number
\[
\floorForm{j}{p}{q}{W}{\omega}=\headDensity{j}{p}{q}{\cL}\cdot\iint\formIntegrand{j}{p}{q}{W}{x}{\omega}{g}\;
\dd\stackMeasure{\cL}_{0,j}(g)\,\dd\groupMeasure{\cL}(x)
\]
(the head density is read as $0$ if it is infinite). Its \emph{modulus} at a point $x\in\GT{T}$ is
\[\formModulus{j}{p}{q}{W}{x}{\omega}=\int\bigl|\formIntegrand{j}{p}{q}{W}{x}{\omega}{g}\bigr|\,\dd\stackMeasure{\cL}_{0,j}(g)\ \in[0,\infty],\]
and for a second function $W'\colon\GT{T}\to\CC^I$ the \emph{majorant} at $x$ is, with $d=\weylRatio{j}{p}{q}$,
\[\begin{aligned}
\formMajorant{j}{p}{q}{W}{W'}{x}=\int{}&\pivotModulus{\cL}{j}{\translatePairing{W}{W}{1}{\cdot}{\cdot}}{g}{gx}\\
&\cdot\pivotModulus{\cL}{j}{\translatePairing{W}{W'}{d}{\cdot}{\cdot}}{\wl{j}{p}^{-1}g}{\wl{j}{p}^{-1}gx}\;\dd\stackMeasure{\cL}_{0,j}(g).
\end{aligned}\]
\end{definition}
\leanline{\lean{floorForm}\sep\lean{formModulus}\sep\lean{formMajorant}}
\begin{definition}[self-norm product]\label{d:snp}
For a level $\cL$ in $\Levels{\Omega}{T}{\mathfrak{m}}$ (an implicit argument), a gauge $\Gamma$ in $\Gauges{\cL}$, a
natural number $j$, sequences $p,q\colon\NN\to A^{\times}$, functions $W,W'\colon\GT{T}\to\CC^I$ and $x\in\GT{T}$,
with $d=\weylRatio{j}{p}{q}$,
\[\begin{aligned}
\selfNormProduct{j}{p}{q}{W}{W'}{x}{\Gamma}={}&\Bigl(\prod_{a<j}\floorModulus{\cL}{a}{p_a^{-1}}\Bigr)^{2}\,
\globalConst{\Gamma}^{3}\;\selfNorm{\cL}{0}{\translatePairing{W}{W}{1}{\cdot}{\cdot}}{x}\\
&\cdot\selfNorm{\cL}{j}{\translatePairing{W}{W'}{d}{\cdot}{\cdot}}{x}\ \in[0,\infty].
\end{aligned}\]
\end{definition}
\leanline{\lean{selfNormProduct}}

\subhead{The complex place.} $\Gone{\mathfrak{s}}$ is $\GT{T}$ for $A=\CC$ and the form $T$ that is the identity on
the $x_i$ and $-1$ at $z$; its Haar measure is written $\dd y$. \emph{Real-place sizes} $\mathfrak{h}$ in
$\RealPlaceSizes{\mathfrak{s}}{\rho}{s_0}{b}$ are the numerical hypotheses of Input~\Inp{C1}; they give the level
$\cLC{\mathfrak{h}}$ and, for every real $R$ with $\rho\le R$, the gauge $\gauge{\mathfrak{h}}{\rho\le R}$ of
Input~\Inp{C3}, and the constant $\haarRatio{\mathfrak{h}}$ of Input~\Inp{C2}.
\begin{definition}[cut Bessel data]\label{d:cutbessel}
Let $\Phi,\Phi_0\colon\Gone{\mathfrak{s}}\to\CC^I$ be continuous, $\mathrm{rep}$ a homomorphism from
$\mathrm{U}(\kappa)$ to the unitary group of $\CC^I$, $\rho$ and $b$ real numbers, $s_0$ and $\mathsf{n}$ natural
numbers. \emph{Cut Bessel data} $\mathfrak{c}$, \m{\IsCutBessel{\Phi}{\mathrm{rep}}{\rho}{s_0}{b}{\Phi_0}{\mathsf{n}}},
are: a proof that $(\Phi_0,\Phi)$ is a cut pair with respect to $b$, $\mathsf{n}$ and $\mathrm{rep}$, with full
function $\Phi_0$ and cut $\Phi$ (Section~\ref{s:inputs}); and real-place sizes $\mathfrak{h}$ in
$\RealPlaceSizes{\mathfrak{s}}{\rho}{s_0}{b}$, the sizes that \emph{underlie} $\mathfrak{c}$.
\end{definition}
\leanline{\lean{IsCutBessel}\sep\lean{volumeExponenteq}}
For such data, $\volumeExponent{\cLC{\mathfrak{h}}}=s_0$: the level has $\mathfrak{m}=0$ and $|\Omega|=k-1$, and
$k=s_0+3n+2$.
\begin{definition}[the constant of the forms]\label{d:formconst}
Let $\Phi,\Phi_0\colon\Gone{\mathfrak{s}}\to\CC^I$, $\mathrm{rep}$, $\rho$, $b$, $s_0$ and $\mathsf{n}$ be as in
Definition~\ref{d:cutbessel}, and $\mathfrak{c}$ cut Bessel data,
\m{\IsCutBessel{\Phi}{\mathrm{rep}}{\rho}{s_0}{b}{\Phi_0}{\mathsf{n}}}, with underlying sizes $\mathfrak{h}$. Then
\[\formConst{\mathfrak{c}}=\frac{\volU{k}}{\volU{r}}\cdot\frac{1}{\bigl(\haarRatio{\mathfrak{h}}\,|I|\bigr)^{2}},
\qquad \volU{N}=\prod_{i<N}\frac{2\pi^{i+1}}{i!}.\]
\end{definition}
\leanline{\lean{formConst}}
\begin{definition}[the vector of the top floor]\label{d:sieve}
For a continuous $\Phi\colon\Gone{\mathfrak{s}}\to\CC^I$ and $p\colon\NN\to\CC^{\times}$,
\[
\sieveVector{\Phi}{p}=\frac{\sqrt{\cellc{\mathfrak{s}}}\,\sqrt{|p_\varrho|^{2}-1}^{\,k-1}}{|I|}
\prod_{a<\varrho}|p_a|^{k+1+2(\varrho-(a+1))}\cdot\corr{\dd y}{\Phi}{\wl{\varrho}{p}\,\boost{\varrho}{p_\varrho}}\ \in\CC^{I\times I},
\]
where the square root of a negative number is $0$, $\cellc{\mathfrak{s}}$, $\corr{\dd y}{\Phi}{\gamma}$ and the boost
$\boost{\varrho}{c}$ are those of Section~\ref{s:inputs}.
\end{definition}
\leanline{\lean{sieveVector}}
\sect{What is taken for granted}{s:inputs}
\subhead{Inputs from Note T}\phantomsection\label{s:inpT}(same unit; \emph{The tower method at a
place: floor measures, gauges and the bounds on each floor}). $\cL$ is a level in $\Levels{\Omega}{T}{\mathfrak{m}}$,
$\Gamma$ a gauge in $\Gauges{\cL}$; functions called $\Phi$, $f$ take values in $[0,\infty]$ and are measurable;
``invariant'' means invariant under left translations.

\Inp{T1} $\stackMeasure{\cL}_{0,j}=\etab{j}*(\floorMeasure{\cL}{j-1}*\dots*\floorMeasure{\cL}{0})$; $\etab{j}$ is
carried by $\Nb{j}{T}$, is invariant under translations by $\Nb{j}{T}$ on both sides, and
$\delta_s*\etab{j}=\etab{j}*\delta_s$ for $s\in\Hl{j}{T}$.

\Inp{T2} If $\groupMeasure{\cL}$ is a Haar measure, $j\le\varrho+2$ and $\Phi$ is $\Hl{j}{T}$-invariant, then
\[\int\Phi(\wl{j}{p}^{-1}g)\,\dd\stackMeasure{\cL}_{0,j}(g)=\int\Phi\,\dd\stackMeasure{\cL}_{0,j}.\]

\Inp{T3} If $\groupMeasure{\cL}=\rho_1*\nu$ is a Haar measure, with $\rho_1,\nu$ $s$-finite, $\rho_1$ carried by a
subgroup $S$ and invariant under right translations by $S$, and $\delta_s*\rho_1=\rho_1*\delta_s$, then
$\int\Phi(sg)\,\dd\nu(g)=\int\Phi\,\dd\nu$ for every $S$-invariant $\Phi$.

\Inp{T4} If a subgroup $S$ contains $\Hl{l}{T}$, $f$ is $(\Ql{l+1}{T}\cap S)$-invariant and $s\in\Ql{l}{T}\cap S$,
then $\int f(\sigma s)\,\dd\floorMeasure{\cL}{l}(\sigma)=\int f\,\dd\floorMeasure{\cL}{l}$.

\Inp{T5} For $j\le\varrho$:
\[\int\Phi(\weylRatio{j}{p}{q}\,x)\,\dd\upperMeasure{\cL}{0}(x)=\prod_{a<j}\modA(p_a/q_a)^{k+2(\varrho-a)}\int\Phi\,\dd\upperMeasure{\cL}{0}.\]

\Inp{T6} For $b\le\varrho$ and $s\in\Hl{b}{T}$, the map $\sigma\mapsto s\sigma^{-1}f_b$ carries
$\floorMeasure{\cL}{b}$ to the quadric measure $\quadric{\cL}{b}$ on $V_b$.

\Inp{T7} Let \m{\IsPairingKernel{\cL}{j}{k}{d}}. (a) $\sigma\mapsto\upperPairing{\cL}{l+1}{k}{\sigma g}{\sigma y}$ is
$\floorMeasure{\cL}{l}$-integrable with integral $\upperPairing{\cL}{l}{k}{g}{y}$; hence
$|\upperPairing{\cL}{l}{k}{g}{y}|\le\pivotModulus{\cL}{l}{k}{g}{y}$. (b) If $j\le l$ and
$s\in\Ql{l}{T}\cap\Hl{j}{T}$, then $\upperPairing{\cL}{l}{k}{sg}{sy}=\upperPairing{\cL}{l}{k}{g}{y}$. (c)
$\weightedSquare{\cL}{j}{k}{x}{\cdot}$ is $\Ql{j+1}{T}$-invariant and vanishes off the points near at $j+1$;
$\weightedSquareSum{\cL}{k}{x}{j}{\cdot}$ is $\Ql{j}{T}$-invariant and vanishes off the points near at $j$. (d) The
pair is a pairing kernel at every $l\ge j$. (e) \emph{Cauchy--Schwarz on a floor}: if \m{\FloorMassLe{\cL}{l}{K}{u}}
and $|\upperPairing{\cL}{l+1}{k}{\sigma g}{\sigma y}|^{2}\le u(\sigma g)\,v(\sigma)$ for $\sigma\in\Hl{l}{T}$, then
$\pivotModulus{\cL}{l}{k}{g}{y}^{2}\le K\int v\,\dd\floorMeasure{\cL}{l}$. (f) For $l\le\varrho$,
$\selfNorm{\cL}{l}{k}{x}=\int\weightedSquare{\cL}{l}{k}{x}{g}\,\dd(\floorMeasure{\cL}{l}*\dots*\floorMeasure{\cL}{0})(g)+\selfNorm{\cL}{l+1}{k}{x}$,
and $\selfNorm{\cL}{l}{k}{x}\le\selfNorm{\cL}{0}{k}{x}$ for $l\le n$. For every measurable $k$: (g) if
$\pivotModulus{\cL}{l}{k}{g}{y}\ne0$ then $k(\sigma g,\sigma y)\ne0$ for some $\sigma\in\Hl{l}{T}$; (h) for
$l\le\varrho$, $\int|\upperPairing{\cL}{l+1}{k}{\sigma g}{\sigma gx}|^{2}\,\dd\floorMeasure{\cL}{l}(\sigma)\le
\weightedSquareSum{\cL}{k}{x}{l}{g}/\floorWeight{\cL}{l}{g}{x}$.

\Inp{T8} Nearness at $j$ is $\Ql{j}{T}$-invariant; if $xg$ is near at $i$, $x\in\Hl{j}{T}$ and $j\le i$, then $g$ is
near at~$j$. For $\xi\in\Ql{a+1}{T}$, $\floorValues{\cL}{a}{\xi x}{y}=\floorValues{\cL}{a}{x}{y}$; for
$\xi\in\Ql{j}{T}$, $\floorWeight{\cL}{j}{\xi x}{y}=\floorWeight{\cL}{j}{x}{y}$ and
$\invWeight{\cL}{j}{x}(\xi g)=\invWeight{\cL}{j}{x}(g)$; $\floorWeight{\cL}{a+1}{x}{y}=\floorWeight{\cL}{a}{x}{y}\,
\vol(\floorValues{\cL}{a}{x}{y})^{\volumeExponent{\cL}}$ for $a<\varrho$; weights are neither $0$ nor $\infty$.

\Inp{T9} (a) For $b\le\varrho$: \m{\FloorMassLe{\cL}{b}{\Lambda(\Gamma)^{1287n}\vol(\Gamma_{11})^{|\Omega|}}{1}}. (b)
If \m{\FloorMassLe{\cL}{a}{V}{1}} for all $a\le\varrho$, then
$\upperMeasure{\cL}{l}\{g:gc\text{ is near at }\varrho+2\}\le V^{n-l}$ for all $l$, $c$. (c) For $x\in\Posn{\cL}$ and
$a<\varrho$: \m{\FloorMassLe{\cL}{a}{\Lambda(\Gamma)^{10^{4}n}}{u}} with
$u(g)=\floorWeight{\cL}{a}{g}{x}/\floorWeight{\cL}{a+1}{g}{x}$.

\Inp{T10} For $j\le\varrho$, $g$ near at $j$ and every $y$:
$\int\mathbf{1}_{\positionsOfSize{\cL}{\Gamma}}(g^{-1}\nu y)\,\floorWeight{\cL}{j}{g}{g^{-1}\nu y}\,\dd\etab{j}(\nu)\le\prod_{l<j}\floorConst{\Gamma}{l}$.

\Inp{T11} Let $j\le\varrho$, \m{\IsHead{j}{p}} for $\Gamma$, $x\in\Posn{\cL}$, $j\le b$, $\rho_1$ an $s$-finite
measure, and $\Phi$ $\Ql{b}{T}$-invariant and zero off the points near at $b$. Then
$\int\Phi(g)\,\invWeight{\cL}{j}{x}(\wl{j}{p}g)\,\dd(\etab{j}*\rho_1)(g)\le\prod_{a<j}\floorModulus{\cL}{a}{p_a^{-1}}\cdot\globalConst{\Gamma}\int\Phi\,\dd\rho_1$.

\Inp{T12} \emph{Plancherel on a floor.} Let $b\le\varrho$, $t\in\Hl{b}{T}$, $x_1,y_1\in\GT{T}$ with
\m{\HasMatSize{i_1}{tx_1}} for $\Gamma$, \m{\IsPairingKernel{\cL}{b}{k}{d}}, and
$F(\sigma)=\upperPairing{\cL}{b+1}{k}{\sigma x_1}{\sigma y_1}$. Then (i) for $\floorMeasure{\cL}{b}$-almost every
$\sigma$ with $F(\sigma)\ne0$, $t\sigma^{-1}f_b\in\vecsOfSize{i_1+6}$; (ii) for every $\xi\in A^{\times}$,
$y'\in V$ and every measure $\mu$ on $V$ with \m{\HasFibreBound{\Gamma}{b}{\mu}{i}{M}},
\[\begin{aligned}\int_{\vecsOfSize{i}}\Bigl|\int\psi_A\bigl(\xi\,h_T(v+y',t\sigma^{-1}f_b)\bigr)F(\sigma)\,\dd\floorMeasure{\cL}{b}(\sigma)\Bigr|^{2}&\dd\mu(v)\\
\le M\,\Lambda(\Gamma)^{99n(i_1+12)}\Bigl(\frac{c_A}{\modA(\xi)}\Bigr)^{|\Omega|}&\int|F|^{2}\,\dd\floorMeasure{\cL}{b};\end{aligned}\]
and \m{\HasFibreBound{\Gamma}{b}{\quadric{\cL}{b}}{i}{\Lambda(\Gamma)^{99n(i+6)}}} for all $b\le\varrho$ and $i$.

\Inp{T13} \emph{Plancherel on a Heisenberg group.} Let $a<\varrho$, $\xi,u\in A^{\times}$, $\varepsilon\in\NN$,
$C\in[0,\infty]$, \m{\IsPairingKernel{\cL}{a+1}{k}{d}}, $x_1,y_1\in\GT{T}$, $S\subseteq\GT{T}$,
$\Phi\colon\GT{T}\to[0,\infty]$, and natural numbers $i_1$, $i_2$. Suppose: for $s\in\Nf{a}{T}$,
$\gL{\cL}^{-1}(\sd(dsd^{-1})-\sd(s^{\varepsilon}))=\xi\,\sd(s)$; for $\nu\in\Nf{a}{T}$ with $\Phi(\nu)\ne0$,
$\nu\in S$ and $\Phi(\nu)\le C\,|\upperPairing{\cL}{a+1}{k}{\nu^{\varepsilon}x_1}{\nu y_1}|^{2}$; for every $m\in S$
there are $t\in\Ql{a+1}{T}\cap\Hl{a}{T}$ and $q\in\Nb{a+1}{T}$ with \m{\HasMatSize{i_1}{tqx_1}} such that for
every $\nu\in S$ some $\nu'\in\Nb{a}{T}$ makes the column at $e^*_a$ of $u^{-1}(\nu'tm(t\nu)^{-1}-1)$ lie in
$\vecsOfSize{i_2}$ (sizes for $\Gamma$). Then
\[\begin{aligned}\int\Phi\,\dd\eta_a\le{}& C\Bigl(\frac{c_A}{\modA(\xi u)}\Bigr)^{|\Omega|}\floorModulus{\cL}{a}{u}\,\Lambda(\Gamma)^{100n(i_1+i_2+31)}\\
&\cdot\int\bigl|\upperPairing{\cL}{a+2}{k}{\sigma x_1}{\sigma y_1}\bigr|^{2}\,\dd\floorMeasure{\cL}{a+1}(\sigma).\end{aligned}\]

\Inp{T14} \emph{Sizes.} Sizes are for $\Gamma$; $j\le\varrho$. (a) If \m{\IsHead{j}{p}}, $s\in\Nb{j}{T}$,
$sg\in\BoxL{\cL}$ and $\wl{j}{p}^{-1}g$ is near at $j$, then \m{\HasMatSize{60n^2}{g}}. (b) If \m{\IsHead{j}{p}}
and $g$, $\wl{j}{p}^{-1}g$ are near at $j$, there is $\sigma\in\Hl{j}{T}$ with \m{\HasMatSize{60n^2}{\sigma g}} and
\m{\HasMatSize{99n^2}{\sigma\wl{j}{p}^{-1}g}}. (c) If $s\in\Ql{j}{T}$, $sg\in\BoxL{\cL}$ and $\wl{j}{p}^{-1}g$ is
near at $j$, then $p_0\cdots p_{l-1}\in\Gamma_{9n}$ for $l\le j$; if moreover \m{\IsHead{j}{p}} and $a<j$, then
$p_a\in\Gamma_{10n}$. (d) If $g\in\BoxL{\cL}$ and $x,y\in\positionsOfSize{\cL}{\Gamma}$, then
\m{\HasMatSize{14}{\gL{\cL}^{-1}(gx(gy)^{-1}-1)}}. (e) Let $\Gamma_S$ be a gauge in $\Gauges{\cL}$ for which $p,q$
are heads of length $j$, with $\gL{\cL}^{-1}c\in\Gamma_1$ for all $c\in(\Gamma_S)_{10n}$ if $j\ne0$. If
$ux,v(xy)\in\BoxL{\cL}$ for some $u,v\in\UT{T}$ with \m{\HasMatSize{2}{\gL{\cL}^{-1}(v(xy)-ux)}}, if
$\wl{j}{p}^{-1}x$ and $\wl{j}{q}^{-1}xy$ are near at~$j$, and if for every $a$ with $j\le a<\varrho$ the set
$\floorValues{\cL}{a}{x}{y}$ contains a unit $c$ with $c^{-1}\in\Gamma_1$, then
\m{\HasMatSize{99n}{\gL{\cL}^{-1}(y-1)}}.

\subhead{Inputs from Note C}\phantomsection\label{s:inpC}(same unit; \emph{The complex place: the
open-cell chart, the Haar measure of the tower, the Bessel cut and the volume weight}).

\Inp{C1} \emph{Real-place sizes} $\mathfrak{h}$ in $\RealPlaceSizes{\mathfrak{s}}{\rho}{s_0}{b}$ are proofs of
$\rho=b^{6n+2}$, $8r\le\rho$, $k=s_0+3n+2$, $10^{9}n^{3}<s_0$, $n^{7}\le k$, $10^{8}\le\varrho$ and $2\le b$; then
$\rho\ge4$. The level $\cLC{\mathfrak{h}}$ at $\CC$ has $\Omega=\kappa\smallsetminus\{k_0\}$, $\mathfrak{m}=0$,
$\gL{\cLC{\mathfrak{h}}}=1$, the ball $\{v:|v|<\tfrac14\}$, the box $\{x:\|x\|\le\rho\}$ (operator norm), all of
$\Gone{\mathfrak{s}}$ as positions, $\floorMeasure{\cLC{\mathfrak{h}}}{n}$ the Haar probability measure of
$\mathrm{U}(\kappa)=\Hl{n}{T}$, and $\floorMeasure{\cLC{\mathfrak{h}}}{\varrho}$ the convolution of the image of
$\frac{2\pi^{k}}{(k-1)!}\max(|c|^{2}-1,0)^{k-1}\dd c$ under $c\mapsto\boost{\varrho}{c}^{-1}$ with
$\floorMeasure{\cLC{\mathfrak{h}}}{n}$.

\Inp{C2} $0<\haarRatio{\mathfrak{h}}<\infty$; $\groupMeasure{\cLC{\mathfrak{h}}}=\haarRatio{\mathfrak{h}}\cdot\dd y$,
so it is a Haar measure; $\int f\,\dd\upperMeasure{\cLC{\mathfrak{h}}}{0}=\haarRatio{\mathfrak{h}}\int f\,\dd\dot y$ for
measurable $\UT{T}$-invariant $f\ge0$, where $\dd\dot y$ is the measure of $\UT{T}\backslash\Gone{\mathfrak{s}}$ of the
background library; and $\frac{\volU{k}}{\volU{r}}\haarRatio{\mathfrak{h}}\le1$.

\Inp{C3} For real $R\ge\rho$, $\gauge{\mathfrak{h}}{\rho\le R}$ is a gauge in $\Gauges{\cLC{\mathfrak{h}}}$ with sets
$\{x\in\CC:|x|\le R^{i}\}$ and constant $\Lambda=R^{2}$.

\Inp{C4} If the Bessel size of $y\in\Gone{\mathfrak{s}}$ is at most $t$, then $\|uy\|\le t$ for some $u\in\UT{T}$.

\subhead{From the unit on the second moment (unit 148) and the background library.} We recall definitions; none of
their properties is used beyond what is said. For $y\in\Gone{\mathfrak{s}}$ the \emph{Bessel size} is
$\inf_s(\|sy\|+1)$ over $s\in\Ql{n}{T}$. A \emph{cut pair} with respect to $b$, $\mathsf{n}$, $\mathrm{rep}$ is a pair
of continuous functions $\Phi_0$ (full) and $\Phi$ (cut) such that: both are measurable, bounded, satisfy
$F(uy)=\psi_{\CC}(\sd(u))F(y)$ for $u\in\UT{T}$ and $F(hy)=\mathrm{rep}(h)F(y)$ for $h\in\mathrm{U}(\kappa)=\Hl{n}{T}$;
$\Phi_0$ is integrable for $\dd\dot y$; $\int|\Phi_0|^2\dd\dot y=|I|$ and $\int|\Phi|^2\dd\dot y\le|I|$; $\Phi(y)\ne0$
implies that the Bessel size of $y$ is at most $b^{6n+2}$; $\int|\Phi-\Phi_0|^2\dd\dot y\le b^{-10r^2}|I|$; and a
majorant exists, which is not used here and involves~$\mathsf{n}$. For $\mu$ a measure on a group and $W$ a
$\CC^I$-valued function, $\corr{\mu}{W}{\gamma}\in\CC^{I\times I}$ has the coordinates
$\int W_i(y)\overline{W_{i'}(\gamma^{-1}y)}\,\dd\mu(y)$. The constants are
$\cellc{\mathfrak{s}}=\frac{\volU{k}}{\volU{r}}\cdot\frac{2\pi^{k}}{(k-1)!}$ and
$\bandc{\mathfrak{s}}=\frac{\volU{k}}{\volU{r}}\,\pi^{(2\varrho+1)k}$. The \emph{boost} $\boost{\varrho}{c}\in\Hl{\varrho}{T}$,
for $c\in\CC^{\times}$, acts in the plane of $x_{k_0}$ and $z$ and has $\ZT{T}{\boost{\varrho}{c}}_{\varrho\varrho}=c$
when $|c|>1$. For a real $t$ and natural numbers $K$, $i$: $\damp{K}{t}=t^{K}\min(1,t)^{6}$ and
$\hw{K}{i}{p}=\prod_{l<i}\damp{K}{|p_0\cdots p_l|^{-1}}$. A function $\Phi_0$ \emph{has Schur's bound with constant
$D$} if $\int|\int\varphi(y)\Phi_0(\gamma y)\,\dd y|^{2}\dd\dot\gamma\le D\int|\varphi|^2\dd y$ for every integrable
$\varphi\colon\Gone{\mathfrak{s}}\to\CC$. Finally, $\Phi$ \emph{has the band limits with respect to $b$, $\deg$, $J$}
if, with $\beta(\xi)=\bandc{\mathfrak{s}}\,b^{2000n^{9}}\,\damp{k-n^{7}}{|\xi|^{-1}}$ and $u(c)$ as in Theorem~B:
\begin{itemize}
\item[(S1)] for $i\le\varrho$, all $p,q$ and $\xi\in\CC$, $f(c)=\psi_{\CC}(\xi c)\,J_{i+1}(p[i\mapsto u(c)],q[i\mapsto u(c)])$
  is integrable on $\CC$; $\int f=J_i(p,q)$ if $\xi=0$; and
  $|\int f|\le\beta(\xi)\,\hw{n^6}{i}{p}\,\hw{n^6}{i}{q}$ if $\xi\ne0$;
\item[(S2)] for finitely many sequences with a common entry at $\varrho$, the matrix of the values of $J_{\varrho+1}$
  is positive semidefinite; $J_{\varrho+1}$ is continuous;
\item[(S3)] $|J_0(p,q)|\le(1+b^{-10})/\deg$;
\item[(S4)] for $i<\varrho$ and $\xi=p_i^{-1}-q_i^{-1}\ne0$:
  $|J_{i+1}(p,q)|\le\beta(\xi)\,\hw{n^6}{i+1}{p}\,\hw{n^6}{i+1}{q}$;
\item[(S5)] for $\sigma\in\Hl{\varrho}{T}$ and $p,q$ with $\ZT{T}{\sigma}_{\varrho\varrho}=p_\varrho=q_\varrho$:
  \[J_{n}(p,q)=N(p)N(q)\ip{\corr{\dd y}{\Phi}{\wl{\varrho}{p}\sigma}}{\corr{\dd y}{\Phi}{\wl{\varrho}{q}\sigma}},\]
  with $N(p)$ the scalar factor of Definition~\ref{d:sieve}.
\end{itemize}
Three classical results enter through these inputs and directly: the \emph{Cauchy--Schwarz inequality} in the form
$(\int h)^2\le\int f_1\int f_2$ for $h^2\le f_1f_2$ (used freely); \emph{Plancherel's inequality at a place} (inside
\Inp{T12}), which the second reader of the artifact relates to J.~Tate, \emph{Fourier analysis in number fields and
Hecke's zeta-functions} (1950), \S2.2, as a weaker form; and the \emph{Cartan decomposition} of
$\Hl{\varrho}{T}$ at $\CC$: if $\sigma\in\Hl{\varrho}{T}$ has $\ZT{T}{\sigma}_{\varrho\varrho}=c$ with $|c|>1$, then
$\sigma=h\,\boost{\varrho}{c}\,h'$ with $h,h'\in\Hl{n}{T}$, which the second reader relates to the $KAK$ theorem
(A.~Knapp, \emph{Structure theory of semisimple Lie groups}, Theorem~4.10) as a special case with the phase fixed. The
places of these references were given by the second reader and not opened for this note.

\sect{The kernels of a cut function}{s:kernels}
\begin{lemma}[the kernels are pairing kernels]\label{l:pk}
In Setting~\ref{set:cut}, let $W'\colon\GT{T}\to\CC^I$ satisfy \m{\IsBessel{W'}{\gL{\cL}^{-1}}}, let $j$ be a natural
number and $p,q\colon\NN\to A^{\times}$. Then, with $d=\weylRatio{j}{p}{q}$,
\[\IsPairingKernel{\cL}{j}{\translatePairing{W}{W'}{d}{\cdot}{\cdot}}{d}.\]
In particular, for $W'=W$ and $p=q$ constant equal to $1$: \m{\IsPairingKernel{\cL}{j}{\translatePairing{W}{W}{1}{\cdot}{\cdot}}{1}}.
\end{lemma}
\begin{proof}
Write $k(y,y')=\ip{W(y)}{W'(dy')}$; it is measurable. \emph{The law.} Let $\nu,\nu'\in\Nb{l}{T}$ and
$\sigma\in\Hl{l+1}{T}$, which is where $\upperMeasure{\cL}{l+1}$ is carried. Then $\sigma\nu\sigma^{-1}\in\Nb{l}{T}
\subseteq\UT{T}$ has the same entries at the places $(f_a,f_{a+1})$ as $\nu$ (a property of the background library),
so $W(\sigma\nu g)=\psi_A(\gL{\cL}^{-1}\sd(\nu))\,W(\sigma g)$. The matrix of $d$ is diagonal, so
$d(\sigma\nu'\sigma^{-1})d^{-1}\in\UT{T}$, its entries at the places $(f_a,f_{a+1})$ are those of $d\nu'd^{-1}$, and
$W'(d\sigma\nu'y)=\psi_A(\gL{\cL}^{-1}\sd(d\nu'd^{-1}))\,W'(d\sigma y)$. As $\ip{\cdot}{\cdot}$ is conjugate linear in
the first variable and $\overline{\psi_A(c)}=\psi_A(-c)$, the integrand of $\upperPairing{\cL}{l+1}{k}{\nu g}{\nu'y}$
is $\psi_A(\gL{\cL}^{-1}(\sd(d\nu'd^{-1})-\sd(\nu)))$ times that of $\upperPairing{\cL}{l+1}{k}{g}{y}$. \emph{The
element $d$ commutes with $\Hl{j}{T}$}, since $\wl{j}{p}$ and $\wl{j}{q}$ do. \emph{Boundedness.} Let $C,C'$ bound
$|W|,|W'|$. The integrand of $\int|k(\sigma g,\sigma y)|\dd\upperMeasure{\cL}{l}(\sigma)$ is at most $CC'$ and
vanishes unless $W(\sigma g)\ne0$, hence unless $\sigma g$ is near at $\varrho+2$. By \Inp{T9}(a),(b) the measure of
these $\sigma$ is at most $V^{n-l}$ with $V=\Lambda(\Gamma)^{1287n}\vol(\Gamma_{11})^{|\Omega|}
\le\Lambda(\Gamma)^{1287n+12|\Omega|}<\infty$. \emph{Support.} If $g$ is not near at $l$, then for
$\sigma\in\Hl{l}{T}$ the point $\sigma g$ is not near at $l$ either (\Inp{T8}), so $W(\sigma g)=0$ and
$\upperPairing{\cL}{l}{k}{g}{y}=0$. The last assertion is the case in which $\wl{j}{p}=\wl{j}{q}$.
\end{proof}
\leanline{\lean{isPairingKerneltranslatePairing}\sep\lean{isPairingKernelself}}

\begin{lemma}[the integrand at weight $1$]\label{l:one}
In Setting~\ref{set:cut}, let $j$ be a natural number, $p,q\colon\NN\to A^{\times}$ and $x\in\GT{T}$. Then, as
functions of $g$,
\[\formIntegrandOne{j}{W}{x}{\weylRatio{j}{p}{q}}{\wl{j}{p}}{\cdot}=\formIntegrand{j}{p}{q}{W}{x}{1}{\cdot}.\]
\end{lemma}
\begin{proof}
Write $k_1$, $k_d$ for the two kernels of Definition~\ref{d:integrand}; both are pairing kernels at $j$
(Lemma~\ref{l:pk}). By \Inp{T7}(a), $\upperPairing{\cL}{j}{k_1}{g}{gx}$ is the integral over $\sigma$ of
$\upperPairing{\cL}{j+1}{k_1}{\sigma g}{\sigma gx}$, and $\upperPairing{\cL}{j}{k_d}{\wl{j}{p}^{-1}g}{\wl{j}{p}^{-1}gx}$
the integral over $\tau$ of the same expression with $k_d$ and $\tau\wl{j}{p}^{-1}g$. Since
$(\sigma\wl{j}{p}\tau^{-1})^{-1}\sigma g=\tau\wl{j}{p}^{-1}g$, the product of the first integral with the conjugate of
the second is the double integral that defines the right-hand side with $\omega=1$.
\end{proof}
\leanline{\lean{formIntegrandOneeq}}

\begin{lemma}[form, modulus, majorant]\label{l:fmm}
\textup{(i)} Let $\cL$ be a level in $\Levels{\Omega}{T}{\mathfrak{m}}$, $j$ a natural number,
$p,q\colon\NN\to A^{\times}$, $W\colon\GT{T}\to\CC^I$ and $\omega\colon\GT{T}\to\CC$. Then
\[\bigl|\floorForm{j}{p}{q}{W}{\omega}\bigr|\le\headDensity{j}{p}{q}{\cL}\int\formModulus{j}{p}{q}{W}{x}{\omega}\,\dd\groupMeasure{\cL}(x).\]
\textup{(ii)} If moreover $W$ is cut for $\cL$, \m{\IsCut{W}}, and $|\omega(y)|\le1$ for all $y$, then for every
$x\in\GT{T}$
\[\formModulus{j}{p}{q}{W}{x}{\omega}\le\formMajorant{j}{p}{q}{W}{W}{x}.\]
\end{lemma}
\begin{proof}
(i) The absolute value of an integral is at most the integral of the absolute value, and the iterated integral
bounds the integral over the product. (ii) Inside the integral over $g$, the same inequality bounds
$|\formIntegrand{j}{p}{q}{W}{x}{\omega}{g}|$ by the double integral of the product of
$|\upperPairing{\cL}{j+1}{k_1}{\sigma g}{\sigma gx}|$ and
$|\upperPairing{\cL}{j+1}{k_d}{\tau\wl{j}{p}^{-1}g}{\tau\wl{j}{p}^{-1}gx}|$, as in the proof of Lemma~\ref{l:one};
it is the product of the two pivots in the majorant. The measurability needed to split the double integral comes
from that of~$W$.
\end{proof}
\leanline{\lean{enormfloorFormle}\sep\lean{formModulusleformMajorant}}

\begin{lemma}[a pivot under a volume cap]\label{l:cap}
Let $\cL$ be a level in $\Levels{\Omega}{T}{\mathfrak{m}}$, $\Gamma$ a gauge in $\Gauges{\cL}$, $j$ a natural number,
and \m{\IsPairingKernel{\cL}{j}{k}{d}}. Let $x\in\Posn{\cL}$, let $l$ be a natural number, $g\in\GT{T}$ and
$\varepsilon\in[0,\infty]$. Suppose that for every $\sigma\in\Hl{l}{T}$ with $k(\sigma g,\sigma gx)\ne0$ there is $a$
with $l\le a<\varrho$ and $\vol(\floorValues{\cL}{a}{\sigma g}{x})\le\varepsilon$. Then
\[\pivotModulus{\cL}{l}{k}{g}{gx}^{2}\ \le\ \varepsilon^{\volumeExponent{\cL}}\,\bigl(\Lambda(\Gamma)^{10^{4}n}\bigr)^{\varrho-l}\;
\frac{\weightedSquareSum{\cL}{k}{x}{l}{g}}{\floorWeight{\cL}{l}{g}{x}}.\]
\end{lemma}
\begin{proof}
Induction on $\varrho-l$, for all $g$. If $l\ge\varrho$ the pivot is $0$: otherwise \Inp{T7}(g) gives $\sigma\in
\Hl{l}{T}$ with $k(\sigma g,\sigma gx)\ne0$, and the hypothesis an $a$ with $l\le a<\varrho$. Let $l<\varrho$, and write
$s=\volumeExponent{\cL}$, $\Lambda=\Lambda(\Gamma)$, $\mathsf{A}=\upperPairing{\cL}{l+1}{k}{\cdot}{\cdot}$. We apply
Cauchy--Schwarz on floor $l$, \Inp{T7}(e), with $u(g')=\floorWeight{\cL}{l}{g'}{x}/\floorWeight{\cL}{l+1}{g'}{x}$,
whose mass is at most $\Lambda^{10^{4}n}$ by \Inp{T9}(c), and with
\[v(\sigma)=\varepsilon^{s}\bigl(\Lambda^{10^{4}n}\bigr)^{\varrho-l-1}\;
\frac{\weightedSquare{\cL}{l}{k}{x}{\sigma g}+\weightedSquareSum{\cL}{k}{x}{l+1}{\sigma g}}{\floorWeight{\cL}{l}{g}{x}}.\]
Then $\Lambda^{10^4n}\int v\,\dd\floorMeasure{\cL}{l}$ is the right-hand side of the lemma, by the definition of the
weighted square sum. It remains to see $|\mathsf{A}(\sigma g,\sigma gx)|^2\le u(\sigma g)v(\sigma)$ for
$\sigma\in\Hl{l}{T}$. By \Inp{T8},
\[\floorWeight{\cL}{l}{\sigma g}{x}=\floorWeight{\cL}{l}{g}{x},\qquad
\floorWeight{\cL}{l+1}{\sigma g}{x}=\floorWeight{\cL}{l}{g}{x}\vol(\floorValues{\cL}{l}{\sigma g}{x})^{s}.\] If $\vol(\floorValues{\cL}{l}{\sigma g}{x})\le\varepsilon$, then $u(\sigma g)\ge\varepsilon^{-s}$, and
$|\mathsf{A}(\sigma g,\sigma gx)|^2$ is $\weightedSquare{\cL}{l}{k}{x}{\sigma g}$ divided by
$\floorWeight{\cL}{l}{g}{x}$, by definition; the inequality follows since $\Lambda\ge1$. Otherwise the hypothesis of the lemma holds for
$(l+1,\sigma g)$: for $\tau\in\Hl{l+1}{T}$ with $k(\tau\sigma g,\tau\sigma gx)\ne0$ the hypothesis at $(l,g)$ gives
some $a\ge l$, and $a=l$ is excluded because the sets $\floorValues{\cL}{l}{\tau\sigma g}{x}$ and
$\floorValues{\cL}{l}{\sigma g}{x}$ are equal (\Inp{T8}). By \Inp{T7}(a), by induction, and since
$\floorWeight{\cL}{l+1}{\sigma g}{x}^{-1}=u(\sigma g)/\floorWeight{\cL}{l}{g}{x}$,
\[|\mathsf{A}(\sigma g,\sigma gx)|^{2}\le\pivotModulus{\cL}{l+1}{k}{\sigma g}{\sigma gx}^{2}\le
\varepsilon^{s}\Lambda^{10^4n(\varrho-l-1)}\frac{\weightedSquareSum{\cL}{k}{x}{l+1}{\sigma g}}{\floorWeight{\cL}{l+1}{\sigma g}{x}}
\le u(\sigma g)v(\sigma).\qedhere\]
\end{proof}
\leanline{\lean{pivotModulussqleofcapped}}

\sect{The local theorem}{s:local}
\begin{lemma}[from a pointwise bound to the self-norm product]\label{l:snp}
In Setting~\ref{set:local}, let $W'\colon\GT{T}\to\CC^I$ satisfy \m{\IsBessel{W'}{\gL{\cL}^{-1}}}, let
$f\colon\GT{T}\to[0,\infty]$ and $K\in[0,\infty]$, and write $k_1=\translatePairing{W}{W}{1}{\cdot}{\cdot}$,
$k_d=\translatePairing{W}{W'}{d}{\cdot}{\cdot}$ with $d=\weylRatio{j}{p}{q}$. Suppose that for every $g$ with
$f(g)\ne0$
\[f(g)^{2}\le K\,\globalConst{\Gamma}\;\frac{\weightedSquareSum{\cL}{k_1}{x}{j}{g}}{\floorWeight{\cL}{j}{g}{x}}\cdot
\frac{\weightedSquareSum{\cL}{k_d}{x}{j}{\wl{j}{p}^{-1}g}}{\floorWeight{\cL}{j}{\wl{j}{p}^{-1}g}{x}}.\]
Then $\bigl(\int f\,\dd\stackMeasure{\cL}_{0,j}\bigr)^{2}\le K\cdot\selfNormProduct{j}{p}{q}{W}{W'}{x}{\Gamma}$.
\end{lemma}
\begin{proof}
Write $B=\weightedSquareSum{\cL}{k_1}{x}{j}{\cdot}$, $B'=\weightedSquareSum{\cL}{k_d}{x}{j}{\cdot}$,
$\mathsf{v}=\invWeight{\cL}{j}{x}$ and $w=\wl{j}{p}$. Both kernels are pairing kernels at $j$ (Lemma~\ref{l:pk}), so
$B,B'$ are $\Ql{j}{T}$-invariant and vanish off the points near at $j$ (\Inp{T7}(c)). If $f(g)\ne0$ the hypothesis
forces $B(g)\ne0\ne B'(w^{-1}g)$, so $g$ and $w^{-1}g$ are near at $j$ and the two weights are
$\mathsf{v}(g)^{-1}$, $\mathsf{v}(w^{-1}g)^{-1}$. Pairing each sum with the weight of the other point,
\[f(g)^2\le K\,\globalConst{\Gamma}\cdot\bigl[B(g)\,\mathsf{v}(w^{-1}g)\bigr]\cdot\bigl[B'(w^{-1}g)\,\mathsf{v}(g)\bigr]\qquad\text{for all }g.\]
By Cauchy--Schwarz, $(\int f\,\dd\stackMeasure{\cL}_{0,j})^2$ is at most $K\globalConst{\Gamma}$ times the product of
the integrals of the two brackets. The second bracket is $\Phi(w^{-1}g)$ for $\Phi(g)=B'(g)\mathsf{v}(wg)$, which is
$\Hl{j}{T}$-invariant because $w$ commutes with $\Hl{j}{T}$; by \Inp{T2}, \Inp{T1} and \Inp{T11} (with $b=j$)
\[\int B'(w^{-1}g)\,\mathsf{v}(g)\,\dd\stackMeasure{\cL}_{0,j}(g)=\int B'(g)\,\mathsf{v}(wg)\,\dd\stackMeasure{\cL}_{0,j}(g)
\le\prod_{a<j}\floorModulus{\cL}{a}{p_a^{-1}}\cdot\globalConst{\Gamma}\cdot\selfNorm{\cL}{j}{k_d}{x}.\]
For the first bracket, $w^{-1}=\wl{j}{\bar p}$ and $\bar p$ is a head of length~$j$, so \Inp{T11} applies with
$\bar p$ and gives the bound $\prod_{a<j}\floorModulus{\cL}{a}{\bar p_a^{-1}}\cdot\globalConst{\Gamma}\cdot
\selfNorm{\cL}{j}{k_1}{x}$; here $\floorModulus{\cL}{a}{\bar p_a^{-1}}=\floorModulus{\cL}{a}{p_a^{-1}}$, and
$\selfNorm{\cL}{j}{k_1}{x}\le\selfNorm{\cL}{0}{k_1}{x}$ by \Inp{T7}(f). The product of the three bounds is $K$ times
the self-norm product.
\end{proof}
\leanline{\lean{sqlintegralleselfNormProduct}}

\begin{proof}[Proof of Theorem A\textup{(a)}]
Apply Lemma~\ref{l:snp} to the function
\[f(y)=\pivotModulus{\cL}{j}{k_1}{y}{yx}\cdot\pivotModulus{\cL}{j}{k_d}{w^{-1}y}{w^{-1}yx},\qquad w=\wl{j}{p},\]
whose integral is the majorant, and to $K=\sigma^{s}V^{2}$, $s=\volumeExponent{\cL}$. Let
$f(y)\ne0$. For the second pivot, Cauchy--Schwarz on floor $j$ (\Inp{T7}(e) with $u=1$, $K=V$) and \Inp{T7}(h) give,
at $g'=w^{-1}y$,
\[\pivotModulus{\cL}{j}{k_d}{g'}{g'x}^2\le V\,\weightedSquareSum{\cL}{k_d}{x}{j}{g'}/\floorWeight{\cL}{j}{g'}{x}.\] If $\sigma\ge1$ the same bound holds for the first pivot, and $\sigma^{s}\ge1$,
$\globalConst{\Gamma}\ge1$, $V\le V^{2}$. If $\sigma<1$ we use Lemma~\ref{l:cap} with $l=j$ and $\varepsilon=\sigma$;
its hypothesis is checked as follows. Since the second pivot is not $0$ there is $\tau\in\Hl{j}{T}$ with
$k_d(\tau w^{-1}y,\tau w^{-1}yx)\ne0$ (\Inp{T7}(g)); thus $W(\tau w^{-1}y)\ne0$ and $W'(d\tau w^{-1}yx)\ne0$, where
$d\tau w^{-1}=\tau\wl{j}{q}^{-1}$. So $\tau w^{-1}y$ and $\tau\wl{j}{q}^{-1}yx$ are near at $j$. For
$\rho_1\in\Hl{j}{T}$ with $k_1(\rho_1y,\rho_1yx)\ne0$ we get $\ip{W(\rho_1y)}{W(\rho_1yx)}\ne0$, and
$w^{-1}\rho_1y=(\rho_1\tau^{-1})\tau w^{-1}y$ and $\wl{j}{q}^{-1}\rho_1yx$ are near at $j$ (\Inp{T8}); the
hypothesis of the theorem at the point $\rho_1y$ gives the $a$ that Lemma~\ref{l:cap} asks for. Its conclusion
has the factor $(\Lambda(\Gamma)^{10^4n})^{\varrho-j}\le\globalConst{\Gamma}$.
\end{proof}

\begin{proof}[Proof of Theorem A\textup{(b)}]
Apply Lemma~\ref{l:snp} with $W'=W$, $f(g)=|\formIntegrand{j}{p}{q}{W}{x}{\omega_\eta}{g}|$ and
$K=(c_A/\modA(\xi))^{|\Omega|}$. Fix $g$ with $f(g)\ne0$ and put $w=\wl{j}{p}$,
$F_1(\sigma)=\upperPairing{\cL}{j+1}{k_1}{\sigma g}{\sigma gx}$,
$F_2(\tau)=\upperPairing{\cL}{j+1}{k_d}{\tau w^{-1}g}{\tau w^{-1}gx}$. By \Inp{T7}(h) it suffices to show
\begin{equation}\label{e:char}
f(g)^{2}\le K\,\globalConst{\Gamma}\int|F_1|^{2}\,\dd\floorMeasure{\cL}{j}\cdot\int|F_2|^{2}\,\dd\floorMeasure{\cL}{j}.
\end{equation}
By definition and since $\ZT{T}{\sigma\tau^{-1}}_{jj}=h_T(\sigma^{-1}f_j,\tau^{-1}f_j)$,
\[\formIntegrand{j}{p}{q}{W}{x}{\omega_\eta}{g}=\iint\psi_A\bigl(\eta\,h_T(\sigma^{-1}f_j,\tau^{-1}f_j)\bigr)F_1(\sigma)\overline{F_2(\tau)}\,
\dd\floorMeasure{\cL}{j}(\tau)\,\dd\floorMeasure{\cL}{j}(\sigma).\]
As this is not $0$, $F_1$ and $F_2$ are not almost everywhere $0$ on $\Hl{j}{T}$, so $g$ and $w^{-1}g$ are near
at~$j$ (the last clause of the definition of a pairing kernel, and \Inp{T8}). By \Inp{T14}(b) there is
$t\in\Hl{j}{T}$ with \m{\HasMatSize{60n^2}{tg}} and \m{\HasMatSize{99n^2}{tw^{-1}g}} for $\Gamma$. By \Inp{T12}(i),
applied to both kernels with this $t$, for almost all $\sigma,\tau$ with $F_1(\sigma)F_2(\tau)\ne0$ the vectors
$t\sigma^{-1}f_j$ and $t\tau^{-1}f_j$ have sizes at most $60n^2+6$ and $99n^2+6$; hence
$h_T(\sigma^{-1}f_j,\tau^{-1}f_j)=h_T(t\sigma^{-1}f_j,t\tau^{-1}f_j)$ lies in $\Gamma_{200n^2}$, by the rules for
sizes of products and sums, and there $\psi_A(\eta\,\cdot)=\psi_A(\xi\,\cdot)$. So, with
$G(v)=\int\psi_A(-\xi\,h_T(v,t\tau^{-1}f_j))F_2(\tau)\,\dd\floorMeasure{\cL}{j}(\tau)$,
\[\formIntegrand{j}{p}{q}{W}{x}{\omega_\eta}{g}=\int_{\{\sigma:\,t\sigma^{-1}f_j\in\vecsOfSize{60n^2+6}\}}F_1(\sigma)\,\overline{G(t\sigma^{-1}f_j)}\;\dd\floorMeasure{\cL}{j}(\sigma).\]
By Cauchy--Schwarz, $f(g)^2\le\int|F_1|^2\dd\floorMeasure{\cL}{j}$ times the integral of
$|G(t\sigma^{-1}f_j)|^2$ over the same set, which by \Inp{T6} is $\int_{\vecsOfSize{60n^2+6}}|G|^2\dd\quadric{\cL}{j}$.
By \Inp{T12}(ii) for the kernel $k_d$, the points $w^{-1}g$, $w^{-1}gx$, the unit $-\xi$, $y'=0$ and the quadric
measure, this is at most $\Lambda(\Gamma)^{99n(60n^2+12)}\Lambda(\Gamma)^{99n(99n^2+12)}K\int|F_2|^2\dd\floorMeasure{\cL}{j}$,
and $99n(159n^2+24)\le10^6n^3$. This is \eqref{e:char}.
\end{proof}

\begin{proof}[Proof of Theorem A\textup{(c)}]
Here $j=i+1$. Write $w=\wl{j}{p}$, $\mathsf{A}_1=\upperPairing{\cL}{j}{k_1}{\cdot}{\cdot}$,
$\mathsf{A}_d=\upperPairing{\cL}{j}{k_d}{\cdot}{\cdot}$ and $K=(c_A/\modA(\xi))^{|\Omega|}$. By
Lemma~\ref{l:one} the modulus is $\int h\,\dd\stackMeasure{\cL}_{0,j}$ with
$h(g)=|\mathsf{A}_1(g,gx)|\cdot|\mathsf{A}_d(w^{-1}g,w^{-1}gx)|$. Define, for $g\in\GT{T}$,
\[\begin{gathered}\mathsf{r}(g)=\floorWeight{\cL}{j}{g}{x}\cdot\invWeight{\cL}{i}{x}(\wl{i}{p}g),\\
\mathsf{u}(g)=\mathsf{r}(g)\,|\mathsf{A}_d(g,gx)|^{2}\ \text{ if $wg$ is near at $j$},\qquad\mathsf{u}(g)=0\ \text{ otherwise.}\end{gathered}\]
\emph{Step 1.} For all $g$: $h(g)^{2}\le\mathsf{u}(w^{-1}g)\cdot\weightedSquare{\cL}{i}{k_1}{x}{g}\cdot\invWeight{\cL}{j}{x}(w^{-1}g)$.
Indeed, if $h(g)\ne0$ then $g$ and $w^{-1}g$ are near at $j$; $\wl{i}{p}w^{-1}=\wo{i}{p_i}^{-1}\in\Hl{i}{T}$, so
$\invWeight{\cL}{i}{x}(\wl{i}{p}w^{-1}g)=\floorWeight{\cL}{i}{g}{x}^{-1}$ (\Inp{T8}); and
$\weightedSquare{\cL}{i}{k_1}{x}{g}=\floorWeight{\cL}{i}{g}{x}|\mathsf{A}_1(g,gx)|^2$. The weights cancel and
equality holds.

\emph{Step 2.} By Cauchy--Schwarz,
\[\Bigl(\int h\Bigr)^2\le\int\mathsf{u}(w^{-1}g)\,\dd\stackMeasure{\cL}_{0,j}(g)\cdot
\int\weightedSquare{\cL}{i}{k_1}{x}{g}\,\invWeight{\cL}{j}{x}(\wl{j}{\bar p}g)\,\dd\stackMeasure{\cL}_{0,j}(g).\]
The
weighted square of floor $i$ is $\Ql{j}{T}$-invariant and zero off the points near at $j$ (\Inp{T7}(c) for the kernel
at $i$); so by \Inp{T1}, \Inp{T11} (head $\bar p$, $b=j$) and \Inp{T7}(f) the second factor is at most
\[\begin{aligned}\prod_{a<j}\floorModulus{\cL}{a}{p_a^{-1}}\cdot\globalConst{\Gamma}\cdot{}&
\int\weightedSquare{\cL}{i}{k_1}{x}{g}\,\dd(\floorMeasure{\cL}{i}*\dots*\floorMeasure{\cL}{0})(g)\\ \le{}&
\prod_{a<j}\floorModulus{\cL}{a}{p_a^{-1}}\cdot\globalConst{\Gamma}\cdot\selfNorm{\cL}{0}{k_1}{x}.\end{aligned}\]

\emph{Step 3.} The function $\mathsf{u}$ is $\Hl{j}{T}$-invariant (\Inp{T7}(b), \Inp{T8}), so by \Inp{T2} and
\Inp{T1}, with $\rho_0=\floorMeasure{\cL}{i}*\dots*\floorMeasure{\cL}{0}$ and $\etab{j}=\eta_i*\etab{i}$,
\[\int\mathsf{u}(w^{-1}g)\,\dd\stackMeasure{\cL}_{0,j}(g)=\int\!\!\int\mathsf{u}(\nu g)\,\dd\eta_i(\nu)\,\dd(\etab{i}*\rho_0)(g).\]
\emph{Step 4: Plancherel on $\Nf{i}{T}$.} Fix $g$. We apply \Inp{T13} with $a=i$, the kernel $k_d$ (a pairing
kernel at $i+1$), $x_1=g$, $y_1=gx$, $\varepsilon=1$, $C=\mathsf{r}(g)$, $\Phi(\nu)=\mathsf{u}(\nu g)$, the unit
$-\xi p_i$ in the place of $\xi$ and $u=p_i^{-1}$, so that $\modA(\xi u)$ there is $\modA(\xi)$. Its three hypotheses
hold. First, an element $s\in\Nf{i}{T}$ has only one non-zero entry at the places $(f_a,f_{a+1})$, the one with
$a=i$; the diagonal matrix $d$ has the entry $p_i/q_i$ at $e_i$ and $1$ at $f_{i+1}$; hence
$\gL{\cL}^{-1}(\sd(dsd^{-1})-\sd(s))=\gL{\cL}^{-1}(p_i/q_i-1)\sd(s)=-\xi p_i\,\sd(s)$. Second, $\mathsf{r}(\nu g)=
\mathsf{r}(g)$ for $\nu\in\Nf{i}{T}$, because $\nu\in\Ql{j}{T}$ and $\nu$ commutes with $\wl{i}{p}$ and lies in
$\Ql{i}{T}$ (\Inp{T8}); so $\mathsf{u}(\nu g)\le\mathsf{r}(g)|\mathsf{A}_d(\nu g,\nu gx)|^2$. Let $S$ be the set of the
$\nu\in\Nf{i}{T}$ with $\mathsf{u}(\nu g)\ne0$. For $\nu\in S$ the points $w\nu g$ and $\nu g$ are near at $j$, so by
\Inp{T14}(b) there is $\tau_\nu\in\Hl{j}{T}$ with \m{\HasMatSize{99n^2}{\tau_\nu\nu g}}, and by \Inp{T14}(c)
$p_i\in\Gamma_{10n}$. Third, for $m\in S$ take $t=\tau_m$ and $q=m$; for $\nu\in S$ take $\nu'=1$: since
$\tau_\nu$ fixes $e^*_i$, the column at $e^*_i$ of $tm(t\nu)^{-1}-1$ is that of
$(\tau_mmg)(\tau_\nu\nu g)^{-1}-1$, of size at most $198n^2+6$, and multiplied by $u^{-1}=p_i$ it has size at
most $198n^2+10n+6$. The conclusion of \Inp{T13}, with \Inp{T7}(h) at the point $g$, reads
\[\begin{aligned}\int\mathsf{u}(\nu g)\,\dd\eta_i(\nu)&\le\mathsf{r}(g)\,K\,\floorModulus{\cL}{i}{p_i^{-1}}\,\Lambda(\Gamma)^{N}
\frac{\weightedSquareSum{\cL}{k_d}{x}{j}{g}}{\floorWeight{\cL}{j}{g}{x}}\\
&=K\,\floorModulus{\cL}{i}{p_i^{-1}}\,\Lambda(\Gamma)^{N}\,\weightedSquareSum{\cL}{k_d}{x}{j}{g}\,\invWeight{\cL}{i}{x}(\wl{i}{p}g),\end{aligned}\]
with $N=100n(297n^2+10n+37)\le10^6n^3$.

\emph{Step 5.} Integrate over $g$ against $\etab{i}*\rho_0$ and use \Inp{T11} with the head $p$ of length $i$, $b=j$
and $\Phi=\weightedSquareSum{\cL}{k_d}{x}{j}{\cdot}$:
\[\int\mathsf{u}(w^{-1}g)\,\dd\stackMeasure{\cL}_{0,j}(g)\le K\,\globalConst{\Gamma}\prod_{a<j}\floorModulus{\cL}{a}{p_a^{-1}}\cdot
\globalConst{\Gamma}\cdot\selfNorm{\cL}{j}{k_d}{x}.\]
With Step 2 this is the claim.
\end{proof}

\sect{Descent}{s:descent}
\begin{lemma}[a diagonal entry of the frame Gram matrix]\label{l:diag}
Let $A$ be a commutative ring with involution, $\mathfrak{s}$ a shape, $T$ an $\mathfrak{s}$-form over $A$, $b$ a
natural number with $b\le\varrho$, and $x\in\Ql{b+1}{T}\cap\Hl{b}{T}$. Then for every $y\in\GT{T}$
\[\ZT{T}{xy}_{bb}=\ZT{T}{y}_{bb}\qquad\text{and}\qquad\ZT{T}{yx}_{bb}=\ZT{T}{y}_{bb}.\]
\end{lemma}
\begin{proof}
The group $\Ql{b+1}{T}\cap\Hl{b}{T}$ is the stabiliser of $f_b$ in $\Hl{b}{T}$, so $xf_b=f_b=x^{-1}f_b$. Hence
$h_T(f_b,yxf_b)=h_T(f_b,yf_b)$ and $h_T(f_b,xyf_b)=h_T(x^{-1}f_b,yf_b)=h_T(f_b,yf_b)$.
\end{proof}
\leanline{\lean{frameGramdiagmul}}

\begin{lemma}[descent on one floor]\label{l:succ}
In Setting~\ref{set:cut}, let $p,q\colon\NN\to A^{\times}$, let $b$ be a natural number with $b\le\varrho$, $\zeta$
an $s$-finite measure on $\GT{T}$, $v\colon A\to\GT{T}$ and $\rho\colon A\to[0,\infty]$ measurable maps,
$\Psi\colon A\to\CC$, $t\in\CC$ and $\eta\in A$. Suppose that
$\int\formMajorant{b}{p}{q}{W}{W}{x}\,\dd\groupMeasure{\cL}(x)<\infty$; that $\zeta$ is carried by
$\Ql{b+1}{T}\cap\Hl{b}{T}$; that $\rho(c)<\infty$ for almost every $c$; that $\floorMeasure{\cL}{b}=\varpi*\zeta$,
where $\varpi$ is the image of $\rho(c)\,\dd c$ under $c\mapsto v(c)^{-1}$; and that for almost every $c\in A$
\[\begin{aligned}\Psi(c)={}&t\cdot\headDensity{b}{p}{q}{\cL}\,\rho(c)\cdot\psi_A\bigl(\eta\,\ZT{T}{v(c)}_{bb}\bigr)\\
&\cdot\iint\formIntegrandOne{b+1}{W}{x}{\weylRatio{b}{p}{q}}{\wl{b}{p}v(c)}{g}\;\dd\mu_b(g)\,\dd\groupMeasure{\cL}(x),\end{aligned}\]
with $\mu_b=(\zeta*\floorMeasure{\cL}{b})*\stackMeasure{\cL}_{0,b}$. Then $\Psi$ is integrable and
\[\int_A\Psi(c)\,\dd c=t\cdot\floorForm{b}{p}{q}{W}{\omega_\eta},\qquad\omega_\eta(y)=\psi_A\bigl(\eta\,\ZT{T}{y}_{bb}\bigr).\]
\end{lemma}
\begin{proof}
Write $d=\weylRatio{b}{p}{q}$, $w=\wl{b}{p}$, and $\mathsf{I}(x,w',g)=\formIntegrandOne{b+1}{W}{x}{d}{w'}{g}$, which
by definition is $\upperPairing{\cL}{b+1}{k_1}{g}{gx}\,\overline{\upperPairing{\cL}{b+1}{k_d}{w'^{-1}g}{w'^{-1}gx}}$.
Both kernels are pairing kernels at $b$ (Lemma~\ref{l:pk}). \emph{Claim:} the function
$(x,g,\sigma,\tau)\mapsto\omega_\eta(\sigma\tau^{-1})\mathsf{I}(x,\sigma w\tau^{-1},\sigma g)$ under
$\groupMeasure{\cL}\times\stackMeasure{\cL}_{0,b}\times\floorMeasure{\cL}{b}\times\floorMeasure{\cL}{b}$ and the
function $(x,g,c)\mapsto\omega_\eta(v(c))\mathsf{I}(x,wv(c),g)$ under $\groupMeasure{\cL}\times\mu_b\times\rho(c)\dd c$
have the same distribution. To see it, fix $x$, $g$, $\sigma\in\Hl{b}{T}$ and a test function. In the integral
over $\tau$ replace $\tau$ by $\tau\sigma$, by \Inp{T4} with $S=\Hl{b}{T}$: this is allowed because
$\tau\mapsto\omega_\eta(\tau^{-1})\mathsf{I}(x,w\tau^{-1},\sigma g)$ is invariant under
$\Ql{b+1}{T}\cap\Hl{b}{T}$, by Lemma~\ref{l:diag} and \Inp{T7}(b), and because $\sigma$ commutes with $w$. Now write
$\tau=v(c)^{-1}\nu$ with $c$ distributed by $\rho(c)\dd c$ and $\nu$ by $\zeta$. Then
$\omega_\eta(\tau^{-1})=\omega_\eta(\nu^{-1}v(c))=\omega_\eta(v(c))$ by Lemma~\ref{l:diag}, and
$\mathsf{I}(x,w\nu^{-1}v(c),\sigma g)=\mathsf{I}(x,wv(c),\nu\sigma g)$, since $\nu$ commutes with $w$ and
$\upperPairing{\cL}{b+1}{k_1}{\nu\sigma g}{\nu\sigma gx}=\upperPairing{\cL}{b+1}{k_1}{\sigma g}{\sigma gx}$
(\Inp{T7}(b)). The law of $\nu\sigma g$ is $\mu_b$. This proves the claim.

The first function is integrable: as in Lemma~\ref{l:fmm}(ii), and because $|\omega_\eta|=1$, its absolute value
integrates to at most the integral of the majorant $\formMajorant{b}{p}{q}{W}{W}{x}$ over $x$, which is finite.
So the second is integrable, and Fubini's theorem gives that $\Psi$ is integrable with integral
$t\,\headDensity{b}{p}{q}{\cL}$ times the integral of the first function, which is the floor form.
\end{proof}
\leanline{\lean{integralsuccflooreq}}

\begin{proof}[Proof of Theorem B]
Apply Lemma~\ref{l:succ} with $b=a<\varrho$, $\zeta=\eta_a$, which is carried by $\Nf{a}{T}\subseteq
\Ql{a+1}{T}\cap\Hl{a}{T}$, with $v(c)=\wo{a}{c}$ (for units $c$; almost every $c$ is a unit) and
$\rho(c)=\modA(c)^{k+1+2(\varrho-a-1)}$; then $\floorMeasure{\cL}{a}=\varpi_a*\eta_a$ is the required
decomposition. It remains to identify $f(c)$ with $\Psi(c)$ for a unit $c$. Put $p'=p[a\mapsto c]$,
$q'=q[a\mapsto c]$. By Lemma~\ref{l:one},
\[\begin{aligned}\floorForm{a+1}{p'}{q'}{W}{1}={}&\headDensity{a+1}{p'}{q'}{\cL}\\
&\cdot\iint\formIntegrandOne{a+1}{W}{x}{d'}{\wl{a+1}{p'}}{g}\,\dd\stackMeasure{\cL}_{0,a+1}(g)\,\dd\groupMeasure{\cL}(x),\end{aligned}\]
with $d'=\weylRatio{a+1}{p'}{q'}$.
Here $\wl{a+1}{p'}=\wo{a}{c}\wl{a}{p}=\wl{a}{p}\wo{a}{c}$, so $d'=\weylRatio{a}{p}{q}$;
$\stackMeasure{\cL}_{0,a+1}=(\eta_a*\floorMeasure{\cL}{a})*\stackMeasure{\cL}_{0,a}$;
$\headDensity{a+1}{p'}{q'}{\cL}=\headDensity{a}{p}{q}{\cL}\,\rho(c)$; and $\ZT{T}{\wo{a}{c}}_{aa}=c$.
\end{proof}

\sect{Integration over positions}{s:positions}
\begin{proof}[Proof of Theorem C]
Write $\mathsf{A}_l=\upperPairing{\cL}{l}{k}{\cdot}{\cdot}$, $P=\positionsOfSize{\cL}{\Gamma}$,
$\rho^{(l)}=\floorMeasure{\cL}{l}*\dots*\floorMeasure{\cL}{0}$, and for $l\le n$
\[\begin{gathered}\mathsf{W}_l(g,y)=\etab{l}\{\nu:g^{-1}\nu y\in P\},\\
\Phi_l(g,y)=|\mathsf{A}_{l+1}(g,y)|^{2}\ \text{ if }\mathsf{W}_l(g,y)>0,\qquad\Phi_l(g,y)=0\ \text{ otherwise}.\end{gathered}\]
We prove, by decreasing induction on $l$ from $n$ to $j$, that there is an $s$-finite measure $\nu_l$ with
\begin{equation}\label{e:ind}\begin{gathered}
l\,c+\int_P\selfNorm{\cL}{l}{k}{x}\,\dd\groupMeasure{\cL}(x)\le n\,c,\qquad\groupMeasure{\cL}=\etab{l}*\nu_l,\\
\prod_{a<\min(l,\varrho)}\floorConst{\Gamma}{a}\cdot\iint\Phi_l(g,y)\,\dd\nu_l(y)\,\dd\rho^{(l)}(g)\le c.
\end{gathered}\end{equation}
The first inequality at $l=j$ is the theorem. For $l=n$ take $\nu_n=\upperMeasure{\cL}{0}=\rho^{(n)}$: the self-norm
at $n$ is $0$; $\groupMeasure{\cL}=\etab{n}*\upperMeasure{\cL}{0}$ by \Inp{T1}; and $\Phi_n\le|k|^2$ because
$\upperMeasure{\cL}{n+1}=\delta_1$. Let $j\le l\le\varrho$ and suppose \eqref{e:ind} for $l+1$. Put
$\nu_l=\eta_l*\nu_{l+1}$; then $\groupMeasure{\cL}=\etab{l}*\nu_l$.

\emph{The third part of \eqref{e:ind} at $l$.} As $\rho^{(l+1)}=\floorMeasure{\cL}{l+1}*\rho^{(l)}$, the induction
hypothesis bounds $\prod_{a\le l,\,a<\varrho}\floorConst{\Gamma}{a}$ times the integral of
$\int\Phi_{l+1}(\sigma g,y)\,\dd\nu_{l+1}(y)$ over $g$ and $\sigma\in\Hl{l+1}{T}$. For such $\sigma$, \Inp{T3}
applies to $\groupMeasure{\cL}=\etab{l+1}*\nu_{l+1}$ and $S=\Nb{l+1}{T}$ (by \Inp{T1}), the function
$y\mapsto\Phi_{l+1}(\sigma g,y)$ being $\Nb{l+1}{T}$-invariant by the law of the kernel and the invariance of
$\etab{l+1}$; so the inner integral equals $\int\Phi_{l+1}(\sigma g,\sigma y)\,\dd\nu_{l+1}(y)$. It therefore
suffices to show, for all $g$, $y$,
\begin{equation}\label{e:step}
\int\Phi_l(g,\nu y)\,\dd\eta_l(\nu)\le C\int\Phi_{l+1}(\sigma g,\sigma y)\,\dd\floorMeasure{\cL}{l+1}(\sigma),\qquad
C=\floorConst{\Gamma}{l}\ (l<\varrho),\quad C=1\ (l=\varrho).
\end{equation}
If $\mathsf{W}_{l+1}(g,y)=0$ the left side is $0$, because $\int\mathsf{W}_l(g,\nu y)\dd\eta_l(\nu)=\mathsf{W}_{l+1}(g,y)$.
Otherwise $\mathsf{W}_{l+1}(\sigma g,\sigma y)=\mathsf{W}_{l+1}(g,y)>0$ for $\sigma\in\Hl{l+1}{T}$ (\Inp{T1}), so
the right side is $C\int|\mathsf{A}_{l+2}(\sigma g,\sigma y)|^2\dd\floorMeasure{\cL}{l+1}(\sigma)$. For $l=\varrho$,
$\eta_\varrho=\delta_1$, and \eqref{e:step} follows from $|\mathsf{A}_{n}|\le\pivotModulus{\cL}{n}{k}{\cdot}{\cdot}$
and Cauchy--Schwarz on floor $n$ (\Inp{T7}(a),(e) with $u=1$, $K=1$; the mass bound is the third requirement on
$\floorMeasure{\cL}{n}$ in the definition of a level). For $l<\varrho$ we apply \Inp{T13} with $a=l$, the unit
$\gL{\cL}^{-1}$ in the place of $\xi$, $u=\gL{\cL}$, $\varepsilon=0$, $C=1$, $x_1=g$, $y_1=y$,
$\Phi(\nu)=\Phi_l(g,\nu y)$ and $S$ the set of the $\nu\in\Nf{l}{T}$ with $\mathsf{W}_l(g,\nu y)>0$ and $g$ near at
$l+1$. Its hypotheses: $d$ commutes with $\Nf{l}{T}\subseteq\Hl{j}{T}$ and $\sd(1)=0$; if $\Phi_l(g,\nu y)\ne0$
then $\nu\in S$, by the last clause of the definition of a pairing kernel; and for $m\in S$ choose
$\xi'\in\Ql{l+1}{T}$ with $\xi'g\in\BoxL{\cL}$ and write $\xi'=t\nu_0$ with $t\in\Hl{l}{T}$,
$\nu_0\in\Nb{l}{T}$; then $t\in\Ql{l+1}{T}\cap\Hl{l}{T}$, $q=\nu_0$ and \m{\HasMatSize{1}{tqg}} for $\Gamma$. For
$\nu\in S$ choose $\nu_1,\nu_2\in\Nb{l}{T}$ with $g^{-1}\nu_1my$ and $g^{-1}\nu_2\nu y$ in $P$; by \Inp{T14}(d) the
element $(\xi'g)(g^{-1}\nu_1my)\bigl((\xi'g)(g^{-1}\nu_2\nu y)\bigr)^{-1}=\xi'\nu_1m\nu^{-1}\nu_2^{-1}\xi'^{-1}$ has
$\gL{\cL}^{-1}(\,\cdot-1)$ of size at most $14$, and it equals $\nu'\,tm(t\nu)^{-1}$ with
$\nu'\in\Nb{l}{T}$, because $\Hl{l}{T}$ normalises $\Nb{l}{T}$. So $i_1=1$, $i_2=14$, and the constant of \Inp{T13} is
$c_A^{|\Omega|}\floorModulus{\cL}{l}{\gL{\cL}}\Lambda(\Gamma)^{4600n}=\floorConst{\Gamma}{l}$.

\emph{The first part of \eqref{e:ind} at $l$.} By \Inp{T7}(f) and by the third part at $l$, just proved, it
suffices to show
\[\int_P\int\weightedSquare{\cL}{l}{k}{x}{g}\,\dd\rho^{(l)}(g)\,\dd\groupMeasure{\cL}(x)\le
\prod_{a<l}\floorConst{\Gamma}{a}\iint\Phi_l(g,y)\,\dd\nu_l(y)\,\dd\rho^{(l)}(g).\]
Fix $g$. Since
$\groupMeasure{\cL}$ is invariant under left translation,
\[\begin{aligned}\int_P{}&\floorWeight{\cL}{l}{g}{x}|\mathsf{A}_{l+1}(g,gx)|^2\dd\groupMeasure{\cL}(x)\\
&=\iint\mathbf{1}_P(g^{-1}\nu y)\,\floorWeight{\cL}{l}{g}{g^{-1}\nu y}\,|\mathsf{A}_{l+1}(g,\nu y)|^2\,\dd\etab{l}(\nu)\,\dd\nu_l(y).\end{aligned}\]
By the law of the kernel, $|\mathsf{A}_{l+1}(g,\nu y)|=|\mathsf{A}_{l+1}(g,y)|$ for $\nu\in\Nb{l}{T}$. If this is not
$0$, then $g$ is near at $l+1$, hence at $l$, and \Inp{T10} bounds the integral over $\nu$ by
$\prod_{a<l}\floorConst{\Gamma}{a}$; and the integral over $\nu$ vanishes unless $\mathsf{W}_l(g,y)>0$. So the inner
integral is at most $\prod_{a<l}\floorConst{\Gamma}{a}\cdot\Phi_l(g,y)$.
\end{proof}

\begin{proof}[Proof of Theorem D]
Both kernels $k_1$, $k_d$ ($d=\weylRatio{j}{p}{q}$) are pairing kernels (Lemma~\ref{l:pk}), at $0$ and at $j$. On $S$,
$f\le e\cdot\selfNormProduct{j}{p}{q}{W}{W'}{x}{\Gamma}^{1/2}$. By Cauchy--Schwarz and Theorem~C for the gauge
$\Gamma'$, with $D=\prod_{l<\varrho}\floorConst{\Gamma'}{l}$,
\[\begin{aligned}\Bigl(\int_{\positionsOfSize{\cL}{\Gamma'}}\bigl(\selfNorm{\cL}{0}{k_1}{x}\,\selfNorm{\cL}{j}{k_d}{x}\bigr)^{1/2}&\dd\groupMeasure{\cL}(x)\Bigr)^{2}\\
\le n^2D^{2}\;\pairingNorm{W}{W}\cdot{}&\iint|k_d|^2\,\dd\upperMeasure{\cL}{0}\,\dd\upperMeasure{\cL}{0},\end{aligned}\]
and by \Inp{T5} the last double integral is $\prod_{a<j}\modA(p_a/q_a)^{k+2(\varrho-a)}\cdot\pairingNorm{W}{W'}$.
Collect the powers of the modules: the square of the head density is
$\prod_{a<j}\modA(p_aq_a)^{k-1+2(\varrho-a)}$, the square of $\prod_{a<j}\floorModulus{\cL}{a}{p_a^{-1}}$ is
$\prod_{a<j}\modA(p_a)^{-2(k-1+2(\varrho-a))-1}$, and with the factor from \Inp{T5} the exponent of $\modA(p_a)$ is $0$
and that of $\modA(q_a)$ is $-1$. As $q$ is a head, $(q_0\cdots q_{j-1})^{-1}\in\Gamma_1\subseteq\Gamma_2$ and its
module is at most $\Lambda(\Gamma)^2$. Altogether the square of the left side of the theorem is at most
\[e^2\,\Lambda(\Gamma)^{2}\,\globalConst{\Gamma}^{3}\,n^2D^2\,\Theta^2,\qquad
D^2=\Bigl(\prod_{l<\varrho}\floorModulus{\cL}{l}{\gL{\cL}}\Bigr)^{2}c_A^{2|\Omega|\varrho}\Lambda(\Gamma')^{9200n\varrho},\]
and $n^2\le4^{n}\le\Lambda(\Gamma')^{2n}$, $9200n\varrho+2n\le10^5n^2$. This is the square of the right side.
\end{proof}

\sect{The complex place}{s:complex}
\begin{setting}[cut Bessel data]\label{set:cb}
Let $\mathfrak{s}$ be a shape, $I$ a finite set~$I$, $\Phi,\Phi_0\colon\Gone{\mathfrak{s}}\to\CC^I$ continuous,
$\mathrm{rep}$ a homomorphism from $\mathrm{U}(\kappa)$ to the unitary group of $\CC^I$, $\rho,b$ real numbers,
$s_0,\mathsf{n}$ natural numbers, and $\mathfrak{c}$ cut Bessel data,
\m{\IsCutBessel{\Phi}{\mathrm{rep}}{\rho}{s_0}{b}{\Phi_0}{\mathsf{n}}}, with underlying sizes $\mathfrak{h}$. Floor
forms, moduli, majorants, nearness and the predicate $\operatorname{Cut}$ refer to the level $\cLC{\mathfrak{h}}$.
\end{setting}

\begin{lemma}[the cut is cut]\label{l:iscut}
In Setting~\ref{set:cb}, $\Phi$ is a cut Bessel function for the level $\cLC{\mathfrak{h}}$: \m{\IsCut{\Phi}}.
\end{lemma}
\begin{proof}
$\Phi$ is measurable, bounded and satisfies the law of frequency $1=\gL{\cLC{\mathfrak{h}}}^{-1}$, as a member of a cut
pair. If $\Phi(y)\ne0$ the Bessel size of $y$ is at most $b^{6n+2}=\rho$, so $\|uy\|\le\rho$ for some $u\in\UT{T}$
by \Inp{C4}, and $uy$ lies in the box.
\end{proof}
\leanline{\lean{isCutcomplexLevel}}

\begin{lemma}[a gauge that makes two sequences heads]\label{l:heads}
In Setting~\ref{set:cb}, let $j$ and $K$ be natural numbers and $p,q\colon\NN\to\CC^{\times}$. Then there is a real
$R_0$ with $\rho\le R_0$ such that $p$ and $q$ are heads of length $j$ for the gauge
$\gauge{\mathfrak{h}}{\rho\le R_0}$, \m{\IsHead{j}{p}} and \m{\IsHead{j}{q}}, and such that for every real
$R\ge\rho$ with $|p_0\cdots p_{l-1}|\le R$ and $|q_0\cdots q_{l-1}|\le R$ for all $l\le j$,
\[R_0^{K}\le R^{(K+6)(2j+1)}\cdot\hw{K}{j}{p}\,\hw{K}{j}{q}.\]
\end{lemma}
\begin{proof}
For this gauge, being a head of length $j$ means $|p_0\cdots p_{l-1}|^{-1}\le R_0$ for $l\le j$. Induction on $j$.
For $j=0$ take $R_0=\rho$; then $\rho^{K}\le R^{K+6}$ as $R\ge\rho\ge1$. Given $R_0$ for $j$, take
$R_0'=\max(R_0,|p_0\cdots p_j|^{-1},|q_0\cdots q_j|^{-1})$ for $j+1$. The inequality follows by two uses of the
elementary fact: if $L^{K}\le B$, $L,R\ge1$ and $0<|y|\le R$, then
$\max(L,|y|^{-1})^{K}\le B\,R^{K+6}\,\damp{K}{|y|^{-1}}$; indeed $R^{K+6}\damp{K}{|y|^{-1}}\ge1$, and if
$|y|^{-1}>L\ge1$ then $\damp{K}{|y|^{-1}}=|y|^{-K}$.
\end{proof}
\leanline{\lean{existsgaugeisHead}}

\begin{lemma}[summing over shells]\label{l:shells}
In Setting~\ref{set:cb}, let $I$ be non-empty, let $j$ be a natural number with $j\le\varrho$, let
$W\colon\Gone{\mathfrak{s}}\to\CC^I$ satisfy \m{\IsBessel{W}{1}}, and let $\varepsilon$ be a real number with
$\varepsilon\ge0$ and $\int|W|^2\dd\dot y\le\varepsilon^2|I|$. Let $R$ be real with $R\ge\rho$, let
$p,q\colon\NN\to\CC^{\times}$ be heads of length $j$ for the gauge $\gauge{\mathfrak{h}}{\rho\le R}$,
\m{\IsHead{j}{p}} and \m{\IsHead{j}{q}}, and suppose that $W(y)\ne0$ implies
\m{\IsNear{\cLC{\mathfrak{h}}}{j}{y}}. Let $t$ be real with $t\ge1$ and $f\colon\Gone{\mathfrak{s}}\to[0,\infty]$
with $f\le\formMajorant{j}{p}{q}{\Phi}{W}{\cdot}$, and suppose that, if $t>1$, then for all $x$
\[f(x)^2\le\Bigl(\frac{(\pi\rho^{24})^{k}}{t^{s_0}}\Bigr)^{2}\,\selfNormProduct{j}{p}{q}{\Phi}{W}{x}{\gauge{\mathfrak{h}}{\rho\le R}}.\]
Then
\[\bigl|\formConst{\mathfrak{c}}\bigr|\cdot\headDensity{j}{p}{q}{\cLC{\mathfrak{h}}}\int f\,\dd\groupMeasure{\cLC{\mathfrak{h}}}
\ \le\ \bandc{\mathfrak{s}}\;R^{10^{9}n^{4}}\,\rho^{24k}\,\frac{t^{9900000n^{3}}}{t^{s_0}}\,\varepsilon.\]
\end{lemma}
\begin{proof}
Write $\cL=\cLC{\mathfrak{h}}$, $\Gamma=\gauge{\mathfrak{h}}{\rho\le R}$, $\lambda=\haarRatio{\mathfrak{h}}$. For
$i\in\NN$ put $T_i=t\,R^{11ni}$, let $\Gamma^{(i)}$ be the gauge of radius $T_{i+1}^{99n}$ (\Inp{C3}) and
$P_i=\positionsOfSize{\cL}{\Gamma^{(i)}}$, the set of the $x$ with all entries of $x-1$ of absolute value at most
$T_{i+1}^{99n}$. The $P_i$ increase and cover $\Gone{\mathfrak{s}}$; let $D_i=P_i\smallsetminus\bigcup_{i'<i}P_{i'}$.
We apply Theorem~D on each $D_i$, with $\Gamma'=\Gamma^{(i)}$, $\Theta=\varepsilon(\lambda|I|)^2$ and
$e_i=(\pi\rho^{24})^{k}\rho^{9999n}T_i^{-s_0}$; its setting holds by Lemma~\ref{l:iscut} and \Inp{C2}.

\emph{The pairing norms.} By the Cauchy--Schwarz inequality in $\CC^I$ and by \Inp{C2}, as $|\Phi|^2$, $|W|^2$ are
$\UT{T}$-invariant, $\pairingNorm{\Phi}{W}\le\lambda\int|\Phi|^2\dd\dot y\cdot\lambda\int|W|^2\dd\dot y\le\lambda^2|I|^2\varepsilon^2$,
and likewise $\pairingNorm{\Phi}{\Phi}\le\lambda^2|I|^2$. So their product is at most $\Theta^2$.

\emph{The pointwise bound on $D_i$.} If $i=0$ and $t>1$ it is the hypothesis, as $T_0=t$. Otherwise we use
Theorem~A(a) with $\sigma=\vol\{c\in\CC:|c|<T_i^{-1}\}$ and $V=\Lambda(\Gamma)^{1287n}\vol(\Gamma_{11})^{k-1}$
(\Inp{T9}(a)). For $i=0$, $t=1$ we have $\sigma=\pi\ge1$ and nothing to check. For $i\ge1$ let $x\in D_i$, so
$x\notin P_{i-1}$, and let $y$ be as in the hypothesis of Theorem~A(a). If no $a$ with $j\le a<\varrho$ had
$\vol(\floorValues{\cL}{a}{y}{x})\le\sigma$, each of these open sets would contain a $c$ with $|c|\ge T_i^{-1}$, a
unit whose inverse has size~$1$ for the gauge of radius $T_i$; then \Inp{T14}(e), with that gauge in the place of
$\Gamma$ and $\Gamma_S=\Gamma$, would give $|(x-1)_{\alpha\beta}|\le T_i^{99n}$ for all entries, that is $x\in P_{i-1}$.
(Its other hypotheses: $R^{10n}\le T_i$; and $\Phi(y)\ne0\ne\Phi(yx)$ puts $uy$, $v(yx)$ in the box for some
$u,v\in\UT{T}$, with entries of absolute value at most $\rho\le T_i$.) Hence
$f(x)^2\le\sigma^{s_0}V^2\,\selfNormProduct{j}{p}{q}{\Phi}{W}{x}{\Gamma}$, and $\sigma^{s_0}V^2\le e_i^2$ is a
comparison of explicit powers.

\emph{The sum.} Theorem~D bounds $\headDensity{j}{p}{q}{\cL}\int_{D_i}f\,\dd\groupMeasure{\cL}$ by
$e_i\,\pi^{2(k-1)\varrho}\,R^{3000000n^3+2}\,T_{i+1}^{9900000n^3}\,\Theta$ (the floor modules at $\gL{\cL}=1$ are $1$,
$c_{\CC}=\pi^2$, $|\Omega|=k-1$). Since $s_0>10^{9}n^{3}$, the factor $T_i^{-s_0}$ in $e_i$ beats
$T_{i+1}^{9900000n^3}$, and the sum over $i$ is at most the stated bound times $\sum_i2^{-(i+1)}=1$ divided by
$|\formConst{\mathfrak{c}}|(\lambda|I|)^2=\volU{k}/\volU{r}$. This last comparison of exponents is carried out by the
formal text and was not redone by hand.
\end{proof}
\leanline{\lean{formConstmullintegralle}}

\begin{lemma}[the majorant is integrable]\label{l:fin}
In Setting~\ref{set:cb}, let $I$ be non-empty, $j\le\varrho$ and $p,q\colon\NN\to\CC^{\times}$. Then
$\int\formMajorant{j}{p}{q}{\Phi}{\Phi}{x}\,\dd\groupMeasure{\cLC{\mathfrak{h}}}(x)<\infty$.
\end{lemma}
\begin{proof}
Take $R_0$ from Lemma~\ref{l:heads} (with $K=0$) and apply Lemma~\ref{l:shells} with $W=\Phi$, $\varepsilon=1$,
$t=1$ and $f$ the majorant; $\formConst{\mathfrak{c}}$ and the head density are not $0$.
\end{proof}
\leanline{\lean{lintegralformMajorantnetop}}

\begin{lemma}[the explicit bound for a form]\label{l:norm}
In Setting~\ref{set:cb}, let $I$ be non-empty, $j\le\varrho$, $p,q\colon\NN\to\CC^{\times}$, $y\in\CC^{\times}$,
and $\omega\colon\Gone{\mathfrak{s}}\to\CC$ with $|\omega(x)|\le1$ for all $x$. Suppose that for every gauge
$\Gamma$ in $\Gauges{\cLC{\mathfrak{h}}}$ for which $p$ is a head of length $j$, and every $x$,
\[\formModulus{j}{p}{q}{\Phi}{x}{\omega}^{2}\le\Bigl(\frac{\pi^{2}}{|y|^{2}}\Bigr)^{k-1}\selfNormProduct{j}{p}{q}{\Phi}{\Phi}{x}{\Gamma}.\]
Then
\[\bigl|\formConst{\mathfrak{c}}\cdot\floorForm{j}{p}{q}{\Phi}{\omega}\bigr|\le\bandc{\mathfrak{s}}\,b^{2000n^{9}}\,
\damp{k-n^{7}}{|y|^{-1}}\cdot\hw{n^6}{j}{p}\,\hw{n^6}{j}{q}.\]
\end{lemma}
\begin{proof}
By Lemma~\ref{l:fmm}(i) the left side is at most $|\formConst{\mathfrak{c}}|$ times the head density times
$\int\formModulus{j}{p}{q}{\Phi}{x}{\omega}\dd\groupMeasure{\cL}(x)$. Take $R_0$ from Lemma~\ref{l:heads} with
$K=n^6$. \emph{Case 1:} some partial product of $p$ or $q$ of length at most $j$ exceeds $\rho^{9n}$ in absolute
value. Then the form is $0$: by Theorem~A(a) with $W'=\Phi$, $\sigma=0$ and $V=\infty$ the majorant vanishes, since
no $y'$ satisfies the conditions of that hypothesis (they would put all partial products into the set of size
$9n$ of the gauge of radius $\rho$, by \Inp{T14}(c)); and $s_0\ne0$. \emph{Case 2:} otherwise
$R_0^{n^6}\le\rho^{9n(n^6+6)(2j+1)}\hw{n^6}{j}{p}\hw{n^6}{j}{q}$. Apply Lemma~\ref{l:shells} with $W=\Phi$,
$\varepsilon=1$, $R=R_0$, $t=\max(1,\rho^{24}|y|)$ and $f$ the modulus, which is at most the majorant
(Lemma~\ref{l:fmm}(ii)); if $t>1$ the hypothesis gives $f^2\le(\pi\rho^{24}/t)^{2(k-1)}$ times the self-norm product,
and $(\pi\rho^{24}/t)^{k-1}\le(\pi\rho^{24})^{k}/t^{s_0}$ because $s_0\le k-1$. The resulting bound
$\bandc{\mathfrak{s}}R_0^{10^9n^4}\rho^{24k}t^{9900000n^3-s_0}$ is at most the right side: $10^9n^4\le n^6$,
$\rho=b^{6n+2}$, $1/t\le\min(1,\rho^{-24}|y|^{-1})$ and $s_0-9900000n^3\ge k-n^7+6$. This comparison of exponents is
that of the formal text.
\end{proof}
\leanline{\lean{normfloorFormle}}

\begin{lemma}[size of the support]\label{l:supp}
In Setting~\ref{set:cut}, let $p\colon\NN\to A^{\times}$ be a head of length $\varrho$ for $\Gamma$,
\m{\IsHead{\varrho}{p}}, let $x\in\Hl{\varrho}{T}$, and let $y\in\GT{T}$ with $W(y)\ne0$ and
$W((\wl{\varrho}{p}x)^{-1}y)\ne0$. Then \m{\HasMatSize{60n^2}{y}} for $\Gamma$.
\end{lemma}
\begin{proof}
There is $u\in\UT{T}=\Nb{\varrho}{T}$ with $uy$ in the box; and $x^{-1}\wl{\varrho}{p}^{-1}y$ is near at $\varrho$,
hence so is $\wl{\varrho}{p}^{-1}y$ (\Inp{T8}). Now use \Inp{T14}(a).
\end{proof}
\leanline{\lean{hasMatSizeofnezero}\sep\lean{isCompacthasMatSize}}

\begin{lemma}[the inner product of two correlations]\label{l:corr}
In Setting~\ref{set:cut}, let $\mu$ be a Haar measure~$\mu$ of $\GT{T}$, suppose that $W$ is continuous, let
$p,q\colon\NN\to A^{\times}$ be heads of length $\varrho$ for $\Gamma$, \m{\IsHead{\varrho}{p}} and
\m{\IsHead{\varrho}{q}}, and let $x\in\Hl{\varrho}{T}$. Then, with $\gamma_p=\wl{\varrho}{p}x$ and
$\gamma_q=\wl{\varrho}{q}x$,
\[\ip{\corr{\mu}{W}{\gamma_p}}{\corr{\mu}{W}{\gamma_q}}=\iint\ip{W(y)}{W(y')}\;
\overline{\ip{W(\gamma_p^{-1}y)}{W(\gamma_q^{-1}y')}}\;\dd\mu(y)\,\dd\mu(y').\]
\end{lemma}
\begin{proof}
In coordinates the left side is $\sum_{i,i'}\overline{\int W_i(y)\overline{W_{i'}(\gamma_p^{-1}y)}\dd\mu(y)}\cdot
\int W_i(y')\overline{W_{i'}(\gamma_q^{-1}y')}\dd\mu(y')$. Each product of two integrals is a double integral,
because the integrands are continuous and vanish outside the set of the $y$ of size at most $60n^2$
(Lemma~\ref{l:supp}), which is compact: it is closed, and the $\Gamma_i$ are compact. Summing over $i,i'$ gives the
integrand on the right.
\end{proof}
\leanline{\lean{innerbesselCorrelation}}

\begin{lemma}[continuity of the top vector]\label{l:cont}
In Setting~\ref{set:cb}, the map $p\mapsto\sieveVector{\Phi}{p}$ is continuous on the space of the sequences
$p\colon\NN\to\CC^{\times}$ with the product topology.
\end{lemma}
\begin{proof}
The scalar factor is continuous. Fix $p$ and a gauge radius $R_0$ for which $p$ is a head of length $\varrho$
(Lemma~\ref{l:heads}); the sequences that are heads for the radius $R_0+1$ form a neighbourhood of $p$. For them the
integrand of each coordinate of $\corr{\dd y}{\Phi}{\wl{\varrho}{p}\boost{\varrho}{p_\varrho}}$ vanishes outside one
compact set (Lemma~\ref{l:supp}, with $x=\boost{\varrho}{p_\varrho}\in\Hl{\varrho}{T}$), and it is continuous in
$(p,y)$ because the boost depends continuously on $p_\varrho$. An integral of a jointly continuous function over a
fixed compact set is continuous in the parameter.
\end{proof}
\leanline{\lean{continuoussieveVector}}

\begin{proof}[Proof of Theorem E]
Write $\cL=\cLC{\mathfrak{h}}$, $\lambda=\haarRatio{\mathfrak{h}}$, and let $J$ be the family of the statement.
$\Phi$ is cut for $\cL$ (Lemma~\ref{l:iscut}), $\groupMeasure{\cL}$ is a Haar measure (\Inp{C2}), every point is a
position, and $|\Omega|=k-1$, $\gL{\cL}=1$, $c_{\CC}/\modA(y)=\pi^2/|y|^2$.

\emph{\textup{(S1)} for $i<\varrho$.} Theorem~B with $t=\formConst{\mathfrak{c}}$ and $\eta=\xi$, whose hypothesis is
Lemma~\ref{l:fin}, shows that $f$ is integrable with integral
$\formConst{\mathfrak{c}}\floorForm{i}{p}{q}{\Phi}{\omega_\xi}$. For $\xi=0$ this is $J_i(p,q)$. For $\xi\ne0$,
Lemma~\ref{l:norm} with $y=\xi$ gives the bound; its hypothesis is Theorem~A(b) with $\eta=\xi$, in every gauge for
which $p$ is a head.

\emph{\textup{(S1)} for $i=\varrho$.} We apply Lemma~\ref{l:succ} with $b=\varrho$,
$\zeta=\floorMeasure{\cL}{n}$, carried by $\Hl{n}{T}=\Ql{n}{T}\cap\Hl{\varrho}{T}$, $v(c)=\boost{\varrho}{c}$ and
$\rho(c)=\frac{2\pi^k}{(k-1)!}\max(|c|^2-1,0)^{k-1}$ (\Inp{C1}), and check the formula for
$\Psi(c)=\psi_{\CC}(\xi c)J_n(p',q')$, $p'=p[\varrho\mapsto c]$, $q'=q[\varrho\mapsto c]$. For $|c|\le1$ both sides
are $0$. Let $|c|>1$; then $\ZT{T}{\boost{\varrho}{c}}_{\varrho\varrho}=c$. On the two top floors
$\eta_\varrho=\eta_n=\delta_1$, so $(\floorMeasure{\cL}{n}*\floorMeasure{\cL}{\varrho})*\stackMeasure{\cL}_{0,\varrho}
=\groupMeasure{\cL}=\lambda\,\dd y$. The averaged kernels at $n$ are averages over $\Hl{n}{T}$ of
$\ip{\Phi(hg)}{\Phi(hg')}=\ip{\Phi(g)}{\Phi(g')}$, as $\mathrm{rep}(h)$ is unitary; so, with
$\gamma_{p'}=\wl{\varrho}{p}\boost{\varrho}{c}$ and $d\gamma_{p'}^{-1}=\gamma_{q'}^{-1}$,
\[\formIntegrandOne{n}{\Phi}{x}{d}{\gamma_{p'}}{g}=\ip{\Phi(g)}{\Phi(gx)}\,\overline{\ip{\Phi(\gamma_{p'}^{-1}g)}{\Phi(\gamma_{q'}^{-1}gx)}}.\]
Integrating over $g$ and $x$ and substituting $y'=gx$ gives $\lambda^2\ip{\corr{\dd y}{\Phi}{\gamma_{p'}}}{\corr{\dd y}{\Phi}{\gamma_{q'}}}$
by Lemma~\ref{l:corr} (with a gauge from Lemma~\ref{l:heads}). The constants agree:
$\formConst{\mathfrak{c}}\lambda^2\cdot\headDensity{\varrho}{p}{q}{\cL}\cdot\rho(c)$ equals the product $N(p')N(q')$ of
the scalar factors of Definition~\ref{d:sieve}, because
$\cellc{\mathfrak{s}}=\frac{\volU{k}}{\volU{r}}\frac{2\pi^k}{(k-1)!}$. The bound for $\xi\ne0$ follows as before.

\emph{\textup{(S2)}.} The matrix is a Gram matrix; continuity is Lemma~\ref{l:cont}.
\emph{\textup{(S4)}.} Lemma~\ref{l:norm} at $j=i+1$ with $\omega=1$, $y=\xi$, and Theorem~A(c).
\emph{\textup{(S5)}.} If $|p_\varrho|\le1$ both sides are $0$. Otherwise the Cartan decomposition gives
$\sigma=h\,\boost{\varrho}{p_\varrho}h'$; write both sides by Lemma~\ref{l:corr}, substitute $(y,y')\mapsto(hy,hy')$ and
use $\Phi(hy)=\mathrm{rep}(h)\Phi(y)$ with $\mathrm{rep}(h)$ unitary to remove $h$ and $h'$.

\emph{\textup{(S3)}, the energy.} For $j=0$ the head density is $1$, $\stackMeasure{\cL}_{0,0}=\delta_1$ and
$\weylRatio{0}{p}{q}=1$, so $J_0(p,q)=\formConst{\mathfrak{c}}\int|\delta_\Phi|^2\dd\groupMeasure{\cL}$ with
$\delta_\Phi(x)=\upperPairing{\cL}{0}{\translatePairing{\Phi}{\Phi}{1}{\cdot}{\cdot}}{1}{x}$. Let
$\delta(x)=\upperPairing{\cL}{0}{\translatePairing{\Phi}{\Phi_0}{1}{\cdot}{\cdot}}{1}{x}$. Then
$\delta_\Phi-\delta$ is the averaged kernel of $\Phi$ and $\Phi-\Phi_0$, and from $|\delta|\ge\bigl||\delta_\Phi|-|\delta_\Phi-\delta|\bigr|$
and \Inp{T7}(a),
\[|\delta_\Phi(x)|^2\le|\delta(x)|^2+2\,\formMajorant{0}{p}{q}{\Phi}{\Phi-\Phi_0}{x}.\]
The second term: $\Phi-\Phi_0$ is a Bessel function of frequency $1$ with
$\int|\Phi-\Phi_0|^2\dd\dot y\le b^{-10r^2}|I|$, and Lemma~\ref{l:shells} with $j=0$, $R=\rho$, $t=1$,
$\varepsilon=b^{-5r^2}$ bounds $|\formConst{\mathfrak{c}}|$ times its integral by
$\bandc{\mathfrak{s}}\rho^{10^9n^4+24k}b^{-5r^2}$; with $\deg\bandc{\mathfrak{s}}\le b^{4r^2}$ and the sizes this is
at most $b^{-10}/(2\deg)$, a comparison of exponents done by the formal text. The first term: we claim
$\int|\delta|^2\dd y\le\lambda^2|I|^2/\deg$; then, by \Inp{C2},
$|\formConst{\mathfrak{c}}|\int|\delta|^2\dd\groupMeasure{\cL}\le\frac{\volU{k}}{\volU{r}}\lambda\cdot\frac1{\deg}\le\frac1{\deg}$.
It suffices to bound $E=\int_{s}|\delta|^2\dd y$ for a set $s$ of finite measure. The function $\delta$ is bounded
(a pairing kernel, Lemma~\ref{l:pk}); let $\beta=\mathbf{1}_s\overline{\delta}$, bounded and integrable, with
$\int|\beta|^2=E<\infty$. By Fubini's theorem and \Inp{C2},
\[E=\int\beta\,\delta\,\dd y=\int\Bigl\langle\Phi(\sigma),\int\beta(x)\Phi_0(\sigma x)\,\dd x\Bigr\rangle\dd\upperMeasure{\cL}{0}(\sigma)
\le\lambda\int|\Phi(\gamma)|\,\Bigl|\int\beta(x)\Phi_0(\gamma x)\,\dd x\Bigr|\,\dd\dot\gamma,\]
the integrand being $\UT{T}$-invariant. By Cauchy--Schwarz, $\int|\Phi|^2\dd\dot y\le|I|$ and Schur's bound for
$\Phi_0$, $E^2\le\lambda^2\,|I|\cdot\frac{|I|}{\deg}\,E$, hence $E\le\lambda^2|I|^2/\deg$.
\end{proof}

\sect{Where the results are used}{s:used}
Theorem~E is the last hypothesis in the construction of a \emph{Bessel system} at a real place, the local datum of
the second moment of Poincar\'e series in the unit on the second moment (unit 148, Chapter~VI of the write-up): there
the band limits (S1)--(S5) at all infinite places and their analogues at the chosen finite places are multiplied
into one sieve system. Theorems A to~D are stated for a general place; the artifact applies them at $\CC$ in this
note, and their statements cover the products of finite places for which the authors' comment announces the same
proof.

\sect{Notes on the formal text}{s:notes}
\textbf{Implicit arguments.} The level is an implicit argument of the forms, of the moduli, of the majorant, of the
pairing norm, of the predicate $\operatorname{Cut}$ and of the self-norm product; the gauge is an implicit argument of
the predicates for sizes and heads. The notation shows the explicit arguments of the Lean definitions, and each
statement names the level and the gauge in words.

\textbf{Argument lists not confirmed by Lean.} For the Bessel predicate, the cut predicate, the pairing norm, the
top vector, and for the Lean theorems behind Theorems B, C, Lemma~\ref{l:fmm}(i), Lemma~\ref{l:succ} and
Lemma~\ref{l:heads}, the reader's dossier derived the list of arguments by Lean's rules for section variables,
because the census misreads a name bound inside the declaration ($x$, $L$) or a line \texttt{variable (F) in}. Nothing
was compiled, so these lists are not confirmed by Lean. The statements above follow the Lean text of each
declaration, read in full; the statement of Theorem~B continues after a \texttt{let} and was read to its end.

\textbf{Differences from the natural statements.} (1) The integrals of complex functions are Bochner integrals that
are $0$ when the function is not integrable; so the forms are defined without an integrability hypothesis, and
Theorem~B carries one. (2) The head density enters the forms through its real value, which is $0$ when it is
infinite. (3) Exponents $\varrho-(a+1)$ are truncated at $0$; in every use $a<\varrho$. (4) Theorem~A(a) allows
$V=\infty$ and $\sigma=0$; this is used in Lemma~\ref{l:norm}. (5) In Theorem~A(b) the weight is written with the
entry $(j,j)$ of the frame Gram matrix, which needs $j\le\varrho$. (6) Theorem~D and Lemma~\ref{l:shells} ask no
measurability of $f$ and $S$: the integrals are upper integrals in the sense of Lebesgue. (7) In
Lemma~\ref{l:shells} the hypothesis on $f^2$ is asked only if $t>1$. (8) Theorem~E is an existence statement in
the Lean text; the family $J$ is the witness of its proof. (9) Lemma~\ref{l:corr} takes any Haar measure~$\mu$ and
does not ask that $\groupMeasure{\cL}$ be one. (10) The functional $\mathrm{rep}$ and the number $\mathsf{n}$ of cut
Bessel data enter only through the cut pair.

\textbf{A declaration outside the census.} The definition of the level $\cLC{\mathfrak{h}}$ at the complex place
is in the Lean text (the file of the gauge at $\CC$, lines 31 to 41), but the census of the artifact does not list
it, because its attribute runs over two lines; so it is in no unit and is not a member of this one. Input~\Inp{C1}
states it from the Lean text, and the second reader recovered the same declaration from the pinned source. Its Lean
name is in Appendix~A.

\textbf{Small results.} Five declarations have no statement of their own: the two top floors of the stack measure
(used in the proof of Theorem~E), compactness of the set of the elements of bounded size (Lemma~\ref{l:corr}), nearness
of the points where a cut function does not vanish (after Definition~\ref{d:cut}), the kernel with $p=q=1$
(Lemma~\ref{l:pk}), and the volume exponent at $\CC$ (after Definition~\ref{d:cutbessel}).
\leanline{\lean{stackMeasuretop}}

\subhead{How this note was written.} From the structured proofs of the reader's dossier for the three files, and from
the Lean text. I read the three Lean files completely, statements and proofs, together with the definitions of the
files of Note~T on which they rest; the proofs behind Theorem~C and Theorem~E, which the reader marked as the least
certain, were read line by line, and the induction of Theorem~C and the energy bound (S3) are written from that
reading. The proofs of Theorem~A, Lemma~\ref{l:snp} and Lemma~\ref{l:succ} were also rebuilt from the Lean terms
(which lemma is applied to which arguments). Not redone by hand: the three final comparisons of explicit exponents
in Lemmas~\ref{l:shells}, \ref{l:norm} and in (S3), which the formal text settles by its arithmetic tactics; they
are named in the proofs. The inputs \Inp{T1}--\Inp{T14} and \Inp{C1}--\Inp{C4} were compared with their Lean
statements; in \Inp{T9}(a) and \Inp{T12} two Lean lemmas are combined and the exponent $1287n=99n\cdot13$ is computed
from them. The program check of the unit's dossier covered \UnitSteps{} leaf steps of \UnitProofs{} proofs.

\head{Appendix A. Concordance with the Lean artifact}
The letters of this note and the names of the formal text: $\cL$ is \texttt{F} (a level) or the complex level of
\texttt{hF}; $\Omega$ is \texttt{ι}; $\mathfrak{m}$ is \texttt{m}; $\Phi$, $\Phi_0$ are \texttt{F}, \texttt{F₀} (or
\texttt{W}) at the complex place; $s_0$ is \texttt{s}; $\mathsf{n}$ is \texttt{n}; $R$ is \texttt{L}; $\modA$,
$\psi_A$, $c_A$ are the fields \texttt{md}, \texttt{ψ}, \texttt{c} of the place; $\Hl{l}{T}$, $\Ql{l}{T}$,
$\Nb{l}{T}$, $\Nf{l}{T}$, $\UT{T}$ are \texttt{GT}, \texttt{unitriGroup}, \texttt{unipotentBelow},
\texttt{floorUnipotent}, \texttt{besselU}; $\sd$ is \texttt{simpleRootSum}; $\ZT{T}{\cdot}$ is \texttt{frameGram}; $\cLC{\mathfrak{h}}$ is
\texttt{RealPlaceSizes.complexLevel}, which the census does not list.
Two statements as the artifact has them:
\quoted{formModulussqleofheads}
\quoted{hasBandLimitsofisCutBessel}
{\fontsize{8.6}{11.6}\selectfont\renewcommand{\leanfont}{\ttfamily\fontsize{7.7}{10}\selectfont\color{leancol}}
\setlength{\LTleft}{0pt}\setlength{\LTright}{\fill}\setlength{\tabcolsep}{0pt}
\begin{longtable}{@{}p{30mm}p{56mm}p{43mm}r@{}}
\textbf{In this note} & \textbf{Lean declaration} & \textbf{File} & \textbf{Lines}\\\midrule\endhead
\ConcordanceRows
\end{longtable}}

\head{Notation}
\begingroup\fontsize{8.6}{10.9}\selectfont\raggedright
\setlength{\columnsep}{9mm}\setlength{\parskip}{0pt}
\begin{multicols}{2}
\nt{$A$, $\bar x$, $\dd x$}{the place: a ring with involution and Haar measure}
\nt{$\modA$, $\psi_A$, $c_A$}{its module, character and Plancherel constant}
\nt{$\tau_A$, $\delta_A$}{its trace-one element and skew unit}
\nt{$\mathfrak{s}=(\varrho,\kappa,k_0)$}{the shape}
\nt{$n$, $k$, $r$}{$\varrho+1$; $|\kappa|$; $2\varrho+k+1$, the rank of $V$}
\nt{$V$, $V_l$, $e_a$, $e^*_a$, $x_i$, $z$}{the space, its floors, its basis}
\nt{$f_a$, $\ZT{T}{g}$}{the frame; the frame Gram matrix of $g$}
\nt{$T$, $\JT{T}$, $h_T$, $\GT{T}$}{the form, its split Gram matrix, the hermitian form, the unitary group}
\nt{$\Hl{l}{T}$, $\Ql{l}{T}$}{stabiliser of the floors below $l$; unitriangular on the first $l$ frame vectors}
\nt{$\Nb{l}{T}$, $\Nf{l}{T}$, $\UT{T}$}{unipotent groups of Section~\ref{s:space}}
\nt{$\sd(u)$}{sum of the entries of $u$ at $(f_a,f_{a+1})$}
\nt{$\wl{l}{p}$, $\wo{a}{c}$}{Weyl element; floor Weyl element}
\nt{$\Omega$, $\mathfrak{m}$}{index set of the orthonormal part and number of extra pairs of a hyperbolic basis}
\nt{$\cL$, $E(\cL)$, $\gL{\cL}$}{a level, its basis matrix, its unit}
\nt{$\Ball{\cL}$, $\BoxL{\cL}$, $\Posn{\cL}$}{its ball, box, positions}
\nt{$\floorMeasure{\cL}{l}$, $\eta_l$, $\etab{j}$, $\varpi_l$}{floor measure, Heisenberg measures, Weyl law}
\nt{$\stackMeasure{\cL}_{a,i}$, $\groupMeasure{\cL}$, $\upperMeasure{\cL}{l}$}{convolutions of floor measures}
\nt{$\quadric{\cL}{b}$, $\delta_s$, $*$}{quadric measure on $V_b$; unit mass; convolution}
\nt{$\Gamma$, $\Gamma_i$, $\Lambda(\Gamma)$}{a gauge, its sets, its constant}
\nt{$\operatorname{sz}$, $\mathrm{Vec}_i$, $\operatorname{Hd}_j$}{size of a matrix, vectors of size $i$, heads}
\nt{$\operatorname{Near}_j$, $\operatorname{Val}_a$, $\mathsf{s}$, $\mathsf{w}_j$, $\mathsf{d}_l$}{nearness, floor values, volume exponent, weight, floor module}
\nt{$\operatorname{Mass}_j$, $\operatorname{Fib}$}{mass of a floor; fibre bound of Note~T}
\nt{$\mathsf{Av}_j$, $\mathsf{Pv}_j$, $\mathsf{Sq}_j$, $\mathsf{Sm}_j$, $\operatorname{SN}_j$}{averaged kernel, pivot, weighted square, its sum, self-norm}
\nt{$\operatorname{PK}_j$, $\operatorname{Pos}$, $C_l$, $C_{\mathrm{gl}}$}{pairing kernels, positions of size $1$, constants of a gauge}
\nt{$I$, $\ip{v}{w}$, $|v|$}{finite set; inner product and norm of $\CC^I$}
\nt{$W$, $W'$, $\operatorname{Bes}$, $\operatorname{Cut}$}{functions; Bessel and cut functions}
\nt{$\operatorname{PN}$, $\mathsf{h}_j$, $\mathsf{I}^1_j$, $\mathsf{I}_j$}{pairing norm, head density, integrands}
\nt{$\mathcal{I}_j$, $\mathcal{I}^{\mathrm{mod}}_j$, $\mathcal{I}^{\mathrm{maj}}_j$, $\operatorname{SNP}_j$}{floor form, modulus, majorant, self-norm product}
\nt{$p$, $q$, $d$, $w$}{sequences of units; Weyl ratio; a Weyl element}
\nt{$\omega$, $\omega_\eta$, $\eta$, $\xi$}{weights; frequencies}
\nt{$j$, $l$, $a$, $b$, $i$, $i_1$, $i_2$}{floors and sizes (at $\CC$, $b$ is also the third scale)}
\nt{$g$, $x$, $y$, $s$, $u$, $v$, $h$, $m$, $t$}{group elements, vectors, units; $t$ also a real or complex number}
\nt{$\sigma$, $\tau$, $\nu$, $\gamma$}{variables of integration; also $\sigma$ a volume, $\gamma$ a group element}
\nt{$k$, $K$, $V$, $C$, $M$, $D$, $E$, $N$}{a kernel (also $k=|\kappa|$); constants}
\nt{$\Phi$, $\Phi_0$, $\mathrm{rep}$}{cut and full function at $\CC$; representation of $\mathrm{U}(\kappa)$}
\nt{$\mathfrak{h}$, $\mathfrak{c}$, $\rho$, $s_0$, $\mathsf{n}$}{sizes, cut Bessel data, their numbers}
\nt{$\cLC{\mathfrak{h}}$, $\lambda$, $\Gamma(\mathfrak{h};\rho\le R)$, $R$, $R_0$}{level, Haar ratio, gauge at $\CC$, radii}
\nt{$\operatorname{RPS}$, $\operatorname{CB}$, $c_{\mathrm{f}}$, $\Xi$}{sizes, cut Bessel data, constant of the forms, top vector}
\nt{$\Gone{\mathfrak{s}}$, $\dd y$, $\dd\dot y$}{the group at $\CC$, its Haar measure, the measure of $\UT{T}\backslash\Gone{\mathfrak{s}}$}
\nt{$\corr{\mu}{W}{\gamma}$, $\boost{\varrho}{c}$}{correlation tensor; boost}
\nt{$\cellc{\mathfrak{s}}$, $\bandc{\mathfrak{s}}$, $\volU{N}$}{constants; volume of a unitary group}
\nt{$\damp{K}{t}$, $\hw{K}{i}{p}$, $\beta$, $\deg$}{damped power, head weight, band function, formal degree}
\nt{$J$, $J_i$, $f$, $F$, $F_1$, $F_2$, $G$, $B$, $P$, $S$, $T_i$, $e$, $c$}{forms and auxiliary functions, sets and numbers of the proofs}
\nt{$\mathsf{r}$, $\mathsf{u}$, $\mathsf{v}$, $\mathsf{A}$, $\mathsf{W}_l$, $\Phi_l$, $\varepsilon$, $\Theta$, $\zeta$, $\Psi$, $\varphi$, $\mu$, $\alpha$}{auxiliary objects of the proofs}
\nt{$\pi$, $\infty$, $\mathbf{1}$, $L$, $H$, $O$, $Q$, $U$, $X$, $Y$, $Z$, $o$, $z$}{the number; infinity; indicator; letters inside names}
\end{multicols}
\endgroup

\end{document}
